%%%
%%% 3画
%%%
\section*{3画}\addcontentsline{toc}{section}{3画}\addcontentsline{loh}{figure}{\#\#\#\# 3画}

%%%%%%%%%% 万 %%%%%%%%%%
\subsection*{万}\addcontentsline{loh}{figure}{万}

\begin{Entry}{万}{3}{⼀}
  \begin{Phonetics}{万}{wan4}[][HSK 2]
    \definition*{s.}{Sobrenome: Wan}
    \definition{adv.}{absolutamente; indica um grau extremamente alto, equivalente a 完全, 绝对 e 极}
    \definition{num.}{dez mil; 10.000; 1.0000 | miríade; um número muito grande}
  \seealsoref{极}{ji2}
  \seealsoref{绝对}{jue2dui4}
  \seealsoref{完全}{wan2quan2}
  \end{Phonetics}
\end{Entry}

\begin{Entry}{万一}{3,1}{⼀,⼀}
  \begin{Phonetics}{万一}{wan4yi1}[][HSK 4]
    \definition{conj.}{por via das dúvidas; se por acaso; só por precaução; expressa uma suposição muito improvável (usado para coisas desagradáveis)}
    \definition{num.}{um décimo milionésimo; uma porcentagem muito pequena}
    \definition{s.}{contingência; eventualidade; contingências muito improváveis}
  \synonymref{如果}{ru2guo3}
  \synonymref{要是}{yao4shi5}
  \synonymref{一旦}{yi2dan4}
  \end{Phonetics}
\end{Entry}

\begin{Entry}{万万}{3,3}{⼀,⼀}
  \begin{Phonetics}{万万}{wan4wan4}[][HSK 7-9]
    \definition{adv.}{totalmente; absolutamente; em qualquer caso; não importa o que aconteça}
    \definition{num.}{cem milhões; 100.000.000; 1.0000.0000}
  \synonymref{绝对}{jue2dui4}
  \synonymref{千万}{qian1wan4}
  \synonymref{完全}{wan2quan2}
  \synonymref{一概}{yi2gai4}
  \end{Phonetics}
\end{Entry}

\begin{Entry}{万分}{3,4}{⼀,⼑}
  \begin{Phonetics}{万分}{wan4fen1}[][HSK 7-9]
    \definition{adv.}{extremamente; muito}
  \synonymref{非常}{fei1chang2}
  \synonymref{极度}{ji2du4}
  \synonymref{极端}{ji2duan1}
  \synonymref{十分}{shi2fen1}
  \synonymref{特别}{te4bie2}
  \synonymref{异常}{yi4chang2}
  \end{Phonetics}
\end{Entry}

\begin{Entry}{万无一失}{3,4,1,5}{⼀,⽆,⼀,⼤}
  \begin{Phonetics}{万无一失}{wan4wu2-yi1shi1}[][HSK 7-9]
    \definition{expr.}{não há perigo de algo dar errado; esteja do lado seguro; \dots não pode falhar em nenhuma circunstância; garantir sucesso total; nenhum risco; não há chance de erro; (faça com que seja mais do que provável que) nada dê errado; perfeitamente seguro; infalível; as chances são de mil para uma, não falharemos}
  \end{Phonetics}
\end{Entry}

\begin{Entry}{万古长青}{3,5,4,8}{⼀,⼝,⾧,⾭}
  \begin{Phonetics}{万古长青}{wan4gu3-chang2qing1}[][HSK 7-9]
    \definition{expr.}{seja perene; seja sempre verde; perene e eterno; sempre vivo; florescer para sempre; durar para sempre; permanecer fresco para sempre; sempre será verde como pinheiros e ciprestes por milhares de gerações; uma metáfora para um espírito nobre ou uma amizade profunda que nunca desaparecerá}
  \synonymref{欣欣向荣}{xin1xin1-xiang4rong2}
  \end{Phonetics}
\end{Entry}

\begin{Entry}{万圣节}{3,5,5}{⼀,⼟,⾋}
  \begin{Phonetics}{万圣节}{wan4sheng4jie2}
    \definition*{s.}{Dia de Todos os Santos}
  \seealsoref{万圣节前夕}{wan4sheng4jie2 qian2xi1}
  \end{Phonetics}
\end{Entry}

\begin{Entry}{万圣节前夕}{3,5,5,9,3}{⼀,⼟,⾋,⼑,⼣}
  \begin{Phonetics}{万圣节前夕}{wan4sheng4jie2 qian2xi1}
    \definition*{s.}{Véspera do Dia de Todos os Santos | Halloween}
  \seealsoref{万圣节}{wan4sheng4jie2}
  \end{Phonetics}
\end{Entry}

\begin{Entry}{万物}{3,8}{⼀,⽜}
  \begin{Phonetics}{万物}{wan4wu4}
    \definition{s.}{toda a criação; todas as coisas na Terra; todos os seres vivos; tudo no universo}
  \synonymref{生物}{sheng1wu4}
  \end{Phonetics}
\end{Entry}

\begin{Entry}{万能}{3,10}{⼀,⾁}
  \begin{Phonetics}{万能}{wan4neng2}[][HSK 7-9]
    \definition{adj.}{onipotente; todo"-poderoso | universal; multifacetado; de múltiplos usos}
  \synonymref{全能}{quan2neng2}
  \end{Phonetics}
\end{Entry}

\begin{Entry}{万福}{3,13}{⼀,⽰}
  \begin{Phonetics}{万福}{wan4fu2}
    \definition{s.}{(antigo) reverência feminina; reverência}
  \end{Phonetics}
\end{Entry}

%%%%%%%%%% 丈 %%%%%%%%%%
\subsection*{丈}\addcontentsline{loh}{figure}{丈}

\begin{Entry}{丈}{3}{⼀}
  \begin{Phonetics}{丈}{zhang4}
    \definition{clas.}{zhang, uma unidade tradicional de comprimento, igual a 10 市尺 e equivalente a 3,333 metros ou 3,65 jardas}
    \definition{s.}{zhang, uma unidade de comprimento (= 3,333\dots metros)}
    \definition{s.}{sênior; ancião | marido (em certos termos de parentesco) | tratamento respeitoso ao idoso na China antiga; um título respeitoso para homens idosos nos tempos antigos | uma forma de tratamento para certos parentes do sexo masculino por casamento}
  \seealsoref{市尺}{shi4chi3}
  \end{Phonetics}
\end{Entry}

\begin{Entry}{丈夫}{3,4}{⼀,⼤}
  \begin{Phonetics}{丈夫}{zhang4fu5}[][HSK 4]
    \definition[个,位,名]{s.}{marido; esposo}
  \end{Phonetics}
\end{Entry}

%%%%%%%%%% 三 %%%%%%%%%%
\subsection*{三}\addcontentsline{loh}{figure}{三}

\begin{Entry}{三}{3}{⼀}
  \begin{Phonetics}{三}{san1}[][HSK 1]
    \definition*{s.}{Sobrenome: San}
    \definition{num.}{três; 3 | muitos; vários; mais de dois; referindo"-se a muitos ou à maioria | alguns; poucos; menos; não muitos}
  \end{Phonetics}
\end{Entry}

\begin{Entry}{三角}{3,7}{⼀,⾓}
  \begin{Phonetics}{三角}{san1jiao3}[][HSK 7-9]
    \definition{adj.}{tripartido; que constitui uma relação tripartite}
    \definition[个,些]{s.}{triângulo; coisas triangulares | trigonometria, abreviação de 三角学}
  \seealsoref{三角学}{san1jiao3xue2}
  \end{Phonetics}
\end{Entry}

\begin{Entry}{三角学}{3,7,8}{⼀,⾓,⼦}
  \begin{Phonetics}{三角学}{san1jiao3xue2}
    \definition{s.}{trigonometria; um ramo da matemática que estuda principalmente as funções trigonométricas e suas propriedades, bem como suas aplicações em geometria}[我三角学学得很好。===Sou muito bom em trigonometria.]
  \end{Phonetics}
\end{Entry}

\begin{Entry}{三角恋爱}{3,7,10,10}{⼀,⾓,⼼,⽖}
  \begin{Phonetics}{三角恋爱}{san1jiao3lian4'ai4}
    \definition[场]{s.}{triângulo amoroso | triângulo eterno}
  \end{Phonetics}
\end{Entry}

\begin{Entry}{三明治}{3,8,8}{⼀,⽇,⽔}
  \begin{Phonetics}{三明治}{san1ming2zhi4}[][HSK 6]
    \definition[个,些,块]{s.}{Empréstimo linguístico: sanduíche, \emph{sandwich}}
  \end{Phonetics}
\end{Entry}

\begin{Entry}{三轮车}{3,8,4}{⼀,⾞,⾞}
  \begin{Phonetics}{三轮车}{san1lun2che1}
    \definition[辆,群,排]{s.}{triciclo; triciclo de passageiros}
  \end{Phonetics}
\end{Entry}

\begin{Entry}{三维}{3,11}{⼀,⽷}
  \begin{Phonetics}{三维}{san1wei2}[][HSK 7-9]
    \definition{s.}{três dimensões; 3D; tridimensional}[我们生活在三维空间。===Vivemos em um espaço tridimensional.]
  \end{Phonetics}
\end{Entry}

\begin{Entry}{三番五次}{3,12,4,6}{⼀,⽥,⼆,⽋}
  \begin{Phonetics}{三番五次}{san1fan1-wu3ci4}[][HSK 7-9]
    \definition{expr.}{repetidamente; de novo e de novo; várias e várias vezes; diversas vezes}
  \synonymref{翻来覆去}{fan1lai2-fu4qu4}
  \synonymref{接二连三}{jie1'er4-lian2san1}
  \end{Phonetics}
\end{Entry}

%%%%%%%%%% 上 %%%%%%%%%%
\subsection*{上}\addcontentsline{loh}{figure}{上}

\begin{Entry}{上}{3}{⼀}
  \begin{Phonetics}{上}{shang3}
    \definition{s.}{tom descendente"-ascendente; significa o segundo tom dos quatro tons do mandarim, e também se refere ao terceiro tom do mandarim padrão}
  \end{Phonetics}
  \begin{Phonetics}{上}{shang4}[][HSK 1]
    \definition{adj.}{mais recente; último; anterior; tempo ou a sequência anterior | superior; mais alto; melhor; indica uma posição elevada em termos de qualidade, nível, etc. | lugar elevado; posição superior}
    \definition{s.}{superior; acima; para cima; um lugar alto ou mais alto do que um determinado local | na superfície de um objeto; usado após um substantivo, indica a superfície de um objeto | indica estar dentro do escopo de algo; usado após um substantivo, indica que algo está dentro do âmbito de determinada coisa | indica um aspecto específico | antigamente, referia"-se ao imperador | usado após palavras que indicam idade, equivale a ``\dots 的时候'' | o primeiro nível da escala da música folclórica chinesa, usado como um símbolo de nota na notação musical, equivalente ao ``1'' na notação simplificada.}
    \definition{v.}{subir; montar; embarcar; entrar | ir para; partir para | estar ocupado (com trabalho, estudos, etc.) em um horário fixo; começar a trabalhar ou estudar na hora marcada, etc. | seguir em frente; prosseguir | encher; abastecer; servir; melhorar; aumentar | aparecer no palco; entrar | colocar algo em posição; ajustar; fixar; montar as duas partes de algo | aplicar; pintar; espalhar | ser registrado; ser publicado (em uma publicação) | atingir; ser suficiente (uma determinada quantidade ou grau) | submeter; enviar; apresentar; submeter à aprovação superior | ventilar; apertar; torcer | trazer; servir; colocar comida, pratos, chá e outras coisas na mesa para os convidados | indicar que uma ação tem um resultado | pesquisar na \emph{Internet} | emaranhar"-se; ficar emaranhado; enredar"-se}
    \definition{v.aux.}{usado após um verbo para indicar início e continuidade}
  \seealsoref{的时候}{de5 shi2hou4}
  \seealsoref{里}{li3}
  \seealsoref{上面}{shang4mian5}
  \seealsoref{中}{zhong4}
  \synonymref{安}{an1}
  \synonymref{出}{chu1}
  \synonymref{登}{deng1}
  \synonymref{里}{li3}
  \synonymref{抹}{ma1}
  \synonymref{上面}{shang4mian5}
  \synonymref{涂}{tu2}
  \synonymref{中}{zhong1}
  \synonymref{装}{zhuang1}
  \antonymref{下}{xia4}
  \end{Phonetics}
\end{Entry}

\begin{Entry}{上下}{3,3}{⼀,⼀}
  \begin{Phonetics}{上下}{shang4xia4}[][HSK 5]
    \definition{adv.}{para cima e para baixo}
    \definition[顶]{s.}{alto e baixo | de cima para baixo; para cima e para baixo | superioridade ou inferioridade relativa | (após números redondos) aproximadamente; mais ou menos; por aí | velhos e jovens; hierarquia em termos de cargo e posição social}
    \definition{v.}{subir ou descer | subir e descer; da alta para a baixa ou da baixa para a alta}
  \synonymref{高低}{gao1di1}
  \antonymref{昼夜}{zhou4ye4}
  \antonymref{左右}{zuo3you4}
  \end{Phonetics}
\end{Entry}

\begin{Entry}{上个月}{3,3,4}{⼀,⼈,⽉}
  \begin{Phonetics}{上个月}{shang4ge4yue4}[][HSK 4]
    \definition{s.}{mês passado; refere"-se à hora de um mês atrás, ou seja, o mês passado mais próximo da hora atual}
  \end{Phonetics}
\end{Entry}

\begin{Entry}{上门}{3,3}{⼀,⾨}
  \begin{Phonetics}{上门}{shang4men2}[][HSK 4]
    \definition{v.}{chamar; visitar; aparecer; ir ou vir para ver alguém; ir até a porta; ir até a casa de alguém | trancar a porta; fechar a porta durante a noite | casar"-se e morar com a família da noiva}
  \synonymref{拜访}{bai4fang3}
  \end{Phonetics}
\end{Entry}

\begin{Entry}{上元节}{3,4,5}{⼀,⼉,⾋}
  \begin{Phonetics}{上元节}{shang4yuan2jie2}
    \definition*{s.}{Festival das Lanternas (o décimo quinto dia do primeiro mês do calendário lunar)}
  \end{Phonetics}
\end{Entry}

\begin{Entry}{上升}{3,4}{⼀,⼗}
  \begin{Phonetics}{上升}{shang4sheng1}[][HSK 3]
    \definition{v.}{elevar; subir; mover"-se para cima; mover de baixo para cima; aumentar em nível, grau, quantidade, etc.}
  \synonymref{上涨}{shang4zhang3}
  \antonymref{回升}{hui2sheng1}
  \antonymref{降落}{jiang4luo4}
  \antonymref{落下}{luo4xia4}
  \antonymref{下降}{xia4jiang4}
  \antonymref{下落}{xia4luo4}
  \end{Phonetics}
\end{Entry}

\begin{Entry}{上午}{3,4}{⼀,⼗}
  \begin{Phonetics}{上午}{shang4wu3}[][HSK 1]
    \definition[个]{s.}{manhã; \emph{ante meridiem} (a.m.); geralmente refere"-se ao período entre a manhã e o meio"-dia}
  \antonymref{下午}{xia4wu3}
  \end{Phonetics}
\end{Entry}

\begin{Entry}{上方}{3,4}{⼀,⽅}
  \begin{Phonetics}{上方}{shang4fang1}[][HSK 7-9]
    \definition{s.}{acima; sobre; em cima de | superjacente}
  \seealsoref{下方}{xia4fang1}
  \end{Phonetics}
\end{Entry}

\begin{Entry}{上火}{3,4}{⼀,⽕}
  \begin{Phonetics}{上火}{shang4/huo3}[][HSK 7-9]
    \definition{v.+compl.}{ter dor de garganta | ter excesso de calor interno; na medicina tradicional chinesa, sintomas como prisão de ventre ou inflamação da mucosa nasal, da mucosa oral ou da conjuntiva são classificados como ``calor interno'' | ficar com raiva}
  \seealsoref{上火儿}{shang4huo3r5}
  \end{Phonetics}
\end{Entry}

\begin{Entry}{上火儿}{3,4,2}{⼀,⽕,⼉}
  \begin{Phonetics}{上火儿}{shang4huo3r5}
    \definition{v.}{Dialeto: ficar com raiva; explodir}
  \end{Phonetics}
\end{Entry}

\begin{Entry}{上车}{3,4}{⼀,⾞}
  \begin{Phonetics}{上车}{shang4che1}[][HSK 1]
    \definition{v.}{entrar; subir (em um ônibus, trem, carro etc.)}
  \antonymref{下车}{xia4che1}
  \end{Phonetics}
\end{Entry}

\begin{Entry}{上去}{3,5}{⼀,⼛}
  \begin{Phonetics}{上去}{shang4qu5}[][HSK 3]
    \definition{v.}{subir (a partir da minha localização) | ascender a um lugar (ou estado) considerado mais elevado (ou acima); usado depois de um verbo para indicar movimento, de baixo para cima ou de perto para longe}
  \synonymref{上来}{shang4lai5}
  \antonymref{下来}{xia4lai5}
  \end{Phonetics}
\end{Entry}

\begin{Entry}{上古}{3,5}{⼀,⼝}
  \begin{Phonetics}{上古}{shang4gu3}
    \definition{s.}{tempos antigos; eras remotas | antiguidade | tempos históricos antigos | o passado distante}
  \antonymref{现代}{xian4dai4}
  \end{Phonetics}
\end{Entry}

\begin{Entry}{上台}{3,5}{⼀,⼝}
  \begin{Phonetics}{上台}{shang4tai2}[][HSK 6]
    \definition{v.}{aparecer no palco; subir na plataforma; ir para o palco ou pódio | assumir o poder; chegar (subir) ao poder; começar a assumir papéis de liderança ou a ganhar algum tipo de poder}
  \synonymref{上任}{shang4/ren4}
  \antonymref{下台}{xia4/tai2}
  \end{Phonetics}
\end{Entry}

\begin{Entry}{上司}{3,5}{⼀,⼝}
  \begin{Phonetics}{上司}{shang4si5}[][HSK 7-9]
    \definition[位,名,个]{s.}{chefe; superior}
  \synonymref{上级}{shang4ji2}
  \antonymref{部属}{bu4shu3}
  \antonymref{部下}{bu4xia4}
  \antonymref{下级}{xia4ji2}
  \antonymref{下属}{xia4shu3}
  \end{Phonetics}
\end{Entry}

\begin{Entry}{上头}{3,5}{⼀,⼤}
  \begin{Phonetics}{上头}{shang4tou2}
    \definition{v.}{(álcool, amor etc.) subir à cabeça; (uma ideia, uma música etc.) entrar na cabeça de alguém; capturar a atenção de alguém | Obsoleto: (uma garota no dia do seu casamento) começar a usar o cabelo preso em um coque (em vez de uma trança)}
  \synonymref{方面}{fang1mian4}
  \synonymref{上面}{shang4mian5}
  \end{Phonetics}
  \begin{Phonetics}{上头}{shang4tou5}[][HSK 7-9]
    \definition{s.}{acima; em cima de; na superfície de; superior}
  \end{Phonetics}
\end{Entry}

\begin{Entry}{上市}{3,5}{⼀,⼱}
  \begin{Phonetics}{上市}{shang4shi4}[][HSK 6]
    \definition{v.}{listar; abrir o capital; ser listado (na bolsa de valores) | estar na estação; estar (aparecer) no mercado | ir ao mercado (para fazer compras)}
  \synonymref{问世}{wen4shi4}
  \end{Phonetics}
\end{Entry}

\begin{Entry}{上边}{3,5}{⼀,⾡}
  \begin{Phonetics}{上边}{shang4bian5}[][HSK 1]
    \definition{s.}{topo; acima; sobre; superior}
  \synonymref{领导}{ling3dao3}
  \synonymref{上方}{shang4fang1}
  \synonymref{上级}{shang4ji2}
  \synonymref{上头}{shang4tou5}
  \synonymref{上面}{shang4mian5}
  \synonymref{上司}{shang4si5}
  \antonymref{底下}{di3xia5}
  \antonymref{下边}{xia4bian5}
  \antonymref{下方}{xia4fang1}
  \antonymref{下级}{xia4ji2}
  \end{Phonetics}
\end{Entry}

\begin{Entry}{上任}{3,6}{⼀,⼈}
  \begin{Phonetics}{上任}{shang4/ren4}[][HSK 7-9]
    \definition[班]{s.}{predecessor; ex"-funcionário}
    \definition{v.+compl.}{assumir o cargo; ocupar um cargo oficial; refere"-se à posse de autoridades}
  \synonymref{出任}{chu1ren4}
  \synonymref{就职}{jiu4/zhi2}
  \synonymref{就任}{jiu4ren4}
  \synonymref{上台}{shang4tai2}
  \antonymref{辞职}{ci2/zhi2}
  \antonymref{离职}{li2/zhi2}
  \end{Phonetics}
\end{Entry}

\begin{Entry}{上场}{3,6}{⼀,⼟}
  \begin{Phonetics}{上场}{shang4/chang3}[][HSK 7-9]
    \definition{v.+compl.}{Esporte: entrar na quadra (ou campo); participar de uma competição | aparecer no palco; subir no palco; entrar em cena}
  \antonymref{退场}{tui4chang3}
  \antonymref{下场}{xia4chang3}
  \end{Phonetics}
\end{Entry}

\begin{Entry}{上当}{3,6}{⼀,⼹}
  \begin{Phonetics}{上当}{shang4/dang4}[][HSK 6]
    \definition{v.+compl.}{ser enganado; ser ludibriado; morder a isca; cair nas mãos de alguém}
  \synonymref{受骗}{shou4/pian4}
  \antonymref{精明}{jing1ming2}
  \antonymref{警惕}{jing3ti4}
  \end{Phonetics}
\end{Entry}

\begin{Entry}{上旬}{3,6}{⼀,⽇}
  \begin{Phonetics}{上旬}{shang4xun2}[][HSK 7-9]
    \definition{s.}{primeiro terço do mês; os dez dias do dia 1 ao dia 10 de cada mês}
  \seealsoref{中旬}{zhong1xun2}
  \antonymref{下旬}{xia4xun2}
  \end{Phonetics}
\end{Entry}

\begin{Entry}{上次}{3,6}{⼀,⽋}
  \begin{Phonetics}{上次}{shang4ci4}[][HSK 1]
    \definition{adv.}{última vez}
  \synonymref{此前}{ci3qian2}
  \synonymref{先前}{xian1qian2}
  \antonymref{下次}{xia4ci4}
  \end{Phonetics}
\end{Entry}

\begin{Entry}{上级}{3,6}{⼀,⽷}
  \begin{Phonetics}{上级}{shang4ji2}[][HSK 5]
    \definition[个,位]{s.}{nível superior; organização ou pessoa em nível superior; organizações ou pessoas de nível superior dentro do mesmo sistema organizacional}
  \synonymref{部属}{bu4shu3}
  \synonymref{上司}{shang4si5}
  \antonymref{部下}{bu4xia4}
  \antonymref{下级}{xia4ji2}
  \end{Phonetics}
\end{Entry}

\begin{Entry}{上网}{3,6}{⼀,⽹}
  \begin{Phonetics}{上网}{shang4/wang3}[][HSK 1]
    \definition{v.+compl.}{conectar"-se à \emph{Internet}; acessar a \emph{Internet}; entrar na \emph{Internet}; acessar a rede; refere"-se especificamente ao computador do usuário conectado à \emph{Internet} para pesquisar e consultar informações, etc.}
  \end{Phonetics}
\end{Entry}

\begin{Entry}{上衣}{3,6}{⼀,⾐}
  \begin{Phonetics}{上衣}{shang4yi1}[][HSK 3]
    \definition[件]{s.}{jaqueta; roupas para a parte superior do corpo}
  \end{Phonetics}
\end{Entry}

\begin{Entry}{上访}{3,6}{⼀,⾔}
  \begin{Phonetics}{上访}{shang4fang3}
    \definition{v.}{apresentar recursos a autoridades superiores em busca de ajuda | solicitar uma audiência com autoridades superiores (especialmente funcionários do governo) para fazer uma petição}
  \antonymref{放任}{fang4ren4}
  \antonymref{忍耐}{ren3nai4}
  \end{Phonetics}
\end{Entry}

\begin{Entry}{上声}{3,7}{⼀,⼠}
  \begin{Phonetics}{上声}{shang3 sheng1}
    \definition{s.}{tom descendente e ascendente | terceiro tom no mandarim moderno}
  \end{Phonetics}
  \begin{Phonetics}{上声}{shang4 sheng1}
    \definition{s.}{tom ascendente (o segundo tom no chinês antigo) | terceiro tom (no chinês moderno)}
  \end{Phonetics}
\end{Entry}

\begin{Entry}{上岗}{3,7}{⼀,⼭}
  \begin{Phonetics}{上岗}{shang4/gang3}[][HSK 7-9]
    \definition{v.+compl.}{estar em período probatório | começar a trabalhar; assumir um cargo}
  \end{Phonetics}
\end{Entry}

\begin{Entry}{上报}{3,7}{⼀,⼿}
  \begin{Phonetics}{上报}{shang4bao4}[][HSK 7-9]
    \definition{v.}{aparecer nos jornais; ser publicado | reportar a um órgão superior; reportar à liderança; reportar"-se aos superiores}
  \synonymref{报到}{bao4/dao4}
  \end{Phonetics}
\end{Entry}

\begin{Entry}{上来}{3,7}{⼀,⽊}
  \begin{Phonetics}{上来}{shang4lai5}[][HSK 3]
    \definition{v.}{subir (para a minha localização) | estar no começo; começar; iniciar | surgir; de um lugar baixo para um lugar alto (o interlocutor está em um lugar alto) | usado após o verbo, indica que algo foi concluído com sucesso}
  \synonymref{上去}{shang4qu5}
  \antonymref{下去}{xia4qu5}
  \end{Phonetics}
\end{Entry}

\begin{Entry}{上诉}{3,7}{⼀,⾔}
  \begin{Phonetics}{上诉}{shang4su4}[][HSK 7-9]
    \definition{s.}{apelação (para um tribunal superior)}
    \definition{v.}{apresentar um recurso; instaurar um recurso}
  \end{Phonetics}
\end{Entry}

\begin{Entry}{上周}{3,8}{⼀,⼝}
  \begin{Phonetics}{上周}{shang4zhou1}[][HSK 2]
    \definition{s.}{semana passada}
  \antonymref{下周}{xia4zhou1}
  \end{Phonetics}
\end{Entry}

\begin{Entry}{上坡路}{3,8,13}{⼀,⼟,⾜}
  \begin{Phonetics}{上坡路}{shang4po1lu4}
    \definition[条,些]{s.}{estrada em subida; aclive; subida | Figurativo: tendência ascendente; progresso constante}
  \end{Phonetics}
\end{Entry}

\begin{Entry}{上学}{3,8}{⼀,⼦}
  \begin{Phonetics}{上学}{shang4xue2}[][HSK 1]
    \definition{v.}{ir à escola; frequentar a escola; estar na escola; ir à escola para estudar | começar a escola; começar a estudar no ensino fundamental}
  \seealsoref{上课}{shang4/ke4}
  \synonymref{读书}{du2/shu1}
  \synonymref{入学}{ru4/xue2}
  \synonymref{上课}{shang4/ke4}
  \antonymref{放学}{fang4/xue2}
  \end{Phonetics}
\end{Entry}

\begin{Entry}{上弦}{3,8}{⼀,⼸}
  \begin{Phonetics}{上弦}{shang4xian2}
    \definition{s.}{quarto crescente (da lua) (oposto à lua minguante)}
    \definition{v.}{apertar a mola de; dar corda (até); dar corda a um relógio}
  \antonymref{下弦}{xia4xian2}
  \end{Phonetics}
\end{Entry}

\begin{Entry}{上空}{3,8}{⼀,⽳}
  \begin{Phonetics}{上空}{shang4kong1}[][HSK 7-9]
    \definition{s.}{no céu; acima da cabeça; no alto, no ar}
  \end{Phonetics}
\end{Entry}

\begin{Entry}{上述}{3,8}{⼀,⾡}
  \begin{Phonetics}{上述}{shang4shu4}[][HSK 7-9]
    \definition{adj.}{mencionado anteriormente; supracitado; acima citado; conforme dito ou narrado acima}
  \synonymref{描述}{miao2shu4}
  \synonymref{摸索}{mo1suo3}
  \synonymref{试验}{shi4yan4}
  \end{Phonetics}
\end{Entry}

\begin{Entry}{上限}{3,8}{⼀,⾩}
  \begin{Phonetics}{上限}{shang4xian4}[][HSK 7-9]
    \definition[个]{s.}{teto; limite superior; refere"-se ao primeiro ou mais alto limite dentro de um determinado conjunto de limites}
  \antonymref{下限}{xia4xian4}
  \end{Phonetics}
\end{Entry}

\begin{Entry}{上帝}{3,9}{⼀,⼱}
  \begin{Phonetics}{上帝}{shang4di4}[][HSK 6]
    \definition*{s.}{Deus; O Deus Supremo no Cristianismo | O Imperador do Céu; um deus na antiga crença chinesa que pode controlar tudo no mundo}
    \definition[个]{s.}{(figurado) cliente; metáfora para consumidores}
  \end{Phonetics}
\end{Entry}

\begin{Entry}{上映}{3,9}{⼀,⽇}
  \begin{Phonetics}{上映}{shang4ying4}[][HSK 7-9]
    \definition{v.}{executar; exibir; mostrar (um filme); (novo filme) lançar para exibição}
  \synonymref{放映}{fang4ying4}
  \end{Phonetics}
\end{Entry}

\begin{Entry}{上面}{3,9}{⼀,⾯}
  \begin{Phonetics}{上面}{shang4mian5}[][HSK 3]
    \definition{s.}{uma posição mais alta que algo; uma posição acima/acima de algo | superfície do objeto | aspecto | a parte acima mencionada; a parte que vem primeiro na ordem; a parte de um artigo ou discurso que vem antes da presente | autoridades superiores | os mais velhos; a geração mais velha da família}
  \synonymref{方面}{fang1mian4}
  \synonymref{上方}{shang4fang1}
  \synonymref{上头}{shang4tou5}
  \antonymref{底下}{di3xia5}
  \antonymref{下面}{xia4mian5}
  \end{Phonetics}
\end{Entry}

\begin{Entry}{上流}{3,10}{⼀,⽔}
  \begin{Phonetics}{上流}{shang4liu2}[][HSK 7-9]
    \definition{adj.}{da classe alta; refinado | pertencente aos círculos superiores; anteriormente, referia"-se a pessoas de alto \emph{status} social}
    \definition{s.}{trecho superior (de um rio); a montante}
  \synonymref{崇高}{chong2gao1}
  \synonymref{高超}{gao1chao1}
  \synonymref{高贵}{gao1gui4}
  \synonymref{高尚}{gao1shang4}
  \end{Phonetics}
\end{Entry}

\begin{Entry}{上海}{3,10}{⼀,⽔}
  \begin{Phonetics}{上海}{shang4hai3}
    \definition*{s.}{Município de Xangai (Shanghai), centro"-leste da China}
  \end{Phonetics}
\end{Entry}

\begin{Entry}{上涨}{3,10}{⼀,⽔}
  \begin{Phonetics}{上涨}{shang4zhang3}[][HSK 5]
    \definition{v.}{subir; ir para cima; ascender}
  \synonymref{高潮}{gao1chao2}
  \synonymref{高涨}{gao1zhang3}
  \synonymref{上升}{shang4sheng1}
  \antonymref{回落}{hui2luo4}
  \antonymref{下跌}{xia4die1}
  \antonymref{下降}{xia4jiang4}
  \antonymref{下落}{xia4luo4}
  \end{Phonetics}
\end{Entry}

\begin{Entry}{上班}{3,10}{⼀,⽟}
  \begin{Phonetics}{上班}{shang4/ban1}[][HSK 1]
    \definition{v.+compl.}{ir trabalhar; começar a trabalhar; estar de plantão; ir trabalhar no local de trabalho regular no horário especificado}
  \synonymref{工作}{gong1zuo4}
  \antonymref{下班}{xia4/ban1}
  \antonymref{休息}{xiu1xi5}
  \end{Phonetics}
\end{Entry}

\begin{Entry}{上班族}{3,10,11}{⼀,⽟,⽅}
  \begin{Phonetics}{上班族}{shang4ban1zu2}
    \definition[本]{s.}{trabalhadores de escritório (como grupo social)}
  \end{Phonetics}
\end{Entry}

\begin{Entry}{上课}{3,10}{⼀,⾔}
  \begin{Phonetics}{上课}{shang4/ke4}[][HSK 1]
    \definition{v.+compl.}{frequentar aulas; ir às aulas; dar uma aula}
  \seealsoref{上学}{shang4xue2}
  \synonymref{讲课}{jiang3ke4}
  \synonymref{教学}{jiao4xue2}
  \synonymref{上学}{shang4xue2}
  \antonymref{下课}{xia4/ke4}
  \end{Phonetics}
\end{Entry}

\begin{Entry}{上调}{3,10}{⼀,⾔}
  \begin{Phonetics}{上调}{shang4tiao2}[][HSK 7-9]
    \definition{v.}{transferir (alguém) para um cargo de nível superior | transferir bens, fundos, etc. para uma unidade de nível superior | ajustar para cima | aumentar (os preços)}
  \antonymref{下调}{xia4tiao2}
  \end{Phonetics}
\end{Entry}

\begin{Entry}{上期}{3,12}{⼀,⽉}
  \begin{Phonetics}{上期}{shang4qi1}[][HSK 7-9]
    \definition{s.}{período anterior}
  \end{Phonetics}
\end{Entry}

\begin{Entry}{上游}{3,12}{⼀,⽔}
  \begin{Phonetics}{上游}{shang4you2}[][HSK 7-9]
    \definition{s.}{a montante; trecho superior de um rio; o trecho de um rio próximo à sua nascente; também se refere à área por onde esse trecho flui | posição avançada; metaforicamente, refere"-se a um \emph{status} ou nível avançado}
  \synonymref{领先}{ling3/xian1}
  \synonymref{先进}{xian1jin4}
  \antonymref{下游}{xia4you2}
  \end{Phonetics}
\end{Entry}

\begin{Entry}{上楼}{3,13}{⼀,⽊}
  \begin{Phonetics}{上楼}{shang4lou2}[][HSK 4]
    \definition{v.}{subir as escadas; ir para o andar de cima}
  \end{Phonetics}
\end{Entry}

\begin{Entry}{上演}{3,14}{⼀,⽔}
  \begin{Phonetics}{上演}{shang4yan3}[][HSK 6]
    \definition{s.}{exibição | encenação}
    \definition{v.}{exibir (um filme); encenar (uma peça); atuar; colocar no palco}
  \synonymref{演出}{yan3chu1}
  \end{Phonetics}
\end{Entry}

\begin{Entry}{上瘾}{3,16}{⼀,⽧}
  \begin{Phonetics}{上瘾}{shang4/yin3}[][HSK 7-9]
    \definition{v.+compl.}{ser viciado (em algo); adquirir o hábito (de fazer algo); gostar muito de algo, a ponto de não conseguir viver sem}
  \synonymref{沉迷}{chen2mi2}
  \synonymref{痴迷}{chi1mi2}
  \synonymref{陶醉}{tao2zui4}
  \antonymref{厌倦}{yan4juan4}
  \end{Phonetics}
\end{Entry}

%%%%%%%%%% 下 %%%%%%%%%%
\subsection*{下}\addcontentsline{loh}{figure}{下}

\begin{Entry}{下}{3}{⼀}
  \begin{Phonetics}{下}{xia4}[][HSK 1,2]
    \definition{clas.}{número de vezes usado para a ação | volume de um contêiner; quantidade de objetos que cabem em um utensílio | usado depois de 两 e 几 para expressar habilidade, capacidade, destreza}
    \definition{s.}{abaixo | próximo; último; segundo; referindo"-se ao que está por vir ou ao que vem em seguida | mais baixo; inferior; de baixo nível ou grau | próximo; último; segundo; em ordem ou em ordem cronológica | indica pertencer a uma determinada faixa, situação, condição, etc. | indica uma determinada época ou estação | usado após um número para indicar posição ou direção | para baixo (após uma preposição) | sob (depois de um substantivo) | para baixo (antes de um verbo)}
    \definition{v.}{desembarcar; descer; sair | cair (chuva, neve, etc.) | enviar; emitir; entregar | ir para | sair; partir; retirar"-se | lançar; colocar | descarregar; desmontar; tirar (fora) | formar (uma opinião, ideia, etc.); tomar decisões, fazer julgamentos, etc. | usar; aplicar | dar à luz (animais) | tomar; capturar; conquistar | ceder | terminar; deixar de lado; terminar o trabalho ou os estudos diários na hora prevista | para negação; ser inferior a; ser menor que}
  \seealsoref{几}{ji3}
  \seealsoref{两}{liang3}
  \seealsoref{下边}{xia4bian5}
  \seealsoref{下面}{xia4mian5}
  \synonymref{差}{chai4}
  \synonymref{次}{ci4}
  \synonymref{低}{di1}
  \synonymref{放}{fang4}
  \synonymref{攻}{gong1}
  \synonymref{克}{ke4}
  \synonymref{取}{qu3}
  \synonymref{少}{shao3}
  \synonymref{退}{tui4}
  \synonymref{卸}{xie4}
  \antonymref{安}{an1}
  \antonymref{多}{duo1}
  \antonymref{高}{gao1}
  \antonymref{好}{hao3}
  \antonymref{进}{jin4}
  \antonymref{良}{liang2}
  \antonymref{拿}{na2}
  \antonymref{上}{shang4}
  \antonymref{优}{you1}
  \antonymref{装}{zhuang1}
  \end{Phonetics}
\end{Entry}

\begin{Entry}{下一代}{3,1,5}{⼀,⼀,⼈}
  \begin{Phonetics}{下一代}{xia4yi2dai4}[][HSK 7-9]
    \definition{s.}{próxima geração | geração mais jovem}[《星际迷航: 下一代》===Star Trek: A Nova Geração | 关心下一代委员会。===Comitê para o Cuidado da Próxima Geração.]
  \end{Phonetics}
\end{Entry}

\begin{Entry}{下个月}{3,3,4}{⼀,⼈,⽉}
  \begin{Phonetics}{下个月}{xia4ge4yue4}[][HSK 4]
    \definition{s.}{próximo mês; mês que vem; refere"-se ao próximo mês do mês atual}
  \end{Phonetics}
\end{Entry}

\begin{Entry}{下乡}{3,3}{⼀,⼄}
  \begin{Phonetics}{下乡}{xia4/xiang1}[][HSK 7-9]
    \definition{v.+compl.}{ir para o campo, zona rural (especialmente para intelectuais)}
  \end{Phonetics}
\end{Entry}

\begin{Entry}{下山}{3,3}{⼀,⼭}
  \begin{Phonetics}{下山}{xia4/shan1}[][HSK 7-9]
    \definition{v.+compl.}{descer uma colina ou montanha | (Sol) pôr"-se; desaparecer abaixo do horizonte}
  \end{Phonetics}
\end{Entry}

\begin{Entry}{下午}{3,4}{⼀,⼗}
  \begin{Phonetics}{下午}{xia4wu3}[][HSK 1]
    \definition[个]{s.}{tarde; \emph{post meridiem} (p.m.); refere"-se ao período entre o meio"-dia e o pôr do sol}
  \synonymref{午后}{wu3hou4}
  \antonymref{上午}{shang4wu3}
  \end{Phonetics}
\end{Entry}

\begin{Entry}{下午茶}{3,4,9}{⼀,⼗,⾋}
  \begin{Phonetics}{下午茶}{xia4wu3cha2}
    \definition[杯]{s.}{chá da tarde (normalmente chás com doces)}
  \end{Phonetics}
\end{Entry}

\begin{Entry}{下巴}{3,4}{⼀,⼰}
  \begin{Phonetics}{下巴}{xia4ba5}
    \definition[个]{s.}{queixo | mandíbula inferior}
  \end{Phonetics}
\end{Entry}

\begin{Entry}{下手}{3,4}{⼀,⼿}
  \begin{Phonetics}{下手}{xia4/shou3}[][HSK 7-9]
    \definition{v.+compl.}{colocar a mão na massa (em uma tarefa, etc.); começar a fazer algo; empreender | fazer algo ruim a alguém}
  \synonymref{出手}{chu1/shou3}
  \synonymref{动手}{dong4/shou3}
  \synonymref{开始}{kai1shi3}
  \synonymref{入手}{ru4shou3}
  \synonymref{着手}{zhuo2shou3}
  \end{Phonetics}
\end{Entry}

\begin{Entry}{下方}{3,4}{⼀,⽅}
  \begin{Phonetics}{下方}{xia4fang1}
    \definition{s.}{parte inferior | abaixo | embaixo | mundo dos mortais}
    \definition{v.}{descer ao mundo dos mortais (deuses)}
  \seealsoref{上方}{shang4fang1}
  \antonymref{上方}{shang4fang1}
  \end{Phonetics}
\end{Entry}

\begin{Entry}{下水道}{3,4,12}{⼀,⽔,⾡}
  \begin{Phonetics}{下水道}{xia4shui3dao4}
    \definition{s.}{esgoto; sarjeta; ralo}
  \end{Phonetics}
\end{Entry}

\begin{Entry}{下车}{3,4}{⼀,⾞}
  \begin{Phonetics}{下车}{xia4che1}[][HSK 1]
    \definition{v.}{descer ou sair de (um ônibus, trem, carro etc.)}
  \antonymref{上车}{shang4che1}
  \end{Phonetics}
\end{Entry}

\begin{Entry}{下令}{3,5}{⼀,⼈}
  \begin{Phonetics}{下令}{xia4/ling4}[][HSK 7-9]
    \definition{v.+compl.}{dar ordens; ordenar}
  \synonymref{命令}{ming4ling4}
  \end{Phonetics}
\end{Entry}

\begin{Entry}{下功夫}{3,5,4}{⼀,⼒,⼤}
  \begin{Phonetics}{下功夫}{xia4gong1fu5}[][HSK 7-9]
    \definition{v.}{concentrar esforços; investir tempo e energia em algo para alcançar bons resultados}
  \end{Phonetics}
\end{Entry}

\begin{Entry}{下去}{3,5}{⼀,⼛}
  \begin{Phonetics}{下去}{xia4qu5}[][HSK 3]
    \definition{part.}{usado depois de verbos para indicar de alto a baixo | usado depois de um verbo para indicar continuação}
    \definition{v.}{descer; baixar (a partir da minha localização) | (após um verbo) continuar (fazendo algo); prosseguir | usado após o verbo, indica uma descida de um ponto alto para um ponto baixo | usado após o verbo, indica continuidade | usado após um adjetivo, indica que o grau continua aumentando}
  \synonymref{出来}{chu1lai5}
  \synonymref{下来}{xia4lai5}
  \antonymref{上来}{shang4lai5}
  \end{Phonetics}
\end{Entry}

\begin{Entry}{下台}{3,5}{⼀,⼝}
  \begin{Phonetics}{下台}{xia4/tai2}[][HSK 7-9]
    \definition{v.+compl.}{descer do palco, pódio ou da plataforma | perder o poder; deixar o cargo | (geralmente na forma negativa) sair de uma situação difícil ou embaraçosa; retirar"-se com dignidade; sair de fininho; saída elegante | renunciar; renunciar ao próprio direito (poder); perder o poder; livrar"-se de uma posição}
  \antonymref{上台}{shang4tai2}
  \end{Phonetics}
\end{Entry}

\begin{Entry}{下边}{3,5}{⼀,⾡}
  \begin{Phonetics}{下边}{xia4bian5}[][HSK 1]
    \definition{s.}{abaixo; sob; por baixo | próximo em ordem; seguinte | nível inferior; subordinado | a parte inferior}
  \seealsoref{下}{xia4}
  \synonymref{底下}{di3xia5}
  \synonymref{下}{xia4}
  \synonymref{下方}{xia4fang1}
  \synonymref{下面}{xia4mian5}
  \antonymref{上边}{shang4bian5}
  \antonymref{上方}{shang4fang1}
  \antonymref{上面}{shang4mian5}
  \end{Phonetics}
\end{Entry}

\begin{Entry}{下决心}{3,6,4}{⼀,⼎,⼼}
  \begin{Phonetics}{下决心}{xia4jue2xin1}[][HSK 7-9]
    \definition{v.}{resolver; determinar; tomar uma decisão}
  \end{Phonetics}
\end{Entry}

\begin{Entry}{下场}{3,6}{⼀,⼟}
  \begin{Phonetics}{下场}{xia4chang3}[][HSK 7-9]
    \definition[个,种]{s.}{final ruim; estado lamentável; destino}
    \definition{v.}{sair de cena; encerrar a noite}
  \seealsoref{结局}{jie2ju2}
  \synonymref{结局}{jie2ju2}
  \synonymref{结束}{jie2shu4}
  \antonymref{上场}{shang4/chang3}
  \end{Phonetics}
\end{Entry}

\begin{Entry}{下旬}{3,6}{⼀,⽇}
  \begin{Phonetics}{下旬}{xia4xun2}[][HSK 7-9]
    \definition[月]{s.}{última dezena do mês; último período de dez dias de um mês; do dia 21 até o final de cada mês}
  \seealsoref{中旬}{zhong1xun2}
  \antonymref{上旬}{shang4xun2}
  \end{Phonetics}
\end{Entry}

\begin{Entry}{下次}{3,6}{⼀,⽋}
  \begin{Phonetics}{下次}{xia4ci4}[][HSK 1]
    \definition{s.}{na próxima vez; na próxima oportunidade ou no próximo evento}
  \synonymref{再次}{zai4ci4}
  \end{Phonetics}
\end{Entry}

\begin{Entry}{下级}{3,6}{⼀,⽷}
  \begin{Phonetics}{下级}{xia4ji2}[][HSK 7-9]
    \definition{s.}{classificação inferior | subordinado | nível inferior (posição)}
  \synonymref{部下}{bu4xia4}
  \synonymref{下属}{xia4shu3}
  \antonymref{上级}{shang4ji2}
  \antonymref{上司}{shang4si5}
  \end{Phonetics}
\end{Entry}

\begin{Entry}{下坠}{3,7}{⼀,⼟}
  \begin{Phonetics}{下坠}{xia4zhui4}[][HSK 7-9]
    \definition{v.}{cair; deixar cair | Medicina: fazer força (ao defecar); ter urgência (na gravidez); ter tenesmo}
  \synonymref{笨重}{ben4zhong4}
  \synonymref{下降}{xia4jiang4}
  \synonymref{坠落}{zhui4luo4}
  \antonymref{飞跃}{fei1yue4}
  \end{Phonetics}
\end{Entry}

\begin{Entry}{下岗}{3,7}{⼀,⼭}
  \begin{Phonetics}{下岗}{xia4/gang3}[][HSK 7-9]
    \definition{v.+compl.}{ser despedido; perder o emprego; (funcionários) que deixam seus cargos devido à falência da empresa, redução de pessoal ou outros motivos | sair de serviço}
  \antonymref{全职}{quan2zhi2}
  \end{Phonetics}
\end{Entry}

\begin{Entry}{下来}{3,7}{⼀,⽊}
  \begin{Phonetics}{下来}{xia4lai5}[][HSK 3]
    \definition{part.}{usado após o verbo, indica que a ação ou o comportamento se dirige para a posição do falante ou que a ação é contínua ou concluída | usado após um adjetivo, indica que uma determinada situação começou a ocorrer e continuará a se desenvolver}
    \definition{v.}{descer (para a minha localização) | (colheitas/frutas/vegetais, etc.) ser colhido; estar maduro o suficiente para ser colhido | (período de tempo) acabar; passar; chegar ao fim; indicar o fim de um período de tempo}
  \synonymref{下去}{xia4qu5}
  \antonymref{上去}{shang4qu5}
  \end{Phonetics}
\end{Entry}

\begin{Entry}{下周}{3,8}{⼀,⼝}
  \begin{Phonetics}{下周}{xia4zhou1}[][HSK 2]
    \definition{s.}{próxima semana}
  \antonymref{上周}{shang4zhou1}
  \end{Phonetics}
\end{Entry}

\begin{Entry}{下弦}{3,8}{⼀,⼸}
  \begin{Phonetics}{下弦}{xia4xian2}
    \definition{s.}{lua minguante ou em quarto minguante; último (ou terceiro) quarto (da lua)}
  \antonymref{上弦}{shang4xian2}
  \end{Phonetics}
\end{Entry}

\begin{Entry}{下线}{3,8}{⼀,⽷}
  \begin{Phonetics}{下线}{xia4xian4}
    \definition{v.}{sair da sessão; desconectar"-se da \emph{Internet}; refere"-se à suspensão temporária das atividades de comunicação \emph{on"-line}, geralmente significando interromper temporariamente o bate"-papo \emph{on"-line} ou os jogos \emph{on"-line} | sair da linha de produção; isso se refere a automóveis, eletrodomésticos, etc., que foram montados na linha de produção e estão prontos para sair da fábrica}
  \synonymref{离开}{li2/kai1}
  \antonymref{登陆}{deng1/lu4}
  \antonymref{问世}{wen4shi4}
  \antonymref{在线}{zai4xian4}
  \end{Phonetics}
\end{Entry}

\begin{Entry}{下降}{3,8}{⼀,⾩}
  \begin{Phonetics}{下降}{xia4jiang4}[][HSK 4]
    \definition{v.}{cair; despencar; declinar; descer; diminuir; ir para baixo}
  \synonymref{低落}{di1luo4}
  \synonymref{降低}{jiang4di1}
  \synonymref{降落}{jiang4luo4}
  \synonymref{下跌}{xia4die1}
  \synonymref{下落}{xia4luo4}
  \synonymref{消沉}{xiao1chen2}
  \antonymref{回升}{hui2sheng1}
  \antonymref{爬升}{pa2sheng1}
  \antonymref{起飞}{qi3fei1}
  \antonymref{上升}{shang4sheng1}
  \antonymref{上涨}{shang4zhang3}
  \antonymref{增加}{zeng1jia1}
  \antonymref{增长}{zeng1zhang3}
  \end{Phonetics}
\end{Entry}

\begin{Entry}{下限}{3,8}{⼀,⾩}
  \begin{Phonetics}{下限}{xia4xian4}
    \definition{s.}{limite mínimo ou mais recente permitido; limite inferior; limiar; mínimo prescrito; piso (nível)}
  \antonymref{上限}{shang4xian4}
  \end{Phonetics}
\end{Entry}

\begin{Entry}{下雨}{3,8}{⼀,⾬}
  \begin{Phonetics}{下雨}{xia4/yu3}[][HSK 1]
    \definition{v.+compl.}{chover}
  \synonymref{暴雨}{bao4yu3}
  \synonymref{大雨}{da4yu3}
  \end{Phonetics}
\end{Entry}

\begin{Entry}{下面}{3,9}{⼀,⾯}
  \begin{Phonetics}{下面}{xia4mian5}[][HSK 3]
    \definition{s.}{em baixo; abaixo; parte de baixo | próximo; seguinte; a parte posterior; a parte posterior de um artigo ou discurso em relação ao que está sendo narrado no momento | subordinado; o nível inferior; homens nos níveis inferiores | por baixo}
  \synonymref{底下}{di3xia5}
  \synonymref{上方}{shang4fang1}
  \antonymref{上面}{shang4mian5}
  \end{Phonetics}
\end{Entry}

\begin{Entry}{下海}{3,10}{⼀,⽔}
  \begin{Phonetics}{下海}{xia4/hai3}[][HSK 7-9]
    \definition{v.+compl.}{ir pescar no mar; pescar para ganhar a vida; ir para o mar | tornar"-se profissional; isso se refere a atores de ópera amadores que se tornam atores profissionais | deixar o emprego original e abrir o próprio negócio; refere"-se a pessoas que não eram originalmente donas de negócios, mas que passaram a empreender}
  \antonymref{登陆}{deng1/lu4}
  \end{Phonetics}
\end{Entry}

\begin{Entry}{下班}{3,10}{⼀,⽟}
  \begin{Phonetics}{下班}{xia4/ban1}[][HSK 1]
    \definition{v.+compl.}{sair do trabalho; bater ponto; terminar o trabalho na hora prevista e sair do local de trabalho}
  \antonymref{开工}{kai1/gong1}
  \antonymref{上班}{shang4/ban1}
  \end{Phonetics}
\end{Entry}

\begin{Entry}{下课}{3,10}{⼀,⾔}
  \begin{Phonetics}{下课}{xia4/ke4}[][HSK 1]
    \definition{v.+compl.}{terminar a aula; sair da aula}
  \synonymref{下台}{xia4/tai2}
  \antonymref{就任}{jiu4ren4}
  \antonymref{上课}{shang4/ke4}
  \antonymref{上台}{shang4tai2}
  \end{Phonetics}
\end{Entry}

\begin{Entry}{下调}{3,10}{⼀,⾔}
  \begin{Phonetics}{下调}{xia4diao4}
    \definition{v.}{rebaixar; diminuir a regulamentação | passar para uma unidade inferior}
  \end{Phonetics}
  \begin{Phonetics}{下调}{xia4tiao2}[][HSK 7-9]
    \definition{v.}{regular para baixo; ajustar para baixo}
  \end{Phonetics}
\end{Entry}

\begin{Entry}{下载}{3,10}{⼀,⾞}
  \begin{Phonetics}{下载}{xia4zai3}[][HSK 4]
    \definition{v.}{\emph{download}; baixar; salvar informações da \emph{Web} em um dispositivo, como um computador}
  \antonymref{传输}{chuan2shu1}
  \end{Phonetics}
\end{Entry}

\begin{Entry}{下蛋}{3,11}{⼀,⾍}
  \begin{Phonetics}{下蛋}{xia4dan4}
    \definition{v.}{botar ovos}
  \end{Phonetics}
\end{Entry}

\begin{Entry}{下雪}{3,11}{⼀,⾬}
  \begin{Phonetics}{下雪}{xia4/xue3}[][HSK 2]
    \definition{v.+compl.}{nevar}
  \end{Phonetics}
\end{Entry}

\begin{Entry}{下属}{3,12}{⼀,⼫}
  \begin{Phonetics}{下属}{xia4shu3}[][HSK 7-9]
    \definition[个,位,名]{s.}{subordinado; subalterno}
  \synonymref{部属}{bu4shu3}
  \synonymref{部下}{bu4xia4}
  \synonymref{下级}{xia4ji2}
  \antonymref{上司}{shang4si5}
  \end{Phonetics}
\end{Entry}

\begin{Entry}{下崽}{3,12}{⼀,⼭}
  \begin{Phonetics}{下崽}{xia4 zai3}
    \definition{v.}{dar à luz à filhotes; parir}
  \end{Phonetics}
\end{Entry}

\begin{Entry}{下期}{3,12}{⼀,⽉}
  \begin{Phonetics}{下期}{xia4qi1}[][HSK 7-9]
    \definition{s.}{da próxima vez; o período ou intervalo de tempo do próximo evento ou atividade.}
  \end{Phonetics}
\end{Entry}

\begin{Entry}{下棋}{3,12}{⼀,⽊}
  \begin{Phonetics}{下棋}{xia4/qi2}[][HSK 7-9]
    \definition{v.+compl.}{jogar xadrez; ter uma partida de xadrez}
  \end{Phonetics}
\end{Entry}

\begin{Entry}{下游}{3,12}{⼀,⽔}
  \begin{Phonetics}{下游}{xia4you2}[][HSK 7-9]
    \definition{s.}{a jusante; rio abaixo; trechos inferiores; o trecho do rio próximo à sua foz | para trás; a posição inferior; referindo"-se metaforicamente a uma posição invertida}
  \antonymref{上游}{shang4you2}
  \end{Phonetics}
\end{Entry}

\begin{Entry}{下落}{3,12}{⼀,⾋}
  \begin{Phonetics}{下落}{xia4luo4}[][HSK 7-9]
    \definition{s.}{localização da pessoa ou coisa que está sendo procurada}
    \definition{v.}{cair; deixar cair}
  \synonymref{落下}{luo4xia4}
  \synonymref{下跌}{xia4die1}
  \synonymref{下降}{xia4jiang4}
  \synonymref{着落}{zhuo2luo4}
  \antonymref{爬升}{pa2sheng1}
  \antonymref{上升}{shang4sheng1}
  \antonymref{上涨}{shang4zhang3}
  \end{Phonetics}
\end{Entry}

\begin{Entry}{下跌}{3,12}{⼀,⾜}
  \begin{Phonetics}{下跌}{xia4die1}[][HSK 7-9]
    \definition{v.}{cair; despencar; descer}
  \synonymref{下降}{xia4jiang4}
  \synonymref{下落}{xia4luo4}
  \antonymref{上涨}{shang4zhang3}
  \end{Phonetics}
\end{Entry}

\begin{Entry}{下意识}{3,13,7}{⼀,⼼,⾔}
  \begin{Phonetics}{下意识}{xia4yi4shi2}[][HSK 7-9]
    \definition{adv.}{instintivamente; subconscientemente}
    \definition{s.}{instinto; subconsciente; a psicologia refere"-se às reações psicológicas inconscientes causadas por determinadas condições}
  \end{Phonetics}
\end{Entry}

\begin{Entry}{下楼}{3,13}{⼀,⽊}
  \begin{Phonetics}{下楼}{xia4lou2}[][HSK 4]
    \definition{v.}{descer as escadas}
  \end{Phonetics}
\end{Entry}

%%%%%%%%%% 与 %%%%%%%%%%
\subsection*{与}\addcontentsline{loh}{figure}{与}

\begin{Entry}{与}{3}{⼀}
  \begin{Phonetics}{与}{yu2}
    \definition{part.}{não; o quê; hein}
    \variantof{欤}
  \end{Phonetics}
  \begin{Phonetics}{与}{yu3}[][HSK 6]
    \definition*{s.}{Sobrenome: Yu}
    \definition{conj.}{e; junto com}
    \definition{prep.}{com}
    \definition{v.}{dar; oferecer; conceder | conviver com; estar em bons termos com; socializar; ser amigável | ajudar; apoiar; patrocinar | Literário: esperar}
  \synonymref{给}{gei3}
  \synonymref{跟}{gen1}
  \synonymref{和}{he2}
  \antonymref{取}{qu3}
  \end{Phonetics}
  \begin{Phonetics}{与}{yu4}
    \definition{v.}{participar de; tomar parte em}
  \synonymref{参加}{can1jia1}
  \synonymref{出席}{chu1/xi2}
  \synonymref{加入}{jia1ru4}
  \synonymref{介入}{jie4ru4}
  \antonymref{旁观}{pang2guan1}
  \end{Phonetics}
\end{Entry}

\begin{Entry}{与日俱增}{3,4,10,15}{⼀,⽇,⼈,⼟}
  \begin{Phonetics}{与日俱增}{yu3ri4-ju4zeng1}[][HSK 7-9]
    \definition{expr.}{aumentando a cada dia; crescer a cada dia que passa; crescendo ao longo do tempo}
  \end{Phonetics}
\end{Entry}

\begin{Entry}{与众不同}{3,6,4,6}{⼀,⼈,⼀,⼝}
  \begin{Phonetics}{与众不同}{yu3zhong4-bu4tong2}[][HSK 7-9]
    \definition{expr.}{fora do comum; diferente dos demais; peculiar; incomum; significa ser diferente dos outros}
  \synonymref{鹤立鸡群}{he4li4ji1qun2}
  \antonymref{司空见惯}{si1kong1-jian4guan4}
  \end{Phonetics}
\end{Entry}

\begin{Entry}{与此同时}{3,6,6,7}{⼀,⽌,⼝,⽇}
  \begin{Phonetics}{与此同时}{yu3ci3-tong2shi2}[][HSK 7-9]
    \definition{expr.}{ao mesmo tempo; entretanto; enquanto isso; simultaneamente; simultaneamente a esse evento; frequentemente usado para conectar duas frases ou orações que estão em relação paralela}[他在学习，与此同时，我在做饭。===Ele estava estudando, e ao mesmo tempo, eu estava cozinhando.]
  \synonymref{不相上下}{bu4xiang1-shang4xia4}
  \end{Phonetics}
\end{Entry}

\begin{Entry}{与否}{3,7}{⼀,⼝}
  \begin{Phonetics}{与否}{yu3fou3}[][HSK 7-9]
    \definition{adv./part.}{independentemente de ser ou não; se ou não (no final de uma frase)}
  \synonymref{是否}{shi4fou3}
  \end{Phonetics}
\end{Entry}

\begin{Entry}{与时俱进}{3,7,10,7}{⼀,⽇,⼈,⾡}
  \begin{Phonetics}{与时俱进}{yu3shi2-ju4jin4}[][HSK 7-9]
    \definition{expr.}{continuar a se desenvolver e avançar com o passar do tempo; manter"-se atualizado; estar a par dos desenvolvimentos modernos; progressivo; oportuno; acompanhar os tempos}
  \end{Phonetics}
\end{Entry}

\begin{Entry}{与其}{3,8}{⼀,⼋}
  \begin{Phonetics}{与其}{yu3qi2}[][HSK 7-9]
    \definition{conj.}{em vez de\dots; orações de ligação, indicando uma decisão tomada após comparação; 与其 é usado para o aspecto a ser abandonado, enquanto o aspecto a ser escolhido é frequentemente expresso por frases como 不如 ou 宁可}
  \seealsoref{不如}{bu4ru2}
  \seealsoref{宁可}{ning4ke3}
  \synonymref{如果}{ru2guo3}
  \synonymref{因为}{yin1wei5}
  \end{Phonetics}
\end{Entry}

\begin{Entry}{与其……不如……}{3,8,4,6}{⼀,⼋,⼀,⼥}
  \begin{Phonetics}{与其……不如……}{yu3qi2 bu4ru2}
    \definition{conj.}{em vez de\dots é melhor\dots}
  \end{Phonetics}
\end{Entry}

\begin{Entry}{与其……宁可……}{3,8,5,5}{⼀,⼋,⼧,⼝}
  \begin{Phonetics}{与其……宁可……}{yu3qi2 ning4ke3}
    \definition{conj.}{em vez de\dots eu prefereria\dots}
  \end{Phonetics}
\end{Entry}

%%%%%%%%%% 个 %%%%%%%%%%
\subsection*{个}\addcontentsline{loh}{figure}{个}

\begin{Entry}{个}{3}{⼈}
  \begin{Phonetics}{个}{ge3}
    \definition{pron.}{usado em 自个儿}
  \seealsoref{自个儿}{zi4ge3r5}
  \end{Phonetics}
  \begin{Phonetics}{个}{ge4}[][HSK 1]
    \definition{adj.}{individual}
    \definition{clas.}{usado antes de substantivos que não têm palavras de medida específicas | usado na frente do divisor; usado na frente do número aproximado | usado após verbos com objeto direto |  usado entre verbos e complementos}
    \definition{part.}{usado após pronomes demonstrativos | adicionado após certas palavras de tempo}
  \seealsoref{口}{kou3}
  \seealsoref{名}{ming2}
  \synonymref{口}{kou3}
  \synonymref{名}{ming2}
  \end{Phonetics}
\end{Entry}

\begin{Entry}{个人}{3,2}{⼈,⼈}
  \begin{Phonetics}{个人}{ge4ren2}[][HSK 3]
    \definition{pron.}{pessoal; si mesmo}
    \definition[个]{s.}{indivíduo; pessoa}
  \antonymref{集体}{ji2ti3}
  \antonymref{团体}{tuan2ti3}
  \end{Phonetics}
\end{Entry}

\begin{Entry}{个儿}{3,2}{⼈,⼉}
  \begin{Phonetics}{个儿}{ge4r5}[][HSK 5]
    \definition{s.}{tamanho; altura; estatura; tamanho do corpo ou do objeto | pessoas ou coisas consideradas isoladamente; referir"-se a uma pessoa ou coisa individualmente}
  \end{Phonetics}
\end{Entry}

\begin{Entry}{个子}{3,3}{⼈,⼦}
  \begin{Phonetics}{个子}{ge4zi5}[][HSK 2]
    \definition[个,种,些]{s.}{altura; estatura; refere"-se ao tamanho do corpo humano e também ao tamanho do corpo dos animais}
  \synonymref{身高}{shen1gao1}
  \end{Phonetics}
\end{Entry}

\begin{Entry}{个头儿}{3,5,2}{⼈,⼤,⼉}
  \begin{Phonetics}{个头儿}{ge4tou2er5}[][HSK 7-9]
    \definition{s.}{tamanho; altura}
  \end{Phonetics}
\end{Entry}

\begin{Entry}{个体}{3,7}{⼈,⼈}
  \begin{Phonetics}{个体}{ge4ti3}[][HSK 4]
    \definition[个,位]{s.}{uma única pessoa ou organismo}
  \synonymref{个别}{ge4bie2}
  \synonymref{个人}{ge4ren2}
  \antonymref{集体}{ji2ti3}
  \antonymref{群体}{qun2ti3}
  \antonymref{系统}{xi4tong3}
  \antonymref{总体}{zong3ti3}
  \end{Phonetics}
\end{Entry}

\begin{Entry}{个别}{3,7}{⼈,⼑}
  \begin{Phonetics}{个别}{ge4bie2}[][HSK 4]
    \definition{adj.}{muito poucos; excepcionais}
    \definition{adv.}{separadamente; individualmente; isoladamente}
  \synonymref{部分}{bu4fen5}
  \synonymref{个人}{ge4ren2}
  \synonymref{个体}{ge4ti3}
  \synonymref{局部}{ju2bu4}
  \synonymref{片面}{pian4mian4}
  \synonymref{一部分}{yi2bu4fen5}
  \synonymref{一面}{yi2mian4}
  \antonymref{多数}{duo1shu4}
  \antonymref{集体}{ji2ti3}
  \antonymref{普遍}{pu3bian4}
  \antonymref{普通}{pu3tong1}
  \antonymref{全都}{quan2dou1}
  \antonymref{全体}{quan2ti3}
  \antonymref{特殊}{te4shu1}
  \antonymref{一般}{yi4ban1}
  \end{Phonetics}
\end{Entry}

\begin{Entry}{个性}{3,8}{⼈,⼼}
  \begin{Phonetics}{个性}{ge4xing4}[][HSK 3]
    \definition[种,点儿]{s.}{individualidade; personalidade; caráter individual; as características relativamente fixas de uma pessoa, formadas sob determinadas condições sociais e influências educacionais | propriedade específica; caráter específico; a propriedade ou característica especial que distingue uma coisa de outras coisas}
  \synonymref{本性}{ben3xing4}
  \synonymref{脾气}{pi2qi5}
  \synonymref{特性}{te4xing4}
  \synonymref{天性}{tian1xing4}
  \synonymref{性格}{xing4ge2}
  \synonymref{性情}{xing4qing2}
  \antonymref{共性}{gong4xing4}
  \end{Phonetics}
\end{Entry}

\begin{Entry}{个案}{3,10}{⼈,⽊}
  \begin{Phonetics}{个案}{ge4'an4}[][HSK 7-9]
    \definition[个,些]{s.}{caso individual (ou especial); caso; caso a caso}
  \end{Phonetics}
\end{Entry}

%%%%%%%%%% 丫 %%%%%%%%%%
\subsection*{丫}\addcontentsline{loh}{figure}{丫}

\begin{Entry}{丫}{3}{⼁}
  \begin{Phonetics}{丫}{ya1}
    \definition{s.}{bifurcação; forquilha | Dialeto: menina}
  \end{Phonetics}
\end{Entry}

\begin{Entry}{丫头}{3,5}{⼁,⼤}
  \begin{Phonetics}{丫头}{ya1tou5}[][HSK 7-9]
    \definition[个]{s.}{garota; menina; o penteado com coque usado por crianças antigamente, com os cabelos espetados em ambos os lados da cabeça; usado metaforicamente para se referir a meninas | empregada doméstica; criada}
  \synonymref{姑娘}{gu1niang5}
  \synonymref{女儿}{nv3'er2}
  \synonymref{女孩}{nv3hai2}
  \end{Phonetics}
\end{Entry}

%%%%%%%%%% 丸 %%%%%%%%%%
\subsection*{丸}\addcontentsline{loh}{figure}{丸}

\begin{Entry}{丸}{3}{⼂}
  \begin{Phonetics}{丸}{wan2}[][HSK 7-9]
    \definition{clas.}{utilizado para medicamentos em comprimido}
    \definition[个]{s.}{bola; grânulo | comprimido; bolo (como em bolo alimentar); pílula}
  \seealsoref{丸儿}{wan2r5}
  \end{Phonetics}
\end{Entry}

\begin{Entry}{丸儿}{3,2}{⼂,⼉}
  \begin{Phonetics}{丸儿}{wan2r5}
    \definition{s.}{bola; grânulo}
  \end{Phonetics}
\end{Entry}

%%%%%%%%%% 久 %%%%%%%%%%
\subsection*{久}\addcontentsline{loh}{figure}{久}

\begin{Entry}{久}{3}{⼃}
  \begin{Phonetics}{久}{jiu3}[][HSK 3]
    \definition{adj.}{por muito tempo; longo período de tempo | duração de tempo especificada}
  \antonymref{暂}{zan4}
  \end{Phonetics}
\end{Entry}

\begin{Entry}{久仰}{3,6}{⼃,⼈}
  \begin{Phonetics}{久仰}{jiu3yang3}[][HSK 7-9]
    \definition{expr.}{``É um prazer conhecê"-lo(a).''; ``Há muito tempo que desejo conhecê"-lo(a).''; ``Há muito tempo que aguardo ansiosamente o nosso encontro.''; ``Já ouvi falar muito de você.''}
  \end{Phonetics}
\end{Entry}

\begin{Entry}{久违}{3,7}{⼃,⾡}
  \begin{Phonetics}{久违}{jiu3wei2}[][HSK 7-9]
    \definition{v.}{não ver há muito tempo; fazer muito tempo desde o último encontro; apenas um comentário educado, muito tempo sem ver}
  \synonymref{重逢}{chong2feng2}
  \antonymref{经常}{jing1chang2}
  \end{Phonetics}
\end{Entry}

%%%%%%%%%% 义 %%%%%%%%%%
\subsection*{义}\addcontentsline{loh}{figure}{义}

\begin{Entry}{义}{3}{⼂}
  \begin{Phonetics}{义}{yi4}
    \definition*{s.}{Sobrenome: Yi}
    \definition{adj.}{justo; equitativo | adotado; adotivo | juramentado | artificial; falso}
    \definition[个]{s.}{justiça; retidão | laços humanos; relacionamento | significado; importância}
  \end{Phonetics}
\end{Entry}

\begin{Entry}{义工}{3,3}{⼂,⼯}
  \begin{Phonetics}{义工}{yi4gong1}[][HSK 7-9]
    \definition{s.}{voluntário; isso se refere a indivíduos e grupos que, sem levar em conta a recompensa material, contribuem com seu tempo, energia e conhecimento técnico pessoal para melhorar a sociedade com base na moralidade, crenças, consciência, compaixão e responsabilidade}
  \synonymref{志愿者}{zhi4yuan4zhe3}
  \end{Phonetics}
\end{Entry}

\begin{Entry}{义务}{3,5}{⼂,⼒}
  \begin{Phonetics}{义务}{yi4wu4}[][HSK 4]
    \definition{adj.}{voluntário; fornecer serviços ou ajuda a outros gratuitamente}
    \definition[项]{s.}{dever; obrigação; responsabilidades perante a lei | obrigação moral; responsabilidade moral}
  \synonymref{负担}{fu4dan1}
  \synonymref{任务}{ren4wu5}
  \synonymref{无偿}{wu2chang2}
  \synonymref{仔肩}{zi1jian1}
  \antonymref{权利}{quan2li4}
  \antonymref{权力}{quan2li4}
  \antonymref{权益}{quan2yi4}
  \antonymref{责任}{ze2ren4}
  \end{Phonetics}
\end{Entry}

%%%%%%%%%% 之 %%%%%%%%%%
\subsection*{之}\addcontentsline{loh}{figure}{之}

\begin{Entry}{之}{3}{⼂}
  \begin{Phonetics}{之}{zhi1}[][HSK 7-9]
    \definition*{s.}{Sobrenome: Zhi}
    \definition{part.}{entre um atributivo e a palavra que ele modifica; equivalente a 的 | usado entre o sujeito e o predicado, em estruturas sujeito"-predicado, de modo a torná"-lo nominalizado}
    \definition{pron.}{substituto de uma pessoa ou coisa, limitado a ser usado como um objeto; substituir a pessoa ou coisa mencionada anteriormente | isto; isso; não substitui uma pessoa ou coisa específica, mas serve apenas para complementar sílabas}
    \definition{v.}{ir; deixar}
  \seealsoref{的}{de5}
  \end{Phonetics}
\end{Entry}

\begin{Entry}{之一}{3,1}{⼂,⼀}
  \begin{Phonetics}{之一}{zhi1yi1}[][HSK 4]
    \definition[分]{s.}{um de (algo); pertence a um ou a todo um grupo de coisas com as mesmas características}
  \synonymref{之中}{zhi1zhong1}
  \end{Phonetics}
\end{Entry}

\begin{Entry}{之下}{3,3}{⼂,⼀}
  \begin{Phonetics}{之下}{zhi1xia4}[][HSK 5]
    \definition{s.}{usado para indicar algo abaixo de um determinado intervalo, posição, grau, etc.; indica um aspecto inferior em termos de alcance, posição, status, nível, Chengdu, etc. | usado para indicar as condições sob as quais algo acontece | usado para indicar o humor, estado em que alguém faz algo; expressa um determinado comportamento em um determinado estado de espírito ou situação}
  \synonymref{以下}{yi3xia4}
  \end{Phonetics}
\end{Entry}

\begin{Entry}{之中}{3,4}{⼂,⼁}
  \begin{Phonetics}{之中}{zhi1zhong1}[][HSK 5]
    \definition{prep.}{em; no meio de; entre}
  \synonymref{之一}{zhi1yi1}
  \end{Phonetics}
\end{Entry}

\begin{Entry}{之内}{3,4}{⼂,⼌}
  \begin{Phonetics}{之内}{zhi1nei4}[][HSK 5]
    \definition{adv.}{em; dentro de; indica dentro de um determinado intervalo, limite ou período de tempo, etc.}
  \synonymref{以内}{yi3nei4}
  \end{Phonetics}
\end{Entry}

\begin{Entry}{之外}{3,5}{⼂,⼣}
  \begin{Phonetics}{之外}{zhi1wai4}[][HSK 5]
    \definition{adv.}{lado de fora; exceto; além de; além disso; refere"-se a algo que excede um determinado limite}
  \synonymref{除外}{chu2wai4}
  \synonymref{以外}{yi3wai4}
  \end{Phonetics}
\end{Entry}

\begin{Entry}{之后}{3,6}{⼂,⼝}
  \begin{Phonetics}{之后}{zhi1hou4}[][HSK 4]
    \definition{s.}{mais tarde; posteriormente; depois; desde então; para indicar que é depois de um determinado tempo ou de uma determinada coisa, 以后 é usado com frequência na linguagem falada; às vezes, também pode indicar que é depois de um determinado lugar ou local,  后面 é usado com frequência na linguagem falada}
  \seealsoref{后面}{hou4mian5}
  \seealsoref{以后}{yi3hou4}
  \synonymref{后来}{hou4lai2}
  \synonymref{继而}{ji4'er2}
  \synonymref{接着}{jie1zhe5}
  \synonymref{其后}{qi2hou4}
  \synonymref{事后}{shi4hou4}
  \synonymref{以后}{yi3hou4}
  \antonymref{之前}{zhi1qian2}
  \end{Phonetics}
\end{Entry}

\begin{Entry}{之间}{3,7}{⼂,⾨}
  \begin{Phonetics}{之间}{zhi1jian1}[][HSK 4]
    \definition{s.}{(depois de um substantivo) entre; dentro de duas delimitações de tempo, local ou quantitativas | colocado após certos verbos ou advérbios de duas sílabas para indicar um curto período de tempo}
  \end{Phonetics}
\end{Entry}

\begin{Entry}{之所以}{3,8,4}{⼂,⼾,⼈}
  \begin{Phonetics}{之所以}{zhi1suo3yi3}[][HSK 7-9]
    \definition{conj.}{porque; a razão pela qual; é utilizado em frases causais complexas para introduzir o resultado na primeira oração e explicar a causa na segunda oração}[他之所以成功是因为努力。===Ele obteve sucesso graças ao seu trabalho árduo.]
  \end{Phonetics}
\end{Entry}

\begin{Entry}{之前}{3,9}{⼂,⼑}
  \begin{Phonetics}{之前}{zhi1qian2}[][HSK 4]
    \definition{adv.}{(referindo"-se ao tempo) antes, antes de, atrás | (referindo"-se ao local físico) na frente de | (usado independentemente) no passado, antes disso}
  \synonymref{曾经}{ceng2jing1}
  \synonymref{此前}{ci3qian2}
  \synonymref{以前}{yi3qian2}
  \antonymref{正在}{zheng4zai4}
  \antonymref{之后}{zhi1hou4}
  \end{Phonetics}
\end{Entry}

\begin{Entry}{之类}{3,9}{⼂,⽶}
  \begin{Phonetics}{之类}{zhi1lei4}[][HSK 6]
    \definition{s.}{usado para dar exemplos (coisas do tipo, desse tipo, assim); uma categoria de pessoas ou coisas que compartilham as mesmas características das pessoas ou coisas mencionadas anteriormente}[我喜欢香蕉、苹果之类的水果。===Eu gosto de frutas como bananas e maçãs.]
  \end{Phonetics}
\end{Entry}

%%%%%%%%%% 乞 %%%%%%%%%%
\subsection*{乞}\addcontentsline{loh}{figure}{乞}

\begin{Entry}{乞}{3}{⼄}
  \begin{Phonetics}{乞}{qi3}
    \definition*{s.}{Sobrenome: Qi}
    \definition{v.}{implorar (por esmolas, etc.); suplicar}
  \end{Phonetics}
\end{Entry}

\begin{Entry}{乞丐}{3,4}{⼄,⼀}
  \begin{Phonetics}{乞丐}{qi3gai4}[][HSK 7-9]
    \definition[个,位,群]{s.}{mendigo; pessoas que não têm meios de subsistência e dependem exclusivamente da mendicância para conseguir comida e dinheiro para sobreviver}
  \synonymref{乞讨}{qi3tao3}
  \antonymref{富人}{fu4ren2}
  \end{Phonetics}
\end{Entry}

\begin{Entry}{乞讨}{3,5}{⼄,⾔}
  \begin{Phonetics}{乞讨}{qi3tao3}[][HSK 7-9]
    \definition{v.}{implorar; pedir dinheiro, pedir comida, etc.}
  \synonymref{乞丐}{qi3gai4}
  \antonymref{恩赐}{en1ci4}
  \end{Phonetics}
\end{Entry}

\begin{Entry}{乞求}{3,7}{⼄,⽔}
  \begin{Phonetics}{乞求}{qi3qiu2}[][HSK 7-9]
    \definition{v.}{implorar; suplicar; mendigar | cair de joelhos}
  \synonymref{哀求}{ai1qiu2}
  \synonymref{恳求}{ken3qiu2}
  \synonymref{请求}{qing3qiu2}
  \antonymref{恩赐}{en1ci4}
  \end{Phonetics}
\end{Entry}

%%%%%%%%%% 也 %%%%%%%%%%
\subsection*{也}\addcontentsline{loh}{figure}{也}

\begin{Entry}{也}{3}{⼄}
  \begin{Phonetics}{也}{ye3}[][HSK 1]
    \definition*{s.}{Sobrenome: Ye}
    \definition{adv.}{também; igualmente; assim como; da mesma forma; usado em frases simples, implica que é igual a outra coisa | assim como (expressar ênfase) | (expressar que as consequências são as mesmas) | também (expressar ufemismo; expressar um tom diplomático) | usado em frases compostas paralelas, indica que duas ou mais coisas têm algo em comum (pode ser usado em todas as frases ou apenas na última frase)}
    \definition{part.}{usado no meio de uma frase, destacando um elemento da frase sobre o qual deve ser feita uma afirmação | usado no final de uma frase, indicando uma explicação ou um julgamento; usado no final da frase, indica tom afirmativo e também pode reforçar o tom interrogativo, exclamativo ou imperativo}
  \seealsoref{都}{dou1}
  \synonymref{都}{dou1}
  \synonymref{亦}{yi4}
  \end{Phonetics}
\end{Entry}

\begin{Entry}{也好}{3,6}{⼄,⼥}
  \begin{Phonetics}{也好}{ye3hao3}[][HSK 5]
    \definition{part.}{pode não ser uma má ideia; também pode ser | (reduplicado) se\dots ou\dots; não importa se | pode não ser uma má ideia | se\dots ou\dots; usado em conjunto, significa que não está condicionado a uma determinada situação}
  \end{Phonetics}
\end{Entry}

\begin{Entry}{也有今天}{3,6,4,4}{⼄,⽉,⼈,⼤}
  \begin{Phonetics}{也有今天}{ye3 you3 jin1tian1}
    \definition{expr.}{Gíria: receber o que se merece | receber a sua parte (de coisas boas ou ruins) | servir alguém corretamente}
    \definition{expr.}{``Todo cachorro tem seu dia.''}
  \end{Phonetics}
\end{Entry}

\begin{Entry}{也许}{3,6}{⼄,⾔}
  \begin{Phonetics}{也许}{ye3xu3}[][HSK 2]
    \definition{adv.}{talvez; provavelmente; estou com medo; para expressar incerteza; para expressar uma alta probabilidade}
  \synonymref{或许}{huo4xu3}
  \synonymref{或者}{huo4zhe3}
  \end{Phonetics}
\end{Entry}

\begin{Entry}{也就是}{3,12,9}{⼄,⼪,⽇}
  \begin{Phonetics}{也就是}{ye3 jiu4shi4}
    \definition{adv.}{i.e., isso é | ou seja}
  \end{Phonetics}
\end{Entry}

\begin{Entry}{也就是说}{3,12,9,9}{⼄,⼪,⽇,⾔}
  \begin{Phonetics}{也就是说}{ye3jiu4shi4shuo1}[][HSK 7-9]
    \definition{adv.}{em outras palavras | então | isto é; ou seja | por isso; assim}
  \end{Phonetics}
\end{Entry}

%%%%%%%%%% 习 %%%%%%%%%%
\subsection*{习}\addcontentsline{loh}{figure}{习}

\begin{Entry}{习}{3}{⼄}
  \begin{Phonetics}{习}{xi2}
    \definition*{s.}{Sobrenome: Xi}
    \definition{s.}{hábito; costume; prática usual; um comportamento que se desenvolve inconscientemente por meio de ações repetidas ao longo de um longo período de tempo}
    \definition{v.}{revisar; praticar; exercitar | acostumado a; familiarizado com; familiarizado com algo por meio de contato frequente | estudar; aprender (pássaro)}
  \end{Phonetics}
\end{Entry}

\begin{Entry}{习尚}{3,8}{⼄,⼩}
  \begin{Phonetics}{习尚}{xi2shang4}
    \definition{s.}{prática comum; costume; prática usual}
  \synonymref{风气}{feng1qi4}
  \synonymref{风尚}{feng1shang4}
  \end{Phonetics}
\end{Entry}

\begin{Entry}{习俗}{3,9}{⼄,⼈}
  \begin{Phonetics}{习俗}{xi2su2}[][HSK 7-9]
    \definition[种,个]{s.}{costume; convenção; hábitos e costumes}
  \synonymref{风气}{feng1qi4}
  \synonymref{风俗}{feng1su2}
  \synonymref{民俗}{min2su2}
  \synonymref{习惯}{xi2guan4}
  \end{Phonetics}
\end{Entry}

\begin{Entry}{习惯}{3,11}{⼄,⼼}
  \begin{Phonetics}{习惯}{xi2guan4}[][HSK 2]
    \definition[个,种]{s.}{hábito; costume; prática usual; comportamentos, tendências ou tendências sociais que se desenvolvem gradualmente ao longo de um longo período de tempo e são difíceis de mudar}
    \definition{v.}{estar acostumado a; ter o hábito de}
  \end{Phonetics}
\end{Entry}

%%%%%%%%%% 乡 %%%%%%%%%%
\subsection*{乡}\addcontentsline{loh}{figure}{乡}

\begin{Entry}{乡}{3}{⼄}
  \begin{Phonetics}{乡}{xiang1}[][HSK 5]
    \definition[个,座,片]{s.}{país; campo; vilarejo; área rural | local de origem; vila ou cidade natal | município (uma unidade administrativa rural subordinada ao condado) | vila natal; cidade natal | terra ou local famoso por produzir algo}
  \antonymref{城}{cheng2}
  \end{Phonetics}
\end{Entry}

\begin{Entry}{乡下}{3,3}{⼄,⼀}
  \begin{Phonetics}{乡下}{xiang1xia4}[][HSK 7-9]
    \definition[个]{s.}{país; campo; interior; área rural}
  \synonymref{村庄}{cun1zhuang1}
  \synonymref{农村}{nong2cun1}
  \synonymref{乡村}{xiang1cun1}
  \antonymref{城市}{cheng2shi4}
  \antonymref{都市}{du1shi4}
  \end{Phonetics}
\end{Entry}

\begin{Entry}{乡巴佬}{3,4,8}{⼄,⼰,⼈}
  \begin{Phonetics}{乡巴佬}{xiang1ba1lao3}
    \definition{s.}{Pejorativo: aldeão; caipira}
  \end{Phonetics}
\end{Entry}

\begin{Entry}{乡村}{3,7}{⼄,⽊}
  \begin{Phonetics}{乡村}{xiang1cun1}[][HSK 5]
    \definition{adj.}{rural | rústico}
    \definition[个]{s.}{vila; campo; área rural; principalmente envolvido na agricultura; áreas com distribuição populacional mais dispersa em relação às cidades}
  \synonymref{农村}{nong2cun1}
  \antonymref{城市}{cheng2shi4}
  \end{Phonetics}
\end{Entry}

\begin{Entry}{乡亲}{3,9}{⼄,⼇}
  \begin{Phonetics}{乡亲}{xiang1qin1}[][HSK 7-9]
    \definition{s.}{pessoa da mesma aldeia ou cidade; conterrâneo ou morador da mesma cidade | população local; aldeões; gente; povo local | termo de tratamento direto para pessoas locais ou moradores de vilarejos}
  \synonymref{故乡}{gu4xiang1}
  \synonymref{家园}{jia1yuan2}
  \end{Phonetics}
\end{Entry}

%%%%%%%%%% 于 %%%%%%%%%%
\subsection*{于}\addcontentsline{loh}{figure}{于}

\begin{Entry}{于}{3}{⼆}
  \begin{Phonetics}{于}{yu2}[][HSK 6]
    \definition*{s.}{Sobrenome: Yu}
    \definition{prep.}{indica hora, lugar, alcance, etc. | indica a direção da ação | usado depois de um verbo para indicar dar, entregar, etc. | apresentar a relação do objeto ou entidade introduzida | indica o ponto de início ou de partida | indica comparação | indica passividade}
  \end{Phonetics}
\end{Entry}

\begin{Entry}{于是}{3,9}{⼆,⽇}
  \begin{Phonetics}{于是}{yu2shi4}[][HSK 4]
    \definition{conj.}{então; portanto; consequentemente; como resultado; indica que o último segue o primeiro e que o último é frequentemente causado pelo primeiro}
  \synonymref{所以}{suo3yi3}
  \synonymref{因而}{yin1'er2}
  \synonymref{因此}{yin1ci3}
  \antonymref{但是}{dan4shi4}
  \antonymref{然而}{ran2'er2}
  \end{Phonetics}
\end{Entry}

%%%%%%%%%% 亏 %%%%%%%%%%
\subsection*{亏}\addcontentsline{loh}{figure}{亏}

\begin{Entry}{亏}{3}{⼆}
  \begin{Phonetics}{亏}{kui1}[][HSK 5]
    \definition{adv.}{felizmente; por sorte; graças a | contrariamente, expressando sarcasmo}
    \definition{s.}{prejuízo; perda; déficit | perda; dano; ferida}
    \definition{v.}{perder dinheiro, etc.; ter um déficit; ter prejuízo | ter falta de; ser deficiente; carecer de | tratar injustamente; causar prejuízo; trair a confiança}
  \synonymref{缺}{que1}
  \synonymref{损}{sun3}
  \antonymref{盈}{ying2}
  \end{Phonetics}
\end{Entry}

\begin{Entry}{亏本}{3,5}{⼆,⽊}
  \begin{Phonetics}{亏本}{kui1/ben3}[][HSK 7-9]
    \definition{v.+compl.}{estar no vermelho; perder o capital; perder dinheiro (nos negócios)}
  \synonymref{亏损}{kui1sun3}
  \antonymref{盈利}{ying2/li4}
  \antonymref{赚钱}{zhuan4qian2}
  \end{Phonetics}
\end{Entry}

\begin{Entry}{亏损}{3,10}{⼆,⼿}
  \begin{Phonetics}{亏损}{kui1sun3}[][HSK 7-9]
    \definition{v.}{perder; esgotar; ter déficit; as despesas excedem as receitas | desidratar; enfraquecer; o corpo está enfraquecido devido a danos ou falta de nutrição}
  \synonymref{不足}{bu4zu2}
  \synonymref{吃亏}{chi1/kui1}
  \synonymref{耗费}{hao4fei4}
  \synonymref{亏本}{kui1/ben3}
  \synonymref{丧失}{sang4shi1}
  \synonymref{损失}{sun3shi1}
  \synonymref{牺牲}{xi1sheng1}
  \antonymref{盈利}{ying2/li4}
  \end{Phonetics}
\end{Entry}

%%%%%%%%%% 亡 %%%%%%%%%%
\subsection*{亡}\addcontentsline{loh}{figure}{亡}

\begin{Entry}{亡}{3}{⼇}
  \begin{Phonetics}{亡}{wang2}
    \definition{adj.}{falecido}
    \definition{v.}{fugir; escapar | perder; ir embora; jogar fora | morrer; perecer; falecer | conquistar; subjugar | ser destruído; morrer}
  \synonymref{灭}{mie4}
  \synonymref{死}{si3}
  \synonymref{卒}{zu2}
  \antonymref{存}{cun2}
  \antonymref{兴}{xing1}
  \end{Phonetics}
\end{Entry}

\begin{Entry}{亡羊补牢}{3,6,7,7}{⼇,⽺,⾐,⼧}
  \begin{Phonetics}{亡羊补牢}{wang2yang2-bu3lao2}[][HSK 7-9]
    \definition{expr.}{consertar a situação antes que seja tarde demais; agir tardiamente após a ocorrência de um acidente; reparar o curral depois que uma ovelha se perde é uma metáfora para encontrar maneiras de fazer as pazes depois de sofrer uma perda, de modo a evitar sofrer perdas novamente no futuro}
  \end{Phonetics}
\end{Entry}

%%%%%%%%%% 亿 %%%%%%%%%%
\subsection*{亿}\addcontentsline{loh}{figure}{亿}

\begin{Entry}{亿}{3}{⼈}
  \begin{Phonetics}{亿}{yi4}[][HSK 2]
    \definition*{s.}{Sobrenome: Yi}
    \definition{num.}{cem milhões; 100.000.000; 1.0000.0000}
  \end{Phonetics}
\end{Entry}

%%%%%%%%%% 凡 %%%%%%%%%%
\subsection*{凡}\addcontentsline{loh}{figure}{凡}

\begin{Entry}{凡}{3}{⼏}
  \begin{Phonetics}{凡}{fan2}[][HSK 7-9]
    \definition*{s.}{Sobrenome: Fan}
    \definition{adj.}{comum; ordinário}
    \definition{adv.}{Literário: qualquer; todos; todo | Literário: em tudo; completamente}
    \definition{s.}{este mundo mortal; a terra | o mundo secular; refere"-se ao mundo humano | uma nota da escala em Gongchepu (工尺谱), correspondente a 4 na notação musical numerada | Literário: ideia geral; esboço}
  \seealsoref{工尺谱}{gong1che3 pu3}
  \end{Phonetics}
\end{Entry}

\begin{Entry}{凡是}{3,9}{⼏,⽇}
  \begin{Phonetics}{凡是}{fan2shi4}[][HSK 6]
    \definition{adv.}{todos; qualquer; cada; resumir tudo dentro de um determinado âmbito}
  \end{Phonetics}
\end{Entry}

%%%%%%%%%% 勺 %%%%%%%%%%
\subsection*{勺}\addcontentsline{loh}{figure}{勺}

\begin{Entry}{勺}{3}{⼓}
  \begin{Phonetics}{勺}{shao2}[][HSK 6]
    \definition{clas.}{shao; uma unidade tradicional de volume, igual a 0,01 市升, e equivalente a 1 centilitro ou 0,018 \emph{pint}}
    \definition{s.}{colher; concha}
  \seealsoref{市升}{shi4 sheng1}
  \end{Phonetics}
\end{Entry}

%%%%%%%%%% 千 %%%%%%%%%%
\subsection*{千}\addcontentsline{loh}{figure}{千}

\begin{Entry}{千}{3}{⼗}
  \begin{Phonetics}{千}{qian1}[][HSK 2]
    \definition*{s.}{Sobrenome: Qian}
    \definition{num.}{mil; 1.000; 1000 | a grande quantidade de; um grande número de}
  \end{Phonetics}
\end{Entry}

\begin{Entry}{千万}{3,3}{⼗,⼀}
  \begin{Phonetics}{千万}{qian1wan4}[][HSK 3]
    \definition{adv.}{(usado para indicar desejos fortes) por todos os meios; sob quaisquer circunstâncias; expressa uma exortação sincera, equivalente a 务必}
    \definition{num.}{dez milhões; 10.000.000; 1000.0000; milhões e milhões; um número aproximado, indicando um grande número}
  \seealsoref{务必}{wu4bi4}
  \end{Phonetics}
\end{Entry}

\begin{Entry}{千千万万}{3,3,3,3}{⼗,⼗,⼀,⼀}
  \begin{Phonetics}{千千万万}{qian1qian1wan4wan4}
    \definition{expr.}{inumerável; números incontáveis; milhares e milhares}
  \end{Phonetics}
\end{Entry}

\begin{Entry}{千方百计}{3,4,6,4}{⼗,⽅,⽩,⾔}
  \begin{Phonetics}{千方百计}{qian1fang1-bai3ji4}[][HSK 7-9]
    \definition{expr.}{por todos os meios; fazer tudo o que for possível; descreve alguém que esgotou todos os meios ou métodos}
  \end{Phonetics}
\end{Entry}

\begin{Entry}{千古}{3,5}{⼗,⼝}
  \begin{Phonetics}{千古}{qian1gu3}
    \definition*{s.}{eternidade (usada em um dístico elegíaco, coroa de flores, etc., dedicada aos mortos); inscrições para dísticos e coroas de flores fúnebres}
    \definition{adj.}{eterno}
    \definition{adv.}{por toda a eternidade; para todo o sempre; através dos tempos}
  \synonymref{千年}{qian1nian2}
  \end{Phonetics}
\end{Entry}

\begin{Entry}{千军万马}{3,6,3,3}{⼗,⼍,⼀,⾺}
  \begin{Phonetics}{千军万马}{qian1jun1-wan4ma3}[][HSK 7-9]
    \definition{expr.}{``Milhares de soldados.''; milhares e milhares de homens e cavalos; um exército poderoso; uma força imensa; todos os cavalos do rei e todos os homens do rei; exército magnífico com milhares de homens e cavalos; demonstração impressionante de força humana}
  \end{Phonetics}
\end{Entry}

\begin{Entry}{千年}{3,6}{⼗,⼲}
  \begin{Phonetics}{千年}{qian1nian2}
    \definition{s.}{milênio; uma metáfora para muito tempo}
  \end{Phonetics}
\end{Entry}

\begin{Entry}{千克}{3,7}{⼗,⼗}
  \begin{Phonetics}{千克}{qian1ke4}[][HSK 2]
    \definition{clas.}{kg; quilo; quilograma; 1 quilograma equivale a 1.000 gramas, ou 2 jin (斤)}
  \seealsoref{斤}{jin1}
  \end{Phonetics}
\end{Entry}

\begin{Entry}{千变万化}{3,8,3,4}{⼗,⼜,⼀,⼔}
  \begin{Phonetics}{千变万化}{qian1bian4-wan4hua4}[][HSK 7-9]
    \definition{expr.}{``Sempre em mudança.''; as miríades de mudanças; mudança caleidoscópica; mudanças intermináveis; em constante transformação; ser infinito em variedade; mudanças infinitas; em constante mudança}
  \end{Phonetics}
\end{Entry}

\begin{Entry}{千钧一发}{3,9,1,5}{⼗,⾦,⼀,⼜}
  \begin{Phonetics}{千钧一发}{qian1jun1-yi1fa4}[][HSK 7-9]
    \definition{expr.}{``Por pouco não deu certo.''; cem pesos pendurados por um fio; em perigo iminente; uma questão de vida ou morte}
  \end{Phonetics}
\end{Entry}

\begin{Entry}{千家万户}{3,10,3,4}{⼗,⼧,⼀,⼾}
  \begin{Phonetics}{千家万户}{qian1jia1-wan4hu4}[][HSK 7-9]
    \definition{expr.}{``Milhares de famílias.''; inúmeras famílias; todas as famílias}
  \end{Phonetics}
\end{Entry}

%%%%%%%%%% 卫 %%%%%%%%%%
\subsection*{卫}\addcontentsline{loh}{figure}{卫}

\begin{Entry}{卫}{3}{⼙}
  \begin{Phonetics}{卫}{wei4}
    \definition*{s.}{Wei, um estado da Dinastia Zhou | Sobrenome: Wei}
    \definition{s.}{uma palavra usada no nome do lugar | outro nome para um burro}
    \definition{v.}{defender; guardar; proteger}
  \end{Phonetics}
\end{Entry}

\begin{Entry}{卫生}{3,5}{⼙,⽣}
  \begin{Phonetics}{卫生}{wei4sheng1}[][HSK 3]
    \definition{adj.}{bom para a saúde; higiênico; limpo; capaz de prevenir doenças e benéfico para a saúde}
    \definition{s.}{higiene; saneamento; situação limpa}
  \end{Phonetics}
\end{Entry}

\begin{Entry}{卫生巾}{3,5,3}{⼙,⽣,⼱}
  \begin{Phonetics}{卫生巾}{wei4sheng1jin1}
    \definition{s.}{absorvente higiênico; absorvente feminino; um produto de higiene íntima usado por mulheres durante a menstruação}
  \end{Phonetics}
\end{Entry}

\begin{Entry}{卫生厅}{3,5,4}{⼙,⽣,⼚}
  \begin{Phonetics}{卫生厅}{wei4sheng1ting1}
    \definition*{s.}{Departamento de Saúde (da Província)}
  \end{Phonetics}
\end{Entry}

\begin{Entry}{卫生防疫}{3,5,6,9}{⼙,⽣,⾩,⽧}
  \begin{Phonetics}{卫生防疫}{wei4sheng1fang2yi4}
    \definition{s.}{prevenção contra a epidemia}
  \end{Phonetics}
\end{Entry}

\begin{Entry}{卫生局}{3,5,7}{⼙,⽣,⼫}
  \begin{Phonetics}{卫生局}{wei4sheng1ju2}
    \definition*{s.}{Departamento de Saúde | Escritório de Saúde}
  \end{Phonetics}
\end{Entry}

\begin{Entry}{卫生纸}{3,5,7}{⼙,⽣,⽷}
  \begin{Phonetics}{卫生纸}{wei4sheng1zhi3}
    \definition{s.}{papel higiênico | lenços desinfetantes específicos para mulheres durante o período menstrual}
  \end{Phonetics}
\end{Entry}

\begin{Entry}{卫生间}{3,5,7}{⼙,⽣,⾨}
  \begin{Phonetics}{卫生间}{wei4sheng1jian1}[][HSK 3]
    \definition[间,个]{s.}{banheiro; sanitário; \emph{toilette}; quartos com instalações sanitárias em hotéis ou residências}
  \end{Phonetics}
\end{Entry}

\begin{Entry}{卫生套}{3,5,10}{⼙,⽣,⼤}
  \begin{Phonetics}{卫生套}{wei4sheng1 tao4}
    \definition[只]{s.}{preservativo; camisinha}
  \end{Phonetics}
\end{Entry}

\begin{Entry}{卫生部}{3,5,10}{⼙,⽣,⾢}
  \begin{Phonetics}{卫生部}{wei4sheng1bu4}
    \definition*{s.}{Ministério da Saúde | Ministério da Saúde Pública}
  \end{Phonetics}
\end{Entry}

\begin{Entry}{卫生球}{3,5,11}{⼙,⽣,⽟}
  \begin{Phonetics}{卫生球}{wei4sheng1qiu2}
    \definition{s.}{naftalina; bolas de naftalina}
  \end{Phonetics}
\end{Entry}

\begin{Entry}{卫生棉}{3,5,12}{⼙,⽣,⽊}
  \begin{Phonetics}{卫生棉}{wei4sheng1mian2}
    \definition{s.}{absorvente | algodão absorvente esterilizado (usado para curativos ou limpeza de feridas) | absorvente tampão; \emph{tampon}}
  \end{Phonetics}
\end{Entry}

\begin{Entry}{卫生署}{3,5,13}{⼙,⽣,⽹}
  \begin{Phonetics}{卫生署}{wei4sheng1shu3}
    \definition*{s.}{Agência de Saúde (ou Escritório, ou Departamento)}
  \end{Phonetics}
\end{Entry}

\begin{Entry}{卫视}{3,8}{⼙,⾒}
  \begin{Phonetics}{卫视}{wei4shi4}[][HSK 7-9]
    \definition[大,家]{s.}{televisão por satélite; abreviação de 卫星电视}
  \seealsoref{卫星电视}{wei4xing1 dian4shi4}
  \end{Phonetics}
\end{Entry}

\begin{Entry}{卫星}{3,9}{⼙,⽇}
  \begin{Phonetics}{卫星}{wei4xing1}[][HSK 5]
    \definition[个,颗]{s.}{satélite; lua; corpos celestes orbitando planetas | satélite artificial | algo que gira em torno de um centro}
  \end{Phonetics}
\end{Entry}

\begin{Entry}{卫星电视}{3,9,5,8}{⼙,⽇,⽥,⾒}
  \begin{Phonetics}{卫星电视}{wei4xing1 dian4shi4}
    \definition{s.}{TV por satélite; televisão por satélite}
  \end{Phonetics}
\end{Entry}

%%%%%%%%%% 叉 %%%%%%%%%%
\subsection*{叉}\addcontentsline{loh}{figure}{叉}

\begin{Entry}{叉}{3}{⼜}
  \begin{Phonetics}{叉}{cha1}[][HSK 5]
    \definition{s.}{garfo; forquilha | símbolo de cruz, ``×''}
    \definition{v.}{trabalhar com um garfo; garfar; pegar coisas com um garfo}
  \end{Phonetics}
  \begin{Phonetics}{叉}{cha2}
    \definition{v.}{bloquear; emperrar; congestionar}
  \end{Phonetics}
  \begin{Phonetics}{叉}{cha3}
    \definition{v.}{separar de modo a formar uma bifurcação; bifurcar}
  \end{Phonetics}
\end{Entry}

\begin{Entry}{叉子}{3,3}{⼜,⼦}
  \begin{Phonetics}{叉子}{cha1zi5}[][HSK 5]
    \definition[把,个]{s.}{garfo; ferramenta com mais de duas pontas em uma extremidade | tridente; forquilha; ferramentas de agricultura antigas}
  \end{Phonetics}
\end{Entry}

%%%%%%%%%% 及 %%%%%%%%%%
\subsection*{及}\addcontentsline{loh}{figure}{及}

\begin{Entry}{及}{3}{⼃}
  \begin{Phonetics}{及}{ji2}[][HSK 7-9]
    \definition*{s.}{Sobrenome: Ji}
    \definition{conj.}{e; bem como; conectando substantivos paralelos ou frases nominais}
    \definition{v.}{alcançar; chegar até | ser comparável a; alcançar (geralmente usado em termos negativos) | chegar a tempo para | estender"-se a; cuidar de; envolver | dar}
  \end{Phonetics}
\end{Entry}

\begin{Entry}{及早}{3,6}{⼃,⽇}
  \begin{Phonetics}{及早}{ji2zao3}[][HSK 7-9]
    \definition{adv.}{o mais cedo possível; antes que seja tarde demais}
  \synonymref{趁早}{chen4zao3}
  \synonymref{赶早}{gan3zao3}
  \end{Phonetics}
\end{Entry}

\begin{Entry}{及时}{3,7}{⼃,⽇}
  \begin{Phonetics}{及时}{ji2shi2}[][HSK 3]
    \definition{adj.}{oportuno; na hora certa; adequado; na ocasião certa}
    \definition{adv.}{prontamente; sem demora; imediatamente}
  \synonymref{随时}{sui2shi2}
  \antonymref{耽误}{dan1wu5}
  \end{Phonetics}
\end{Entry}

\begin{Entry}{及其}{3,8}{⼃,⼋}
  \begin{Phonetics}{及其}{ji2qi2}[][HSK 7-9]
    \definition{conj.}{(conjunção que liga dois substantivos) e seu\dots.; e seus\dots.; e dele\dots.; e dela\dots; usado para conectar duas ou mais coisas para indicar que elas são de igual importância ou existem da mesma maneira}[文化及其发展影响社会。===A cultura e seu desenvolvimento influenciam a sociedade.]
  \end{Phonetics}
\end{Entry}

\begin{Entry}{及格}{3,10}{⼃,⽊}
  \begin{Phonetics}{及格}{ji2/ge2}[][HSK 4]
    \definition{v.+compl.}{passar; passar em um teste, exame, etc.}
  \synonymref{合格}{he2ge2}
  \end{Phonetics}
\end{Entry}

%%%%%%%%%% 口 %%%%%%%%%%
\subsection*{口}\addcontentsline{loh}{figure}{口}

\begin{Entry}{口}{3}{⼝}[Kangxi 30]
  \begin{Phonetics}{口}{kou3}[][HSK 1]
    \definition*{s.}{Sobrenome: Kou}
    \definition{clas.}{usado para coisas com bocas (pessoas, animais domésticos, canhões, etc.) | usado para mordidas ou bocados | usado para idiomas}
    \definition{s.}{boca | borda; boca; o espaço externo ao recipiente | saída; entrada; local de entrada e saída | o gosto de alguém | corte; buraco; ferida |  a borda de uma faca; lâminas de facas, espadas, tesouras, etc. | a idade de um animal de tração | seção; departamento; sistema integrado de departamentos relacionados | conversa, discurso; pronunciamento; referência à fala | um portão da Grande Muralha (frequentemente usado em nomes de lugares)}
  \seealsoref{名}{ming2}
  \synonymref{个}{ge4}
  \synonymref{味}{wei4}
  \synonymref{嘴}{zui3}
  \end{Phonetics}
\end{Entry}

\begin{Entry}{口子}{3,3}{⼝,⼦}
  \begin{Phonetics}{口子}{kou3zi5}[][HSK 7-9]
    \definition{clas.}{referente a pessoas}[你家有几口子？===Quantas pessoas há na sua família?]
    \definition{s.}{Coloquial: pessoa | marido ou esposa | abertura; buraco; corte; rasgo; uma grande lacuna; uma fenda | Figurativo: abertura; oportunidade}
  \end{Phonetics}
\end{Entry}

\begin{Entry}{口才}{3,3}{⼝,⼿}
  \begin{Phonetics}{口才}{kou3cai2}[][HSK 7-9]
    \definition{s.}{eloquência; persuasão; a capacidade de se expressar verbalmente; o talento para a fala}
  \end{Phonetics}
\end{Entry}

\begin{Entry}{口气}{3,4}{⼝,⽓}
  \begin{Phonetics}{口气}{kou3qi4}[][HSK 7-9]
    \definition{s.}{maneira de falar; a força e o ritmo do tom; a dinâmica do discurso | implicação; o que realmente se quer dizer; o significado não dito nas palavras | tom; nota; tom emocional na fala}
  \end{Phonetics}
\end{Entry}

\begin{Entry}{口水}{3,4}{⼝,⽔}
  \begin{Phonetics}{口水}{kou3shui3}[][HSK 7-9]
    \definition{s.}{saliva; baba; o termo geral para saliva}
  \end{Phonetics}
\end{Entry}

\begin{Entry}{口令}{3,5}{⼝,⼈}
  \begin{Phonetics}{口令}{kou3ling4}[][HSK 7-9]
    \definition[串]{s.}{palavra de comando | senha; palavra"-chave; contra"-senha | palavra; comando; lema; comando verbal | códigos verbais para identificar amigo ou inimigo}
  \end{Phonetics}
\end{Entry}

\begin{Entry}{口号}{3,5}{⼝,⼝}
  \begin{Phonetics}{口号}{kou3hao4}[][HSK 5]
    \definition[个,条,些]{s.}{\emph{slogan}; palavra de ordem; lema}
  \end{Phonetics}
\end{Entry}

\begin{Entry}{口头}{3,5}{⼝,⼤}
  \begin{Phonetics}{口头}{kou3tou2}[][HSK 7-9]
    \definition{s.}{oral; verbal; expressa"-se através da fala | boca (referindo"-se à boca ao falar); lábios}
  \synonymref{表面}{biao3mian4}
  \antonymref{思想}{si1xiang3}
  \antonymref{行动}{xing2dong4}
  \antonymref{行为}{xing2wei2}
  \end{Phonetics}
\end{Entry}

\begin{Entry}{口吃}{3,6}{⼝,⼝}
  \begin{Phonetics}{口吃}{kou3chi1}[][HSK 7-9]
    \definition{s.}{gagueira; espasmofemia; balbucinato; mogilalia; battarismo; battarismo; iscnofonia; pselismo; o fenômeno de repetir palavras ou interromper frases ao falar é um defeito habitual de linguagem comumente conhecido como gagueira}
  \end{Phonetics}
\end{Entry}

\begin{Entry}{口吃病}{3,6,10}{⼝,⼝,⽧}
  \begin{Phonetics}{口吃病}{kou3chi1 bing4}
    \definition{s.}{doença da gagueira}
  \end{Phonetics}
\end{Entry}

\begin{Entry}{口味}{3,8}{⼝,⼝}
  \begin{Phonetics}{口味}{kou3wei4}[][HSK 7-9]
    \definition[个,种]{s.}{sabor da comida; o gosto da comida | o gosto de uma pessoa; preferência de cada um em termos de sabor | gostos; metáfora para interesses e hobbies pessoais}
  \end{Phonetics}
\end{Entry}

\begin{Entry}{口径}{3,8}{⼝,⼻}
  \begin{Phonetics}{口径}{kou3jing4}[][HSK 7-9]
    \definition{s.}{calibre; diâmetro; diâmetro da boca circular do vaso | requisitos; especificações; geralmente se refere às especificações, desempenho, etc., exigidos | declaração; pensamento; linha de ação ou fala; conteúdo metafórico da fala}
  \end{Phonetics}
\end{Entry}

\begin{Entry}{口试}{3,8}{⼝,⾔}
  \begin{Phonetics}{口试}{kou3shi4}[][HSK 6]
    \definition{s.}{exame oral (ou teste); um tipo de exame que exige que os candidatos respondam a perguntas oralmente}
    \definition{v.}{examinar oralmente}
  \antonymref{笔试}{bi3shi4}
  \end{Phonetics}
\end{Entry}

\begin{Entry}{口语}{3,9}{⼝,⾔}
  \begin{Phonetics}{口语}{kou3yu3}[][HSK 4]
    \definition[门]{s.}{linguagem oral; linguagem falada; linguagem coloquial; linguagem usada em conversas}
  \end{Phonetics}
\end{Entry}

\begin{Entry}{口音}{3,9}{⼝,⾳}
  \begin{Phonetics}{口音}{kou3yin1}[][HSK 7-9]
    \definition[种]{s.}{sotaques; sons da fala oral (linguística); os hábitos de fala de uma pessoa, especialmente seus hábitos de pronúncia, podem revelar sua origem linguística}
  \end{Phonetics}
  \begin{Phonetics}{口音}{kou3yin5}
    \definition[种]{s.}{voz | sotaque; dialeto}
  \end{Phonetics}
\end{Entry}

\begin{Entry}{口香糖}{3,9,16}{⼝,⾹,⽶}
  \begin{Phonetics}{口香糖}{kou3xiang1tang2}[][HSK 7-9]
    \definition{s.}{goma de mascar; chiclete; um tipo de doce feito da seiva viscosa que escorre do tronco da sapotilha (uma árvore perene cujos frutos têm o tamanho de peras e formato de coração, daí o nome), juntamente com açúcar e aromatizantes; é para ser mastigado apenas e não deve ser engolido}
  \end{Phonetics}
\end{Entry}

\begin{Entry}{口哨}{3,10}{⼝,⼝}
  \begin{Phonetics}{口哨}{kou3shao4}[][HSK 7-9]
    \definition{s.}{apito; assobio}
    \definition{v.}{assobiar}
  \end{Phonetics}
\end{Entry}

\begin{Entry}{口袋}{3,11}{⼝,⾐}
  \begin{Phonetics}{口袋}{kou3dai5}[][HSK 4]
    \definition[个,只]{s.}{bolso | saco; sacola; artigos de tecido ou couro}
  \end{Phonetics}
\end{Entry}

\begin{Entry}{口袋妖怪}{3,11,7,8}{⼝,⾐,⼥,⼼}
  \begin{Phonetics}{口袋妖怪}{kou3dai4 yao1guai4}
    \definition*{s.}{Pokémon (franquia de mídia japonesa)}
  \end{Phonetics}
\end{Entry}

\begin{Entry}{口腔}{3,12}{⼝,⾁}
  \begin{Phonetics}{口腔}{kou3qiang1}[][HSK 7-9]
    \definition{s.}{cavidade oral; a cavidade oral é um espaço oco composto pelos lábios, bochechas, palato duro e palato mole; contém órgãos como dentes, língua e glândulas salivares}
  \end{Phonetics}
\end{Entry}

\begin{Entry}{口感}{3,13}{⼝,⼼}
  \begin{Phonetics}{口感}{kou3gan3}[][HSK 7-9]
    \definition{s.}{textura (dos alimentos); sabor; sensação que o alimento proporciona na boca}
  \end{Phonetics}
\end{Entry}

\begin{Entry}{口碑}{3,13}{⼝,⽯}
  \begin{Phonetics}{口碑}{kou3bei1}[][HSK 7-9]
    \definition{s.}{elogio público; isso se refere à avaliação verbal que as pessoas fazem de alguém (antigamente, elogios a uma pessoa eram frequentemente gravados em tábuas de pedra)}
  \end{Phonetics}
\end{Entry}

\begin{Entry}{口罩}{3,13}{⼝,⽹}
  \begin{Phonetics}{口罩}{kou3zhao4}[][HSK 7-9]
    \definition[副]{s.}{máscara cirúrgica; produtos de higiene; feitos de gaze, etc.; usados sobre a boca e o nariz para evitar a entrada de poeira e germes}
  \end{Phonetics}
\end{Entry}

%%%%%%%%%% 土 %%%%%%%%%%
\subsection*{土}\addcontentsline{loh}{figure}{土}

\begin{Entry}{土}{3}{⼟}[Kangxi 32]
  \begin{Phonetics}{土}{tu3}[][HSK 3,6]
    \definition*{s.}{Sobrenome: Tu}
    \definition{adj.}{local; nativo; local com características regionais| caseiro; indígena; o que é tradicional no país; popular | não refinado; não esclarecido; não está na moda; não é popular}
    \definition[堆,捧,层]{s.}{solo; terra | terra; território | ópio bruto | cidade natal; terra natal; pátria}
  \end{Phonetics}
\end{Entry}

\begin{Entry}{土生土长}{3,5,3,4}{⼟,⽣,⼟,⾧}
  \begin{Phonetics}{土生土长}{tu3sheng1-tu3zhang3}[][HSK 7-9]
    \definition{expr.}{nativo; nascido e criado; cultivado localmente}
  \end{Phonetics}
\end{Entry}

\begin{Entry}{土地}{3,6}{⼟,⼟}
  \begin{Phonetics}{土地}{tu3di4}[][HSK 4]
    \definition[片,块,顷]{s.}{terra; solo; chão; superfície terrestre da Terra usada para cultivar, construir edifícios e viver | território; território de um país}
  \end{Phonetics}
  \begin{Phonetics}{土地}{tu3di5}
    \definition[片,块,顷]{s.}{deus da audeia; deus local; \emph{genius loci} deidade protetora de um local; Superstição: refere"-se ao deus da terra que governa uma pequena área}
  \end{Phonetics}
\end{Entry}

\begin{Entry}{土豆}{3,7}{⼟,⾖}
  \begin{Phonetics}{土豆}{tu3dou4}[][HSK 5]
    \definition[颗,斤,个,棵]{s.}{batata; denominação comum da batata}
  \end{Phonetics}
\end{Entry}

\begin{Entry}{土豆泥}{3,7,8}{⼟,⾖,⽔}
  \begin{Phonetics}{土豆泥}{tu3dou4ni2}
    \definition{s.}{purê de batata}[她的土豆泥确实不错。===O purê de batatas dela estava realmente muito bom.]
  \end{Phonetics}
\end{Entry}

\begin{Entry}{土鸡}{3,7}{⼟,⿃}
  \begin{Phonetics}{土鸡}{tu3ji1}
    \definition[只]{s.}{galinha criada em casa; galinha caipira}
  \end{Phonetics}
\end{Entry}

\begin{Entry}{土匪}{3,10}{⼟,⼕}
  \begin{Phonetics}{土匪}{tu3fei3}[][HSK 7-9]
    \definition{s.}{bandido; salteador}
  \seealsoref{胡匪}{hu2fei3}
  \synonymref{强盗}{qiang2dao4}
  \end{Phonetics}
\end{Entry}

\begin{Entry}{土家}{3,10}{⼟,⼧}
  \begin{Phonetics}{土家}{tu3jia1}
    \definition*{s.}{Tujia}
  \end{Phonetics}
\end{Entry}

\begin{Entry}{土壤}{3,20}{⼟,⼟}
  \begin{Phonetics}{土壤}{tu3rang3}[][HSK 7-9]
    \definition{s.}{solo; uma camada solta de material na superfície terrestre, composta por diversos minerais granulares, matéria orgânica, água, ar, microrganismos, etc., que permite o crescimento de plantas}
  \synonymref{泥土}{ni2tu3}
  \end{Phonetics}
\end{Entry}

%%%%%%%%%% 士 %%%%%%%%%%
\subsection*{士}\addcontentsline{loh}{figure}{士}

\begin{Entry}{士}{3}{⼠}[Kangxi 33]
  \begin{Phonetics}{士}{shi4}
    \definition*{s.}{Sobrenome: Shi}
    \definition[位,名,个]{s.}{soldado; militar | oficial não comissionado; primeira classe de soldados | pessoa treinada em uma determinada área; algum tipo de técnico | pessoa (louvável) | bacharel (na China antiga) | classe social, entre os oficiais, 大夫, e o povo comum, 庶民 | estudioso | guarda"-costas, uma das peças do xadrez chinês}
  \seealsoref{大夫}{da4fu1}
  \seealsoref{庶民}{shu4min2}
  \end{Phonetics}
\end{Entry}

\begin{Entry}{士气}{3,4}{⼠,⽓}
  \begin{Phonetics}{士气}{shi4qi4}[][HSK 7-9]
    \definition{s.}{moral; o espírito de luta do exército e, de forma mais ampla, o espírito de luta das massas}
  \synonymref{斗志}{dou4zhi4}
  \synonymref{气魄}{qi4po4}
  \synonymref{气势}{qi4shi4}
  \synonymref{勇气}{yong3qi4}
  \end{Phonetics}
\end{Entry}

\begin{Entry}{士兵}{3,7}{⼠,⼋}
  \begin{Phonetics}{士兵}{shi4bing1}[][HSK 4]
    \definition[个,名,位,批,群]{s.}{soldado; militar; termo coletivo para oficiais não comissionados e soldados; os membros mais jovens do exército}
  \end{Phonetics}
\end{Entry}

%%%%%%%%%% 夕 %%%%%%%%%%
\subsection*{夕}\addcontentsline{loh}{figure}{夕}

\begin{Entry}{夕}{3}{⼣}[Kangxi 36]
  \begin{Phonetics}{夕}{xi1}
    \definition*{s.}{Sobrenome: Xi}
    \definition{s.}{pôr do sol; crepúsculo | tarde; noite}
  \end{Phonetics}
\end{Entry}

\begin{Entry}{夕阳}{3,6}{⼣,⾩}
  \begin{Phonetics}{夕阳}{xi1yang2}
    \definition{s.}{pôr do sol | ocaso; metaforicamente, refere"-se à velhice; ou a indústrias tradicionais e não competitivas, sem perspectivas de futuro}
  \seealsoref{日出}{ri4chu1}
  \synonymref{落日}{luo4ri4}
  \synonymref{斜阳}{xie2yang2}
  \antonymref{旭日}{xu4ri4}
  \end{Phonetics}
\end{Entry}

%%%%%%%%%% 大 %%%%%%%%%%
\subsection*{大}\addcontentsline{loh}{figure}{大}

\begin{Entry}{大}{3}{⼤}[Kangxi 37]
  \begin{Phonetics}{大}{da4}[][HSK 1]
    \definition*{s.}{Sobrenome: Da}
    \definition{adj.}{grande; amplo; grande em volume, área, etc. | mais velho; em primeiro lugar no ranking | tamanho; descreve o grau de grandeza | usado em certas épocas do ano, condições climáticas, feriados ou antes de um determinado momento, para enfatizar | o tempo mais distante; há muito tempo}
    \definition{adv.}{grandemente; totalmente; expressa um grau muito profundo | não muito; não frequentemente; usado após 不, indica um grau baixo ou poucas vezes}
    \definition{s.}{adulto; crescido; pessoas idosas | pai | irmão do pai de alguém; tio}
  \seealsoref{不}{bu4}
  \synonymref{广}{guang3}
  \synonymref{巨}{ju4}
  \synonymref{阔}{kuo4}
  \synonymref{庞}{pang2}
  \synonymref{硕}{shuo4}
  \antonymref{小}{xiao3}
  \end{Phonetics}
  \begin{Phonetics}{大}{dai4}
    \definition{s.}{usado em 大夫: médico, doutor | usado em 大王: grande rei}
  \seealsoref{大夫}{dai4fu5}
  \seealsoref{大王}{dai4wang5}
  \end{Phonetics}
\end{Entry}

\begin{Entry}{大人}{3,2}{⼤,⼈}
  \begin{Phonetics}{大人}{da4ren2}[][HSK 2]
    \definition[个,位]{s.}{senhor; ilustre; sua excelência; antigo título honorífico para funcionários públicos | adulto; crescido; maduro;}
  \end{Phonetics}
\end{Entry}

\begin{Entry}{大力}{3,2}{⼤,⼒}
  \begin{Phonetics}{大力}{da4li4}[][HSK 6]
    \definition{adv.}{energicamente; vigorosamente; indica uso de grande força}
    \definition{s.}{grande força, poder}
  \end{Phonetics}
\end{Entry}

\begin{Entry}{大于}{3,3}{⼤,⼆}
  \begin{Phonetics}{大于}{da4yu2}[][HSK 5]
    \definition{v.}{ser maior, mais numeroso, mais importante, etc. do que}
  \end{Phonetics}
\end{Entry}

\begin{Entry}{大口}{3,3}{⼤,⼝}
  \begin{Phonetics}{大口}{da4 kou3}
    \definition{s.}{grande bocado (de comida, bebida, fumaça etc.); macrostoma | boca aberta; boquiaberto | devoração | deglutição}
  \end{Phonetics}
\end{Entry}

\begin{Entry}{大大}{3,3}{⼤,⼤}
  \begin{Phonetics}{大大}{da4da5}[][HSK 2]
    \definition{adv.}{grandemente; enormemente; enfatizar grande quantidade ou grau profundo}
  \end{Phonetics}
\end{Entry}

\begin{Entry}{大大咧咧}{3,3,9,9}{⼤,⼤,⼝,⼝}
  \begin{Phonetics}{大大咧咧}{da4da5lie1lie1}[][HSK 7-9]
    \definition{expr.}{despreocupado; descuidado; casual}
  \end{Phonetics}
\end{Entry}

\begin{Entry}{大小}{3,3}{⼤,⼩}
  \begin{Phonetics}{大小}{da4xiao3}[][HSK 2]
    \definition{adv.}{no mínimo; grande ou pequeno (geralmente pequeno), significa que ainda pode ser considerado}
    \definition[家]{s.}{tamanho; o grau de tamanho | ordem de senioridade; hierarquia | adultos e crianças | grande ou pequeno}
  \end{Phonetics}
\end{Entry}

\begin{Entry}{大门}{3,3}{⼤,⾨}
  \begin{Phonetics}{大门}{da4men2}[][HSK 2]
    \definition{s.}{portão; entrada; portão grande, referindo"-se especificamente ao portão principal de um edifício (como uma casa, pátio ou parque) que dá para a rua (em contraste com o segundo portão e as portas das várias divisões)}
  \end{Phonetics}
\end{Entry}

\begin{Entry}{大马}{3,3}{⼤,⾺}
  \begin{Phonetics}{大马}{da4ma3}
    \definition*{s.}{Malásia}
  \end{Phonetics}
\end{Entry}

\begin{Entry}{大公无私}{3,4,4,7}{⼤,⼋,⽆,⽲}
  \begin{Phonetics}{大公无私}{da4gong1-wu2si1}[][HSK 7-9]
    \definition{expr.}{altruísta; generoso; desinteressado | perfeitamente imparcial | nenhuma consideração pessoal (egoísta); não pensar em si mesmo; justo e altruísta; grande imparcialidade exclui consideração de si mesmo; desinteresse e altruísmo}
  \end{Phonetics}
\end{Entry}

\begin{Entry}{大厅}{3,4}{⼤,⼚}
  \begin{Phonetics}{大厅}{da4ting1}[][HSK 5]
    \definition{s.}{\emph{hall}; saguão, uma sala grande para reuniões ou atividades em um edifício de grande porte}
  \end{Phonetics}
\end{Entry}

\begin{Entry}{大夫}{3,4}{⼤,⼤}
  \begin{Phonetics}{大夫}{da4fu1}
    \definition[个,位,名]{s.}{oficial sênior (na China Imperial)}
  \end{Phonetics}
  \begin{Phonetics}{大夫}{dai4fu5}[][HSK 3]
    \definition[个,位,名]{s.}{médico, doutor}
  \end{Phonetics}
\end{Entry}

\begin{Entry}{大巴}{3,4}{⼤,⼰}
  \begin{Phonetics}{大巴}{da4ba1}[][HSK 4]
    \definition{s.}{ônibus}
  \end{Phonetics}
\end{Entry}

\begin{Entry}{大方}{3,4}{⼤,⽅}
  \begin{Phonetics}{大方}{da4fang1}
    \definition{s.}{generosidades; liberalidades | estudioso; pessoas com conhecimento especializado | um tipo de chá verde, produzido principalmente no Condado de Shexian, Província de Anhui, Condado de Chun'an, Província de Zhejiang, etc.}
  \end{Phonetics}
  \begin{Phonetics}{大方}{da4fang5}[][HSK 4]
    \definition{adj.}{generoso | não afetado; natural e equilibrado |  de bom gosto}
  \end{Phonetics}
\end{Entry}

\begin{Entry}{大气}{3,4}{⼤,⽓}
  \begin{Phonetics}{大气}{da4qi4}[][HSK 7-9]
    \definition{adj.}{generoso; magnânimo; refere"-se a um porte extraordinário, não convencional}
    \definition{s.}{ar; atmosfera | respiração pesada}
  \end{Phonetics}
\end{Entry}

\begin{Entry}{大片}{3,4}{⼤,⽚}
  \begin{Phonetics}{大片}{da4pian4}[][HSK 7-9]
    \definition{adj.}{grande; enorme; imenso}
    \definition{s.}{um filme de grande orçamento; um sucesso de bilheteria; \emph{block"-buster}; refere"-se a um filme caro e bem feito | grande área; vasta extensão; ampla extensão}
  \end{Phonetics}
\end{Entry}

\begin{Entry}{大王}{3,4}{⼤,⽟}
  \begin{Phonetics}{大王}{da4wang2}
    \definition{s.}{rei; magnata | pessoa da mais alta classe ou habilidade em algo; ás | barões | pessoa com habilidade especializada em algo}
  \end{Phonetics}
  \begin{Phonetics}{大王}{dai4wang5}
    \definition{s.}{magnata; barões | barão ladrão (em ópera, histórias antigas)}
  \end{Phonetics}
\end{Entry}

\begin{Entry}{大队}{3,4}{⼤,⾩}
  \begin{Phonetics}{大队}{da4dui4}[][HSK 7-9]
    \definition{adj.}{grande (contingente de tropas, corpo de manifestantes, número de pessoas, etc.)}
    \definition{s.}{unidade militar correspondente ao batalhão ou regimento; grupo | História: brigada de produção (de uma comuna popular rural); brigada}
  \end{Phonetics}
\end{Entry}

\begin{Entry}{大包大揽}{3,5,3,12}{⼤,⼓,⼤,⼿}
  \begin{Phonetics}{大包大揽}{da4bao1-da4lan3}[][HSK 7-9]
    \definition{v.}{assumir o controle total}
  \end{Phonetics}
\end{Entry}

\begin{Entry}{大众}{3,6}{⼤,⼈}
  \begin{Phonetics}{大众}{da4zhong4}[][HSK 4]
    \definition{s.}{massas; população; pessoas comuns; público em geral}
  \end{Phonetics}
\end{Entry}

\begin{Entry}{大伙儿}{3,6,2}{⼤,⼈,⼉}
  \begin{Phonetics}{大伙儿}{da4huo3r5}[][HSK 5]
    \definition{pron.}{todos nós; todos vocês; todo mundo; todos; equivalente a 大家}
  \seealsoref{大家}{da4jia1}
  \end{Phonetics}
\end{Entry}

\begin{Entry}{大会}{3,6}{⼤,⼈}
  \begin{Phonetics}{大会}{da4hui4}[][HSK 4]
    \definition[场,次,个,届]{s.}{sessão plenária; reunião geral de membros; reuniões convocadas por partidos políticos socialistas | reunião de massa; comício de massa}
  \end{Phonetics}
\end{Entry}

\begin{Entry}{大全}{3,6}{⼤,⼊}
  \begin{Phonetics}{大全}{da4quan2}
    \definition{s.}{coleção completa; obras completas | obras completas de; uma coleção completa de; um volume completo sobre; refere"-se a conteúdo rico e completo, e é frequentemente usado em títulos de livros, como ``Coleção Completa de Artigos de Uso Diário Rural'' 《农村日用大全》 e ``Coleção Completa de Ópera Chinesa'' 《中国戏曲大全》}
  \end{Phonetics}
\end{Entry}

\begin{Entry}{大吃一惊}{3,6,1,11}{⼤,⼝,⼀,⼼}
  \begin{Phonetics}{大吃一惊}{da4chi1-yi1jing1}[][HSK 7-9]
    \definition{expr.}{ficar assustado com; ficar espantado com; ficar completamente surpreso; ficar muito surpreso; ficar completamente chocado; ser pego de surpresa; levar um choque; ter (levar) um susto; levar um susto com; ficar pasmo}
  \end{Phonetics}
\end{Entry}

\begin{Entry}{大同小异}{3,6,3,6}{⼤,⼝,⼩,⼶}
  \begin{Phonetics}{大同小异}{da4tong2-xiao3yi4}[][HSK 7-9]
    \definition{expr.}{muito parecidos, mas com pequenas diferenças; semelhante, exceto por pequenas diferenças; ser o mesmo em aspectos essenciais, embora diferindo em pontos menores; uma semelhança geral com pequenas (menores) diferenças; ser semelhante em todos os aspectos essenciais importantes, sendo as diferenças de natureza menor; ser em grande parte idêntico com apenas pequenas diferenças; diferir apenas em (em) pequenos pontos; essencialmente o mesmo, embora diferindo em pontos menores; na maior parte, eles são os mesmos; idênticos em questões importantes, embora com pequenas diferenças; principalmente semelhantes, exceto (por) pequenas diferenças; virtualmente os mesmos}
  \end{Phonetics}
\end{Entry}

\begin{Entry}{大名鼎鼎}{3,6,12,12}{⼤,⼝,⿍,⿍}
  \begin{Phonetics}{大名鼎鼎}{da4ming2-ding3ding3}[][HSK 7-9]
    \definition{expr.}{``O nome de alguém é conhecido em toda parte.''; ser muito famoso; ser amplamente conhecido; celebrado; desfrutar de um grande nome; bem conhecido}
  \end{Phonetics}
\end{Entry}

\begin{Entry}{大后天}{3,6,4}{⼤,⼝,⼤}
  \begin{Phonetics}{大后天}{da4hou4tian1}
    \definition{s.}{daqui a três dias}
  \end{Phonetics}
\end{Entry}

\begin{Entry}{大地}{3,6}{⼤,⼟}
  \begin{Phonetics}{大地}{da4di4}[][HSK 7-9]
    \definition*{s.}{A Terra}[大地正在遭受污染。===A Terra está sendo poluída.]
    \definition[片,块,方]{s.}{terreno espaçoso; terra}
  \end{Phonetics}
\end{Entry}

\begin{Entry}{大多}{3,6}{⼤,⼣}
  \begin{Phonetics}{大多}{da4duo1}[][HSK 4]
    \definition{adv.}{majoritariamente; em sua maior parte; em sua maioria; em grande parte}
  \end{Phonetics}
\end{Entry}

\begin{Entry}{大多数}{3,6,13}{⼤,⼣,⽁}
  \begin{Phonetics}{大多数}{da4duo1shu4}[][HSK 2]
    \definition{s.}{grande maioria; vasta maioria; a maior parte; mais da metade, um número significativo}
  \end{Phonetics}
\end{Entry}

\begin{Entry}{大妈}{3,6}{⼤,⼥}
  \begin{Phonetics}{大妈}{da4ma1}[][HSK 4]
    \definition[个,位]{s.}{tia; esposa do irmão mais velho do pai | tratamento respeitoso às mulheres idosas}
  \end{Phonetics}
\end{Entry}

\begin{Entry}{大师}{3,6}{⼤,⼱}
  \begin{Phonetics}{大师}{da4shi1}[][HSK 6]
    \definition*{s.}{Grande Mestre, título de cortesia usado para se dirigir a um monge budista}
    \definition{s.}{grande mestre; mestre; maestro; uma pessoa com realizações profundas}
  \end{Phonetics}
\end{Entry}

\begin{Entry}{大戏}{3,6}{⼤,⼽}
  \begin{Phonetics}{大戏}{da4xi4}
    \definition*{s.}{Dialeto: Ópera de Pequim}
    \definition{s.}{ópera tradicional em grande escala; grande produção dramática (filme, série de TV etc.)}
  \end{Phonetics}
\end{Entry}

\begin{Entry}{大有可为}{3,6,5,4}{⼤,⽉,⼝,⼂}
  \begin{Phonetics}{大有可为}{da4you3-ke3wei2}[][HSK 7-9]
    \definition{expr.}{ter um futuro brilhante; valer a pena; poder realizar grandes coisas; ter perspectivas brilhantes; pode ser bem administrado; há amplo escopo para nossas habilidades em\dots; valer a pena fazer; ter perspectivas brilhantes}
  \end{Phonetics}
\end{Entry}

\begin{Entry}{大爷}{3,6}{⼤,⽗}
  \begin{Phonetics}{大爷}{da4ye2}
    \definition[个,位]{s.}{Coloquial: irmão mais velho do pai; tio | tratamento respeitoso para um homem mais velho}
  \end{Phonetics}
  \begin{Phonetics}{大爷}{da4ye5}[][HSK 4]
    \definition[个,位]{s.}{irmão mais velho do pai; tio | tio (homenagem aos homens mais velhos)}
  \end{Phonetics}
\end{Entry}

\begin{Entry}{大米}{3,6}{⼤,⽶}
  \begin{Phonetics}{大米}{da4mi3}[][HSK 6]
    \definition[颗,粒,斤,包,袋]{s.}{arroz; arroz descascado; arroz bom}
  \end{Phonetics}
\end{Entry}

\begin{Entry}{大约}{3,6}{⼤,⽷}
  \begin{Phonetics}{大约}{da4yue1}[][HSK 3]
    \definition{adv.}{aproximadamente; sobre; estimativa não muito precisa | provavelmente; expressar suposições sobre a situação}
  \end{Phonetics}
\end{Entry}

\begin{Entry}{大臣}{3,6}{⼤,⾂}
  \begin{Phonetics}{大臣}{da4chen2}[][HSK 7-9]
    \definition[位,个]{s.}{secretário; ministro (de uma monarquia); altos funcionários da monarquia}
  \end{Phonetics}
\end{Entry}

\begin{Entry}{大自然}{3,6,12}{⼤,⾃,⽕}
  \begin{Phonetics}{大自然}{da4zi4ran2}[][HSK 2]
    \definition{s.}{natureza}
  \end{Phonetics}
\end{Entry}

\begin{Entry}{大衣}{3,6}{⼤,⾐}
  \begin{Phonetics}{大衣}{da4yi1}[][HSK 2]
    \definition[件,个]{s.}{sobretudo; casaco; casaco ocidental mais comprido}
  \end{Phonetics}
\end{Entry}

\begin{Entry}{大伯子}{3,7,3}{⼤,⼈,⼦}
  \begin{Phonetics}{大伯子}{da4bai3zi5}
    \definition{s.}{Coloquial: irmão mais velho do marido; cunhado}
  \end{Phonetics}
\end{Entry}

\begin{Entry}{大体}{3,7}{⼤,⼈}
  \begin{Phonetics}{大体}{da4ti3}[][HSK 7-9]
    \definition{adv.}{aproximadamente; em geral; mais ou menos}
    \definition{s.}{princípio primordial; interesse geral; um princípio basilar baseado na situação geral; um princípio relativo à situação geral}
  \end{Phonetics}
\end{Entry}

\begin{Entry}{大体上}{3,7,3}{⼤,⼈,⼀}
  \begin{Phonetics}{大体上}{da4ti3shang4}[][HSK 7-9]
    \definition{adv.}{em geral; grosso modo; de um modo geral}
  \end{Phonetics}
\end{Entry}

\begin{Entry}{大声}{3,7}{⼤,⼠}
  \begin{Phonetics}{大声}{da4sheng1}[][HSK 2]
    \definition{adj.}{alto; volume alto; em voz alta}
  \end{Phonetics}
\end{Entry}

\begin{Entry}{大局}{3,7}{⼤,⼫}
  \begin{Phonetics}{大局}{da4ju2}[][HSK 7-9]
    \definition{s.}{situação geral (ou inteira); toda a situação}
  \end{Phonetics}
\end{Entry}

\begin{Entry}{大批}{3,7}{⼤,⼿}
  \begin{Phonetics}{大批}{da4pi1}[][HSK 6]
    \definition{num.}{grandes quantidades de; exércitos; inundações}[大批书籍被印刷出来。===Grandes quantidades de livros foram impressas.]
  \end{Phonetics}
\end{Entry}

\begin{Entry}{大纲}{3,7}{⼤,⽷}
  \begin{Phonetics}{大纲}{da4gang1}[][HSK 5]
    \definition{s.}{esboço; compêndio; programa de estudos; resumo; fundamentos da organização sistemática de conteúdos (livros, discursos, programas, etc.)}
  \end{Phonetics}
\end{Entry}

\begin{Entry}{大豆}{3,7}{⼤,⾖}
  \begin{Phonetics}{大豆}{da4dou4}
    \definition{s.}{soja; grão de soja}
  \end{Phonetics}
\end{Entry}

\begin{Entry}{大陆}{3,7}{⼤,⾩}
  \begin{Phonetics}{大陆}{da4lu4}[][HSK 4]
    \definition*{s.}{China continental; refere"-se especificamente à vasta porção terrestre do território da China}
    \definition[个,块]{s.}{terra firme; continente; vasta extensão de terra}
  \end{Phonetics}
\end{Entry}

\begin{Entry}{大事}{3,8}{⼤,⼅}
  \begin{Phonetics}{大事}{da4shi4}[][HSK 5]
    \definition{adv.}{em grande escala; em grande estilo; em grande parte}
    \definition[件,桩]{s.}{grande evento; grande acontecimento; assunto importante; grande questão; algo importante | situação geral}
  \end{Phonetics}
\end{Entry}

\begin{Entry}{大使}{3,8}{⼤,⼈}
  \begin{Phonetics}{大使}{da4shi3}[][HSK 6]
    \definition[位,任]{s.}{embaixador; o representante diplomático de mais alto nível enviado por um país a outro país}
  \seealsoref{全称特命全权大使}{quan2cheng1 te4ming4 quan2quan2da4shi3}
  \end{Phonetics}
\end{Entry}

\begin{Entry}{大使馆}{3,8,11}{⼤,⼈,⾷}
  \begin{Phonetics}{大使馆}{da4shi3guan3}[][HSK 3]
    \definition[座,个]{s.}{embaixada; uma representação diplomática de um país em outro país, chefiada por um embaixador}
  \end{Phonetics}
\end{Entry}

\begin{Entry}{大姐}{3,8}{⼤,⼥}
  \begin{Phonetics}{大姐}{da4jie3}[][HSK 4]
    \definition[个,位]{s.}{irmã mais velha (também um termo educado para se dirigir a uma garota ou mulher um pouco mais velha do que a pessoa que fala)}
  \end{Phonetics}
\end{Entry}

\begin{Entry}{大学}{3,8}{⼤,⼦}
  \begin{Phonetics}{大学}{da4xue2}[][HSK 1]
    \definition[所,座]{s.}{universidade; faculdade; tipo de instituição de ensino superior que, na China, geralmente se refere a uma universidade abrangente}
  \synonymref{学院}{xue2yuan4}
  \antonymref{小学}{xiao3xue2}
  \end{Phonetics}
\end{Entry}

\begin{Entry}{大学生}{3,8,5}{⼤,⼦,⽣}
  \begin{Phonetics}{大学生}{da4xue2sheng1}[][HSK 1]
    \definition[名,个,班]{s.}{estudante universitário; estudante de faculdade; estudantes de graduação ou cursos técnicos em instituições de ensino superior}
  \antonymref{小学生}{xiao3xue2sheng5}
  \end{Phonetics}
\end{Entry}

\begin{Entry}{大宗}{3,8}{⼤,⼧}
  \begin{Phonetics}{大宗}{da4zong1}[][HSK 7-9]
    \definition{adj.}{grande; volumoso}
    \definition{s.}{itens principais; produtos ou bens básicos}
  \end{Phonetics}
\end{Entry}

\begin{Entry}{大抵}{3,8}{⼤,⼿}
  \begin{Phonetics}{大抵}{da4di3}
    \definition{adv.}{no geral; de um modo geral; provavelmente; principalmente}
  \end{Phonetics}
\end{Entry}

\begin{Entry}{大规模}{3,8,14}{⼤,⾒,⽊}
  \begin{Phonetics}{大规模}{da4gui1mo2}[][HSK 4]
    \definition{adj.}{em larga escala; extensivo; maciço; massivo}
    \definition{adv.}{em larga escala; extensivo; maciço; massa}
  \end{Phonetics}
\end{Entry}

\begin{Entry}{大雨}{3,8}{⼤,⾬}
  \begin{Phonetics}{大雨}{da4yu3}
    \definition[场]{s.}{chuva pesada, forte}
  \synonymref{暴雨}{bao4yu3}
  \synonymref{下雨}{xia4/yu3}
  \antonymref{多云}{duo1yun2}
  \antonymref{晴朗}{qing2lang3}
  \antonymref{晴天}{qing2tian1}
  \end{Phonetics}
\end{Entry}

\begin{Entry}{大前儿}{3,9,2}{⼤,⼑,⼉}
  \begin{Phonetics}{大前儿}{da4qian2r5}
    \definition{s.}{três dias atrás}
  \end{Phonetics}
\end{Entry}

\begin{Entry}{大前天}{3,9,4}{⼤,⼑,⼤}
  \begin{Phonetics}{大前天}{da4qian2tian1}
    \definition{adv.}{três dias atrás}
  \seealsoref{大前儿}{da4qian2r5}
  \antonymref{大后天}{da4hou4tian1}
  \end{Phonetics}
\end{Entry}

\begin{Entry}{大型}{3,9}{⼤,⼟}
  \begin{Phonetics}{大型}{da4xing2}[][HSK 4]
    \definition{adj.}{grande; em larga escala; tamanho e volume grandes | larga escala (importante e influente)}
  \end{Phonetics}
\end{Entry}

\begin{Entry}{大城}{3,9}{⼤,⼟}
  \begin{Phonetics}{大城}{da4cheng2}
    \definition*[个,座]{s.}{Condado de Dacheng em Langfang 廊坊, Hebei | Município de Tacheng no condado de Changhua | Condado de Changhua, Taiwan}
  \seealsoref{廊坊}{lang2fang2}
  \end{Phonetics}
\end{Entry}

\begin{Entry}{大奖赛}{3,9,14}{⼤,⼤,⾙}
  \begin{Phonetics}{大奖赛}{da4jiang3sai4}[][HSK 5]
    \definition{s.}{grande competição; grande prêmio; \emph{grand prix}}
  \end{Phonetics}
\end{Entry}

\begin{Entry}{大姨}{3,9}{⼤,⼥}
  \begin{Phonetics}{大姨}{da4yi2}
    \definition{s.}{irmã mais velha da mãe; tia | Coloquial: irmã mais velha da esposa | cunhada}
  \seealsoref{大姨儿}{da4yi2r5}
  \synonymref{阿姨}{a1yi2}
  \end{Phonetics}
\end{Entry}

\begin{Entry}{大姨儿}{3,9,2}{⼤,⼥,⼉}
  \begin{Phonetics}{大姨儿}{da4yi2r5}
    \definition{s.}{tia}
  \end{Phonetics}
\end{Entry}

\begin{Entry}{大战}{3,9}{⼤,⼽}
  \begin{Phonetics}{大战}{da4zhan4}
    \definition{v.}{guerrear | lutar em uma guerra}
    \definition{v.}{travar guerra}
  \synonymref{战争}{zhan4zheng1}
  \end{Phonetics}
\end{Entry}

\begin{Entry}{大洋洲}{3,9,9}{⼤,⽔,⽔}
  \begin{Phonetics}{大洋洲}{da4yang2zhou1}
    \definition*{s.}{Oceania (Continente)}
  \end{Phonetics}
\end{Entry}

\begin{Entry}{大神}{3,9}{⼤,⽰}
  \begin{Phonetics}{大神}{da4shen2}
    \definition{s.}{Gíria: (Internet) guru | divindade | especialista; \emph{expert} | gênio}
  \synonymref{牛人}{niu2 ren2}
  \end{Phonetics}
\end{Entry}

\begin{Entry}{大禹}{3,9}{⼤,⽱}
  \begin{Phonetics}{大禹}{da4yu3}
    \definition*{s.}{Yu, o Grande (c. século XXI a.C.), líder mítico que domou as inundações}
  \end{Phonetics}
\end{Entry}

\begin{Entry}{大胆}{3,9}{⼤,⾁}
  \begin{Phonetics}{大胆}{da4dan3}[][HSK 5]
    \definition{adj.}{ousado; atrevido; audacioso; corajoso; destemido}
  \end{Phonetics}
\end{Entry}

\begin{Entry}{大选}{3,9}{⼤,⾡}
  \begin{Phonetics}{大选}{da4xuan3}[][HSK 7-9]
    \definition[次]{s.}{eleição geral; refere"-se à eleição de membros do parlamento ou presidente em alguns países}
  \end{Phonetics}
\end{Entry}

\begin{Entry}{大面积}{3,9,10}{⼤,⾯,⽲}
  \begin{Phonetics}{大面积}{da4mian4ji1}[][HSK 7-9]
    \definition{s.}{grande área}
  \end{Phonetics}
\end{Entry}

\begin{Entry}{大哥}{3,10}{⼤,⼝}
  \begin{Phonetics}{大哥}{da4ge1}[][HSK 4]
    \definition{s.}{irmão mais velho | \emph{big brother}; tratamento educado para um homem da mesma idade que você | líder de gangue; pessoa mais poderosa em uma organização que realiza atividades ilegais na sociedade}
  \end{Phonetics}
\end{Entry}

\begin{Entry}{大娘}{3,10}{⼤,⼥}
  \begin{Phonetics}{大娘}{da4niang2}
    \definition[个,位]{s.}{a esposa do irmão mais velho do pai; tia | tia; termo respeitoso para se referir a uma mulher mais velha}
  \synonymref{阿姨}{a1yi2}
  \synonymref{大妈}{da4ma1}
  \end{Phonetics}
\end{Entry}

\begin{Entry}{大家}{3,10}{⼤,⼧}
  \begin{Phonetics}{大家}{da4jia1}[][HSK 2]
    \definition{pron.}{todos; toda a gente; refere"-se a todas as pessoas dentro de um determinado âmbito}
    \definition{s.}{grande mestre; autoridade; especialista renomado | família nobre; família rica e influente; família tradicional}
  \end{Phonetics}
\end{Entry}

\begin{Entry}{大家庭}{3,10,9}{⼤,⼧,⼴}
  \begin{Phonetics}{大家庭}{da4jia1ting2}[][HSK 7-9]
    \definition[个,些]{s.}{grande família; comunidade | uma grande família é frequentemente usada para descrever um grupo com muitos membros e harmonia interna}
  \end{Phonetics}
\end{Entry}

\begin{Entry}{大海}{3,10}{⼤,⽔}
  \begin{Phonetics}{大海}{da4hai3}[][HSK 2]
    \definition{s.}{o mar; o oceano; o mar aberto, ou seja, a parte do oceano que não está fechada entre cabos nem incluída em estreitos}
  \end{Phonetics}
\end{Entry}

\begin{Entry}{大笔}{3,10}{⼤,⽵}
  \begin{Phonetics}{大笔}{da4bi3}[][HSK 7-9]
    \definition{adj.}{grande quantidade (de dinheiro) | grande soma (de capital, dinheiro, etc.)}
    \definition{s.}{caneta | Cortês: sua escrita; sua caligrafia}
  \end{Phonetics}
\end{Entry}

\begin{Entry}{大脑}{3,10}{⼤,⾁}
  \begin{Phonetics}{大脑}{da4nao3}[][HSK 5]
    \definition{s.}{cérebro; encéfalo}
  \end{Phonetics}
\end{Entry}

\begin{Entry}{大致}{3,10}{⼤,⾄}
  \begin{Phonetics}{大致}{da4zhi4}[][HSK 5]
    \definition{adj.}{geral; no todo}
    \definition{adv.}{grosso modo; aproximadamente; mais ou menos; indica uma estimativa aproximada da situação}
  \end{Phonetics}
\end{Entry}

\begin{Entry}{大部分}{3,10,4}{⼤,⾢,⼑}
  \begin{Phonetics}{大部分}{da4bu4fen5}[][HSK 2]
    \definition[把]{s.}{a maioria; a maior parte; em grande parte; refere"-se a uma quantidade superior a metade do total}
  \end{Phonetics}
\end{Entry}

\begin{Entry}{大都}{3,10}{⼤,⾢}
  \begin{Phonetics}{大都}{da4dou1}
  \end{Phonetics}
  \begin{Phonetics}{大都}{da4du1}[][HSK 5]
    \definition*{s.}{Dadu, capital da China durante a Dinastia Yuan (1280--1368), atual Pequim}
    \definition{adv.}{em sua maior parte; na maior parte; indica que a maioria das pessoas ou coisas em um determinado intervalo tem a mesma natureza e características; também pronunciado como \dpy{da4dou1} na língua falada}
  \end{Phonetics}
\end{Entry}

\begin{Entry}{大惊小怪}{3,11,3,8}{⼤,⼼,⼩,⼼}
  \begin{Phonetics}{大惊小怪}{da4jing1-xiao3guai4}[][HSK 7-9]
    \definition{expr.}{ficar animado com uma coisa pequena; uma tempestade em um copo d'água; ficar surpreso com algo perfeitamente normal; ficar alarmado com algo perfeitamente normal; sentir"-se surpreso; ficar alarmado com; ficar todo agitado por nada; grande alarme com um pequeno fantasma; fazer um alarido; fazer um alarido por nada; fazer um grande alarido sobre (algo); fazer um raro alarido sobre (algo); fazer muito alarido; fazer barulho por nada}
  \end{Phonetics}
\end{Entry}

\begin{Entry}{大象}{3,11}{⼤,⾗}
  \begin{Phonetics}{大象}{da4xiang4}[][HSK 5]
    \definition[只,头,群,个]{s.}{elefante}
  \end{Phonetics}
\end{Entry}

\begin{Entry}{大黄}{3,11}{⼤,⿈}
  \begin{Phonetics}{大黄}{da4huang2}
    \definition{s.}{ruibarbo chinês}
  \end{Phonetics}
\end{Entry}

\begin{Entry}{大厦}{3,12}{⼤,⼚}
  \begin{Phonetics}{大厦}{da4sha4}[][HSK 7-9]
    \definition[栋,座,撞]{s.}{prédio; edifício grande (frequentemente usado em nomes de grandes edifícios)}
  \end{Phonetics}
\end{Entry}

\begin{Entry}{大幅度}{3,12,9}{⼤,⼱,⼴}
  \begin{Phonetics}{大幅度}{da4fu2du4}[][HSK 7-9]
    \definition{adv.}{drasticamente; substancial; por uma ampla margem; a extensão da mudança ou alteração é significativa}
  \end{Phonetics}
\end{Entry}

\begin{Entry}{大棚}{3,12}{⼤,⽊}
  \begin{Phonetics}{大棚}{da4peng2}[][HSK 7-9]
    \definition{s.}{grande abrigo em arco coberto com lona plástica usada na agricultura; grande estufa | estufa}
  \end{Phonetics}
\end{Entry}

\begin{Entry}{大款}{3,12}{⼤,⽋}
  \begin{Phonetics}{大款}{da4kuan3}[][HSK 7-9]
    \definition[个]{s.}{o rico; pessoa muito rica}
  \end{Phonetics}
\end{Entry}

\begin{Entry}{大猩猩}{3,12,12}{⼤,⽝,⽝}
  \begin{Phonetics}{大猩猩}{da4xing1xing5}
    \definition{s.}{gorila}
  \end{Phonetics}
\end{Entry}

\begin{Entry}{大腕儿}{3,12,2}{⼤,⾁,⼉}
  \begin{Phonetics}{大腕儿}{da4wan4r5}[][HSK 7-9]
    \definition[位]{s.}{grande nome, figurão; uma pessoa que tem grande habilidade, fama e influência em um determinado setor ou aspecto}
  \end{Phonetics}
\end{Entry}

\begin{Entry}{大街}{3,12}{⼤,⾏}
  \begin{Phonetics}{大街}{da4jie1}[][HSK 6]
    \definition[条,个]{s.}{avenida; rua; rua principal}
  \end{Phonetics}
\end{Entry}

\begin{Entry}{大街小巷}{3,12,3,9}{⼤,⾏,⼩,⼰}
  \begin{Phonetics}{大街小巷}{da4jie1-xiao3xiang4}[][HSK 7-9]
    \definition{expr.}{grandes ruas e pequenos becos; em todos os lugares da cidade}
  \end{Phonetics}
\end{Entry}

\begin{Entry}{大道}{3,12}{⼤,⾡}
  \begin{Phonetics}{大道}{da4dao4}[][HSK 6]
    \definition*{s.}{O Grande Tao; O Grande Caminho}
    \definition[条]{s.}{estrada principal | o caminho da justiça | avenida | rua principal}
  \end{Phonetics}
\end{Entry}

\begin{Entry}{大量}{3,12}{⼤,⾥}
  \begin{Phonetics}{大量}{da4liang4}[][HSK 2]
    \definition{adj.}{numeroso; em grande quantidade; grande em número ou quantidade | generoso; magnânimo; descreve uma pessoa que não fica zangada quando os outros cometem erros e que costuma perdoar os outros}
  \end{Phonetics}
\end{Entry}

\begin{Entry}{大雁}{3,12}{⼤,⾫}
  \begin{Phonetics}{大雁}{da4yan4}[][HSK 7-9]
    \definition[只,群,行]{s.}{ganso selvagem}
  \end{Phonetics}
\end{Entry}

\begin{Entry}{大意}{3,13}{⼤,⼼}
  \begin{Phonetics}{大意}{da4yi4}[][HSK 7-9]
    \definition{s.}{essência; teor; ideia geral; pontos principais; o significado principal}
  \end{Phonetics}
  \begin{Phonetics}{大意}{da4yi5}[][HSK 7-9]
    \definition{adj.}{descuidado; negligente; desatento}
  \end{Phonetics}
\end{Entry}

\begin{Entry}{大数据}{3,13,11}{⼤,⽁,⼿}
  \begin{Phonetics}{大数据}{da4shu4ju4}[][HSK 7-9]
    \definition*{s.}{\emph{Big Data}}
  \end{Phonetics}
\end{Entry}

\begin{Entry}{大楼}{3,13}{⼤,⽊}
  \begin{Phonetics}{大楼}{da4lou2}[][HSK 4]
    \definition[座,幢]{s.}{edifício; mansão; edifício de vários andares disponível para uso residencial e comercial}
  \end{Phonetics}
\end{Entry}

\begin{Entry}{大概}{3,13}{⼤,⽊}
  \begin{Phonetics}{大概}{da4gai4}[][HSK 3]
    \definition{adj.}{geral; grosseiro; aproximado; não é muito preciso ou muito detalhado}
    \definition{adv.}{sobre; provavelmente; estimativas ou suposições imprecisas sobre eventos, quantidades, tempo, localização, etc. | geralmente; brevemente; não muito seriamente, casualmente; não muito cuidadosamente}
    \definition{s.}{ideia geral; esboço geral; conteúdo geral ou situação}
  \end{Phonetics}
\end{Entry}

\begin{Entry}{大肆}{3,13}{⼤,⾀}
  \begin{Phonetics}{大肆}{da4si4}[][HSK 7-9]
    \definition{adv.}{desenfreadamente; violentamente; imprudentemente; sem restrições; sem escrúpulos (geralmente referindo"-se a fazer coisas ruins)}
  \end{Phonetics}
\end{Entry}

\begin{Entry}{大腿}{3,13}{⼤,⾁}
  \begin{Phonetics}{大腿}{da4tui3}
    \definition{s.}{coxa}
  \end{Phonetics}
\end{Entry}

\begin{Entry}{大蒜}{3,13}{⼤,⾋}
  \begin{Phonetics}{大蒜}{da4suan4}
    \definition[瓣,头,把]{s.}{alho}
  \end{Phonetics}
\end{Entry}

\begin{Entry}{大模大样}{3,14,3,10}{⼤,⽊,⼤,⽊}
  \begin{Phonetics}{大模大样}{da4mu2-da4yang4}[][HSK 7-9]
    \definition{expr.}{pomposamente; corajosamente; ostensivamente; equilibrado; autoconfiante}
  \end{Phonetics}
\end{Entry}

\begin{Entry}{大熊猫}{3,14,11}{⼤,⽕,⽝}
  \begin{Phonetics}{大熊猫}{da4xiong2mao1}[][HSK 5]
    \definition{s.}{panda gigante (\emph{Ailuropoda melanoleuca})}
  \synonymref{猫熊}{mao1xiong2}
  \synonymref{熊猫}{xiong2mao1}
  \end{Phonetics}
\end{Entry}

\begin{Entry}{大赛}{3,14}{⼤,⾙}
  \begin{Phonetics}{大赛}{da4sai4}[][HSK 6]
    \definition{s.}{grande torneio; competição importante; um evento de grande porte e alto nível; um grande evento}
  \end{Phonetics}
\end{Entry}

%%%%%%%%%% 女 %%%%%%%%%%
\subsection*{女}\addcontentsline{loh}{figure}{女}

\begin{Entry}{女}{3}{⼥}[Kangxi 38]
  \begin{Phonetics}{女}{nv3}[][HSK 1]
    \definition{adj.}{mulher; feminino | fêmea (de certos animais)}
    \definition{s.}{menina; filha | nü, uma das mansões lunares | mulher}
  \synonymref{妇女}{fu4nv3}
  \synonymref{姑娘}{gu1niang5}
  \antonymref{男}{nan2}
  \end{Phonetics}
\end{Entry}

\begin{Entry}{女人}{3,2}{⼥,⼈}
  \begin{Phonetics}{女人}{nv3ren5}[][HSK 1]
    \definition[个,位]{s.}{mulher adulta}
  \end{Phonetics}
\end{Entry}

\begin{Entry}{女儿}{3,2}{⼥,⼉}
  \begin{Phonetics}{女儿}{nv3'er2}[][HSK 1]
    \definition[个,位]{s.}{menina; filha}
  \seealsoref{儿子}{er2zi5}
  \synonymref{姑娘}{gu1niang5}
  \synonymref{闺女}{gui1nv5}
  \synonymref{丫头}{ya1tou5}
  \antonymref{儿子}{er2zi5}
  \end{Phonetics}
\end{Entry}

\begin{Entry}{女士}{3,3}{⼥,⼠}
  \begin{Phonetics}{女士}{nv3shi4}[][HSK 4]
    \definition{pron.}{Sra.; Senhorita; Senhora; título honorífico para mulheres (agora usado em contextos diplomáticos)}
    \definition[位,名,个,些]{s.}{senhora; madame}
  \end{Phonetics}
\end{Entry}

\begin{Entry}{女子}{3,3}{⼥,⼦}
  \begin{Phonetics}{女子}{nv3zi3}[][HSK 3]
    \definition[位,名,个]{s.}{mulher; feminino; pessoa do sexo feminino}
  \end{Phonetics}
\end{Entry}

\begin{Entry}{女王}{3,4}{⼥,⽟}
  \begin{Phonetics}{女王}{nv3wang2}
    \definition{s.}{rainha; soberana}
  \end{Phonetics}
\end{Entry}

\begin{Entry}{女生}{3,5}{⼥,⽣}
  \begin{Phonetics}{女生}{nv3sheng1}[][HSK 1]
    \definition[个,位,名]{s.}{estudante; aluna; estudante do sexo feminino | menina; jovem mulher}
  \synonymref{姑娘}{gu1niang5}
  \synonymref{女孩}{nv3hai2}
  \synonymref{女性}{nv3xing4}
  \synonymref{女子}{nv3zi3}
  \synonymref{女人}{nv3ren5}
  \antonymref{男人}{nan2ren2}
  \antonymref{男生}{nan2sheng1}
  \antonymref{男士}{nan2shi4}
  \antonymref{男性}{nan2xing4}
  \antonymref{男子}{nan2zi3}
  \end{Phonetics}
\end{Entry}

\begin{Entry}{女性}{3,8}{⼥,⼼}
  \begin{Phonetics}{女性}{nv3xing4}[][HSK 5]
    \definition[个,位,名]{s.}{mulher; feminino; feminilidade}
  \antonymref{男性}{nan2xing4}
  \end{Phonetics}
\end{Entry}

\begin{Entry}{女朋友}{3,8,4}{⼥,⽉,⼜}
  \begin{Phonetics}{女朋友}{nv3peng2you5}[][HSK 1]
    \definition[个,位]{s.}{namorada}
  \synonymref{爱人}{ai4ren5}
  \synonymref{伴侣}{ban4lv3}
  \synonymref{男朋友}{nan2peng2you5}
  \synonymref{情侣}{qing2lv3}
  \end{Phonetics}
\end{Entry}

\begin{Entry}{女孩}{3,9}{⼥,⼦}
  \begin{Phonetics}{女孩}{nv3hai2}
    \definition{s.}{menina; garota; uma menor de idade}
  \synonymref{闺女}{gui1nv5}
  \synonymref{女生}{nv3sheng1}
  \synonymref{女子}{nv3zi3}
  \synonymref{丫头}{ya1tou5}
  \end{Phonetics}
\end{Entry}

\begin{Entry}{女孩儿}{3,9,2}{⼥,⼦,⼉}
  \begin{Phonetics}{女孩儿}{nv3hai2r5}[][HSK 1]
    \definition{s.}{garota; menina; atualmente também se refere a mulher adolescente | filha}
  \end{Phonetics}
\end{Entry}

\begin{Entry}{女婿}{3,12}{⼥,⼥}
  \begin{Phonetics}{女婿}{nv3xu5}[][HSK 7-9]
    \definition[个,位]{s.}{genro; marido da filha | em algumas regiões, refere"-se ao marido}
  \end{Phonetics}
\end{Entry}

%%%%%%%%%% 子 %%%%%%%%%%
\subsection*{子}\addcontentsline{loh}{figure}{子}

\begin{Entry}{子}{3}{⼦}[Kangxi 39]
  \begin{Phonetics}{子}{zi3}
    \definition*{s.}{Sobrenome: Zi}
    \definition{adj.}{pequeno; jovem; tenro | subsidiário; subordinado; derivado}
    \definition{clas.}{usado para objetos finos que podem ser pinçados com os dedos}
    \definition{pron.}{você;  antigamente, era uma forma de tratamento respeitosa para se referir a outras pessoas, equivalente a 您}
    \definition[个,位,名]{s.}{filho, criança; antigamente, referia"-se aos filhos, mas atualmente refere"-se especificamente aos filhos homens | pessoa | antigo título de respeito para um homem culto ou virtuoso; na antiguidade, referia"-se especificamente a homens eruditos | visconde; o quarto posto na hierarquia dos cinco títulos feudais da nobreza | ovo | semente | coisas pequenas e duras; pequenos fragmentos ou grãos duros e sólidos | cobre; moeda de cobre | o primeiro dos doze ramos terrestres}
  \seealsoref{您}{nin2}
  \antonymref{母}{mu3}
  \end{Phonetics}
  \begin{Phonetics}{子}{zi5}[][HSK 1]
    \definition{suf.}{sufixo para substantivos | sufixos de palavras de medida individuais; anexado a certas palavras classificadoras}
  \end{Phonetics}
\end{Entry}

\begin{Entry}{子女}{3,3}{⼦,⼥}
  \begin{Phonetics}{子女}{zi3nv3}[][HSK 3]
    \definition[个]{s.}{crianças; descendentes; filhos e filhas}
  \end{Phonetics}
\end{Entry}

\begin{Entry}{子孙}{3,6}{⼦,⼦}
  \begin{Phonetics}{子孙}{zi3sun1}[][HSK 7-9]
    \definition{s.}{filhos e netos; descendentes | prole | posteridade}
  \synonymref{儿女}{er2nv3}
  \synonymref{后代}{hou4dai4}
  \end{Phonetics}
\end{Entry}

\begin{Entry}{子弟}{3,7}{⼦,⼸}
  \begin{Phonetics}{子弟}{zi3di4}[][HSK 7-9]
    \definition{s.}{filhos e irmãos mais novos; juniores; crianças}
  \end{Phonetics}
\end{Entry}

\begin{Entry}{子弹}{3,11}{⼦,⼸}
  \begin{Phonetics}{子弹}{zi3dan4}[][HSK 5]
    \definition[粒,颗,发]{s.}{bala; cartucho; munição}
  \end{Phonetics}
\end{Entry}

%%%%%%%%%% 寸 %%%%%%%%%%
\subsection*{寸}\addcontentsline{loh}{figure}{寸}

\begin{Entry}{寸}{3}{⼨}[Kangxi 41]
  \begin{Phonetics}{寸}{cun4}[][HSK 5]
    \definition*{s.}{Sobrenome: Cun}
    \definition{adj.}{muito pouco; muito curto; pequeno | Dialeto: coincidência}
    \definition{clas.}{cun, uma unidade tradicional de comprimento, igual a 0,1 市尺 e equivalente a 3,333 centímetros ou 1,312 polegadas | cun, uma unidade de comprimento (=13 decímetros)}
  \seealsoref{市尺}{shi4chi3}
  \end{Phonetics}
\end{Entry}

%%%%%%%%%% 小 %%%%%%%%%%
\subsection*{小}\addcontentsline{loh}{figure}{小}

\begin{Entry}{小}{3}{⼩}[Kangxi 42]
  \begin{Phonetics}{小}{xiao3}[][HSK 1,2]
    \definition*{s.}{Sobrenome: Xiao}
    \definition{adj.}{menor; pequeno; insignificante; pouco; volume, área, quantidade, intensidade, etc. não são grandes | jovem | expressões humildes, referindo"-se a si mesmo ou a pessoas ou coisas relacionadas a si mesmo | por um tempo; por um curto período; por um curto período de tempo | o mais novo; o último na ordem de antiguidade; em último lugar na classificação}
    \definition{pref.}{usado antes do sobrenome, nome, posição na família, etc.}
    \definition{s.}{os jovens; pessoas mais jovens | concubina}
  \synonymref{略}{lve4}
  \synonymref{浅}{qian3}
  \synonymref{稍}{shao1}
  \synonymref{微}{wei1}
  \synonymref{细}{xi4}
  \synonymref{纤}{xian1}
  \synonymref{幼}{you4}
  \synonymref{暂}{zan4}
  \antonymref{大}{da4}
  \antonymref{巨}{ju4}
  \antonymref{老}{lao3}
  \end{Phonetics}
\end{Entry}

\begin{Entry}{小人}{3,2}{⼩,⼈}
  \begin{Phonetics}{小人}{xiao3ren2}[][HSK 7-9]
    \definition[个]{s.}{pessoa mesquinha; homem pequeno; homem ordinário; uma pessoa de caráter desprezível | ``seu humilde servo''; pessoa de baixa posição social; eu (usado por um homem de posição social inferior ao falar com seus superiores) | na antiguidade, referia"-se a pessoas de baixa condição social; era frequentemente usado como autorreferência | vilão; pessoa desprezível; personagem vil; pessoas imorais}
  \antonymref{君子}{jun1zi3}
  \end{Phonetics}
\end{Entry}

\begin{Entry}{小于}{3,3}{⼩,⼆}
  \begin{Phonetics}{小于}{xiao3yu2}[][HSK 6]
    \definition{prep.}{menor que; menos que; indica que um número ou quantidade é menor que outro}
  \end{Phonetics}
\end{Entry}

\begin{Entry}{小小}{3,3}{⼩,⼩}
  \begin{Phonetics}{小小}{xiao3xiao3}
    \definition{adj.}{muito pouco | muito pequeno}
  \end{Phonetics}
\end{Entry}

\begin{Entry}{小丑}{3,4}{⼩,⼀}
  \begin{Phonetics}{小丑}{xiao3chou3}[][HSK 7-9]
    \definition[个,名,位,帮]{s.}{palhaço; bufão; patife; na ópera tradicional chinesa, o palhaço é uma metáfora para alguém que se comporta de maneira indigna e é hábil em bajular os outros | um miserável desprezível; referindo"-se a uma pessoa insignificante}
  \end{Phonetics}
\end{Entry}

\begin{Entry}{小区}{3,4}{⼩,⼖}
  \begin{Phonetics}{小区}{xiao3qu1}[][HSK 7-9]
    \definition[个]{s.}{bairro; conjunto habitacional; comunidade residencial; áreas residenciais relativamente independentes na cidade, com instalações de serviços bem equipadas}
  \synonymref{社区}{she4qu1}
  \end{Phonetics}
\end{Entry}

\begin{Entry}{小心}{3,4}{⼩,⼼}
  \begin{Phonetics}{小心}{xiao3xin5}[][HSK 2]
    \definition{adj.}{cuidadoso; atento; com cautela}
    \definition{v.}{ter cuidado; ser cauteloso; estar atento; tomar cuidado; prestar atenção}
  \end{Phonetics}
\end{Entry}

\begin{Entry}{小心翼翼}{3,4,17,17}{⼩,⼼,⽻,⽻}
  \begin{Phonetics}{小心翼翼}{xiao3xin1-yi4yi4}[][HSK 7-9]
    \definition{expr.}{cautelosamente; cuidadosamente; com muita cautela; com grande cautela; originalmente significando solene e reverente, agora é usado para descrever ações extremamente cautelosas, sem espaço para negligência}
  \antonymref{粗心大意}{cu1xin1-da4yi4}
  \end{Phonetics}
\end{Entry}

\begin{Entry}{小气}{3,4}{⼩,⽓}
  \begin{Phonetics}{小气}{xiao3qi5}[][HSK 7-9]
    \definition{adj.}{mesquinho; apertado; avarento; parcimonioso; valorizar excessivamente os próprios bens | mesquinho; pequeno; de mente fechada}
  \antonymref{大量}{da4liang4}
  \antonymref{大气}{da4qi4}
  \antonymref{大方}{da4fang5}
  \antonymref{慷慨}{kang1kai3}
  \end{Phonetics}
\end{Entry}

\begin{Entry}{小气鬼}{3,4,9}{⼩,⽓,⿁}
  \begin{Phonetics}{小气鬼}{xiao3qi4gui3}
    \definition{adj.}{avarento; mesquinho; mão"-de"-vaca; miserável; pão"-duro}
  \end{Phonetics}
\end{Entry}

\begin{Entry}{小白菜}{3,5,11}{⼩,⽩,⾋}
  \begin{Phonetics}{小白菜}{xiao3bai2cai4}
    \definition[棵,些]{s.}{repolho chinês (\emph{Brassica chinensis}) | uma variedade de repolho chinês; \emph{pak choi} | \emph{bok choy}}
  \seealsoref{青菜}{qing1cai4}
  \seealsoref{小白菜儿}{xiao3bai2cai4r5}
  \end{Phonetics}
\end{Entry}

\begin{Entry}{小白菜儿}{3,5,11,2}{⼩,⽩,⾋,⼉}
  \begin{Phonetics}{小白菜儿}{xiao3bai2cai4r5}
    \definition{s.}{uma variedade de repolho chinês; \emph{pak choi}}
  \end{Phonetics}
\end{Entry}

\begin{Entry}{小众}{3,6}{⼩,⼈}
  \begin{Phonetics}{小众}{xiao3 zhong4}
    \definition{s.}{uma elite; uma minoria; um pequeno círculo; um grupo relativamente pequeno de pessoas; um pequeno grupo (no contexto de 大众)}
  \seealsoref{大众}{da4zhong4}
  \end{Phonetics}
\end{Entry}

\begin{Entry}{小伙子}{3,6,3}{⼩,⼈,⼦}
  \begin{Phonetics}{小伙子}{xiao3huo3zi5}[][HSK 4]
    \definition[位]{s.}{rapaz jovem; jovem colega}
  \end{Phonetics}
\end{Entry}

\begin{Entry}{小吃}{3,6}{⼩,⼝}
  \begin{Phonetics}{小吃}{xiao3chi1}[][HSK 4]
    \definition[家]{s.}{lanche; petiscos; comida com especialidades locais, não muito para uma porção | prato frio; prato feito; cortes de frios na culinária ocidental | pratos pequenos e baratos; pratos simples em restaurantes com porções pequenas e preços baixos}
  \end{Phonetics}
\end{Entry}

\begin{Entry}{小曲}{3,6}{⼩,⽈}
  \begin{Phonetics}{小曲}{xiao3qu1}[][HSK 7-9]
    \definition[首,支,段]{s.}{levedura para fabricação de vinho de arroz ou fermentação de arroz glutinoso | balada; melodia | melodia folclórica | música popular}
  \end{Phonetics}
\end{Entry}

\begin{Entry}{小声}{3,7}{⼩,⼠}
  \begin{Phonetics}{小声}{xiao3sheng1}[][HSK 2]
    \definition{v.}{falar em voz baixa; falar baixinho; sussurar}
  \end{Phonetics}
\end{Entry}

\begin{Entry}{小时}{3,7}{⼩,⽇}
  \begin{Phonetics}{小时}{xiao3shi2}[][HSK 1]
    \definition{clas.}{hora; unidade de medida legal do tempo, 1 hora equivale a 60 minutos, é 1/24 de um dia}
    \definition[个,半,些]{s.}{hora; refere"-se a um período de uma hora}
  \seealsoref{点钟}{dian3zhong1}
  \synonymref{点钟}{dian3zhong1}
  \synonymref{时刻}{shi2ke4}
  \synonymref{钟头}{zhong1tou2}
  \end{Phonetics}
\end{Entry}

\begin{Entry}{小时候}{3,7,10}{⼩,⽇,⼈}
  \begin{Phonetics}{小时候}{xiao3shi2hou5}[][HSK 2]
    \definition{s.}{na infância; quando alguém era jovem; refere"-se à infância}
  \end{Phonetics}
\end{Entry}

\begin{Entry}{小麦}{3,7}{⼩,⿆}
  \begin{Phonetics}{小麦}{xiao3mai4}[][HSK 6]
    \definition[粒,公斤,吨,棵]{s.}{trigo}
  \end{Phonetics}
\end{Entry}

\begin{Entry}{小卒}{3,8}{⼩,⼗}
  \begin{Phonetics}{小卒}{xiao3zu2}[][HSK 7-9]
    \definition{s.}{soldado de infantaria |Figurativo: ninguém | Xadrez: peão | particular | homem comum; um homem sem importância; pessoas comuns | um mero peão}
  \seealsoref{小卒子}{xiao3zu2 zi3}
  \antonymref{大王}{da4wang2}
  \antonymref{大王}{dai4wang5}
  \end{Phonetics}
\end{Entry}

\begin{Entry}{小卒子}{3,8,3}{⼩,⼗,⼦}
  \begin{Phonetics}{小卒子}{xiao3zu2 zi3}
    \definition{s.}{peão}
  \end{Phonetics}
\end{Entry}

\begin{Entry}{小姐}{3,8}{⼩,⼥}
  \begin{Phonetics}{小姐}{xiao3jie3}[][HSK 1]
    \definition[位,名,个,些]{s.}{jovem senhora; anteriormente, era assim que se referiam às filhas de famílias ricas. | senhorita; título honorífico para mulheres jovens | (gíria) prostituta}
  \seealsoref{姑娘}{gu1niang5}
  \synonymref{姑娘}{gu1niang5}
  \synonymref{女士}{nv3shi4}
  \end{Phonetics}
\end{Entry}

\begin{Entry}{小学}{3,8}{⼩,⼦}
  \begin{Phonetics}{小学}{xiao3xue2}[][HSK 1]
    \definition[所]{s.}{escola primária (ou fundamental); escolas que oferecem ensino fundamental básico | estudos filológicos; antigamente, referia"-se ao estudo da escrita, da fonética e da exegese}
  \antonymref{大学}{da4xue2}
  \end{Phonetics}
\end{Entry}

\begin{Entry}{小学生}{3,8,5}{⼩,⼦,⽣}
  \begin{Phonetics}{小学生}{xiao3xue2sheng5}[][HSK 1]
    \definition[个,位,名,些]{s.}{aluno; estudante; estudante do sexo masculino (男); estudante do sexo feminino (女) | um aluno mais novo (do que os outros da sua turma) | (dialeto) um menino pequeno}
  \seealsoref{男}{nan2}
  \seealsoref{女}{nv3}
  \end{Phonetics}
\end{Entry}

\begin{Entry}{小朋友}{3,8,4}{⼩,⽉,⼜}
  \begin{Phonetics}{小朋友}{xiao3peng2you3}[][HSK 1]
    \definition[个,名,位]{s.}{criança; crianças; refere"-se a crianças e adolescentes | (termo usado por um adulto para se dirigir a uma criança) amiguinho; menino (ou menina); termo carinhoso para se referir a crianças e adolescentes}
  \synonymref{少儿}{shao4'er2}
  \synonymref{小孩儿}{xiao3hai2r5}
  \antonymref{成人}{cheng2ren2}
  \antonymref{大人}{da4ren2}
  \end{Phonetics}
\end{Entry}

\begin{Entry}{小狗}{3,8}{⼩,⽝}
  \begin{Phonetics}{小狗}{xiao3gou3}
    \definition{s.}{filhote de cachorro; cães de pequeno porte, geralmente filhotes ou raças miniatura, são mantidos principalmente como animais de estimação}
  \end{Phonetics}
\end{Entry}

\begin{Entry}{小组}{3,8}{⼩,⽷}
  \begin{Phonetics}{小组}{xiao3zu3}[][HSK 2]
    \definition[个,名,位]{s.}{grupo; um pequeno grupo de pessoas}
  \end{Phonetics}
\end{Entry}

\begin{Entry}{小贩}{3,8}{⼩,⾙}
  \begin{Phonetics}{小贩}{xiao3fan4}[][HSK 7-9]
    \definition[个,名,位,帮]{s.}{vendedor ambulante; vendedor; mascate; comerciantes com muito pouco capital}
  \end{Phonetics}
\end{Entry}

\begin{Entry}{小品}{3,9}{⼩,⼝}
  \begin{Phonetics}{小品}{xiao3pin3}[][HSK 7-9]
    \definition[个,部,篇]{s.}{criação literária ou artística curta e simples; ensaio; esboço; esquete | esboço curto; ato curto; uma criação literária ou artística breve e simples; originalmente referindo"-se a versões abreviadas das escrituras budistas, agora se refere a ensaios curtos ou outras formas concisas de expressão}
  \end{Phonetics}
\end{Entry}

\begin{Entry}{小型}{3,9}{⼩,⼟}
  \begin{Phonetics}{小型}{xiao3xing2}[][HSK 4]
    \definition{adj.}{de tamanho pequeno; em pequena escala; miniatura; tipo pequeno; tamanho de bolso; tipo compacto}
    \definition{s.}{Mediterrâneo: escunas, pequenos veleiros de pesca ou turismo | pequeno \emph{rover} lunar (duas pessoas)}
  \end{Phonetics}
\end{Entry}

\begin{Entry}{小孩儿}{3,9,2}{⼩,⼦,⼉}
  \begin{Phonetics}{小孩儿}{xiao3hai2r5}[][HSK 1]
    \definition[个,名,位]{s.}{criança; bebê}
  \synonymref{儿童}{er2tong2}
  \synonymref{儿子}{er2zi5}
  \synonymref{孩子}{hai2zi5}
  \synonymref{女儿}{nv3'er2}
  \synonymref{小朋友}{xiao3peng2you3}
  \synonymref{子女}{zi3nv3}
  \antonymref{成人}{cheng2ren2}
  \antonymref{大人}{da4ren2}
  \end{Phonetics}
\end{Entry}

\begin{Entry}{小屋}{3,9}{⼩,⼫}
  \begin{Phonetics}{小屋}{xiao3wu1}
    \definition[间]{s.}{chalé; cabana; pousada | casita; cabana; casa geminada | cabine}
  \end{Phonetics}
\end{Entry}

\begin{Entry}{小树}{3,9}{⼩,⽊}
  \begin{Phonetics}{小树}{xiao3shu4}
    \definition[棵]{s.}{muda; árvore jovem; arbusto}
  \end{Phonetics}
\end{Entry}

\begin{Entry}{小洋白菜}{3,9,5,11}{⼩,⽔,⽩,⾋}
  \begin{Phonetics}{小洋白菜}{xiao3 yang2bai2cai4}
    \definition{s.}{couve de bruxelas}
  \end{Phonetics}
\end{Entry}

\begin{Entry}{小看}{3,9}{⼩,⽬}
  \begin{Phonetics}{小看}{xiao3kan4}[][HSK 7-9]
    \definition{v.}{subestimar; desprezar; menosprezar; desdenhar}
  \synonymref{鄙视}{bi3shi4}
  \synonymref{忽视}{hu1shi4}
  \synonymref{看不起}{kan4bu5qi3}
  \synonymref{歧视}{qi2shi4}
  \synonymref{无视}{wu2shi4}
  \antonymref{重视}{zhong4shi4}
  \end{Phonetics}
\end{Entry}

\begin{Entry}{小说}{3,9}{⼩,⾔}
  \begin{Phonetics}{小说}{xiao3shuo1}[][HSK 2]
    \definition[本,部,篇,章]{s.}{história; romance; ficção; uma forma literária que reflete a vida social por meio da descrição de personagens, ambiente e enredo}
  \end{Phonetics}
\end{Entry}

\begin{Entry}{小费}{3,9}{⼩,⾙}
  \begin{Phonetics}{小费}{xiao3fei4}[][HSK 6]
    \definition[笔]{s.}{gorjeta; gratificação; dinheiro extra pago por clientes e viajantes a funcionários de serviços em setores de serviços, como hotéis e pousadas}
  \end{Phonetics}
\end{Entry}

\begin{Entry}{小偷儿}{3,11,2}{⼩,⼈,⼉}
  \begin{Phonetics}{小偷儿}{xiao3tou1er5}[][HSK 5]
    \definition{s.}{ladrão insignificante (ou furtivo); ladrãozinho | ladrão}
  \end{Phonetics}
\end{Entry}

\begin{Entry}{小康}{3,11}{⼩,⼴}
  \begin{Phonetics}{小康}{xiao3kang1}[][HSK 7-9]
    \definition{adj.}{bem"-sucedido; (em relação ao padrão de vida) bem de vida; refere"-se à situação econômica de uma família que consegue manter um padrão de vida de classe média}
  \synonymref{富裕}{fu4yu4}
  \end{Phonetics}
\end{Entry}

\begin{Entry}{小提琴}{3,12,12}{⼩,⼿,⽟}
  \begin{Phonetics}{小提琴}{xiao3ti2qin2}[][HSK 7-9]
    \definition[个,些,把]{s.}{violino; um tipo de violino, o menor em tamanho e com o som mais agudo}
  \end{Phonetics}
\end{Entry}

\begin{Entry}{小溪}{3,13}{⼩,⽔}
  \begin{Phonetics}{小溪}{xiao3xi1}[][HSK 7-9]
    \definition[条]{s.}{riacho; córrego; um pequeno riacho raso, geralmente que flui por um vale ou floresta}
  \end{Phonetics}
\end{Entry}

\begin{Entry}{小腿}{3,13}{⼩,⾁}
  \begin{Phonetics}{小腿}{xiao3tui3}
    \definition{s.}{perna; parte inferior da perna (do joelho ao tornozelo)}
  \end{Phonetics}
\end{Entry}

\begin{Entry}{小路}{3,13}{⼩,⾜}
  \begin{Phonetics}{小路}{xiao3lu4}[][HSK 7-9]
    \definition[条]{s.}{caminho; trilha | atalho | faixa | estrada secundária | via}
  \antonymref{大道}{da4dao4}
  \end{Phonetics}
\end{Entry}

%%%%%%%%%% 尸 %%%%%%%%%%
\subsection*{尸}\addcontentsline{loh}{figure}{尸}

\begin{Entry}{尸}{3}{⼫}[Kangxi 44]
  \begin{Phonetics}{尸}{shi1}
    \definition*{s.}{Sobrenome: Shi}
    \definition[具]{s.}{cadáver; corpo morto; restos mortais | Arcaico: pessoa que se sentava atrás do altar, representando o falecido durante a realização de ritos sacrificiais}
    \definition{v.}{manter um emprego sem fazer nada (como um cadáver) | dirigir; agir como responsável | dispor}
  \end{Phonetics}
\end{Entry}

\begin{Entry}{尸体}{3,7}{⼫,⼈}
  \begin{Phonetics}{尸体}{shi1ti3}[][HSK 7-9]
    \definition[具,个]{s.}{cadáver; restos mortais; corpo morto; os corpos de humanos e animais após a morte}
  \end{Phonetics}
\end{Entry}

%%%%%%%%%% 山 %%%%%%%%%%
\subsection*{山}\addcontentsline{loh}{figure}{山}

\begin{Entry}{山}{3}{⼭}[Kangxi 46]
  \begin{Phonetics}{山}{shan1}[][HSK 1]
    \definition*{s.}{Sobrenome: Shan}
    \definition[座]{s.}{colina; maciço; montanha | qualquer coisa que se assemelhe a uma montanha | arbustos nos quais os bichos"-da"-seda tecem seus casulos; referindo"-se a casulos de bicho"-da"-seda | eco; metáfora para um som muito alto}
  \end{Phonetics}
\end{Entry}

\begin{Entry}{山川}{3,3}{⼭,⼮}
  \begin{Phonetics}{山川}{shan1chuan1}[][HSK 7-9]
    \definition{s.}{montanhas e rios; paisagem}
  \end{Phonetics}
\end{Entry}

\begin{Entry}{山冈}{3,4}{⼭,⼌}
  \begin{Phonetics}{山冈}{shan1gang1}[][HSK 7-9]
    \definition[座]{s.}{colina baixa; pequeno morro}
  \end{Phonetics}
\end{Entry}

\begin{Entry}{山区}{3,4}{⼭,⼖}
  \begin{Phonetics}{山区}{shan1qu1}[][HSK 5]
    \definition[片]{s.}{área montanhosa; região montanhosa | colina; serra; montanha | distrito montanhoso}
  \end{Phonetics}
\end{Entry}

\begin{Entry}{山东}{3,5}{⼭,⼀}
  \begin{Phonetics}{山东}{shan1dong1}
    \definition*{s.}{Província de Shandong (Shantung) no nordeste da China}
  \end{Phonetics}
\end{Entry}

\begin{Entry}{山羊}{3,6}{⼭,⽺}
  \begin{Phonetics}{山羊}{shan1yang2}
    \definition{s.}{cabra | (ginástica) cavalo de volteio de pequeno porte}
  \end{Phonetics}
\end{Entry}

\begin{Entry}{山西}{3,6}{⼭,⾑}
  \begin{Phonetics}{山西}{shan1xi1}
    \definition*{s.}{Província de Shanxi (Shansi) no norte da China entre Hebei e Shaanxi, abreviada como 晋 (Shansi)}
  \seealsoref{晋}{jin4}
  \end{Phonetics}
\end{Entry}

\begin{Entry}{山阴}{3,6}{⼭,⾩}
  \begin{Phonetics}{山阴}{shan1yin1}
    \definition*{s.}{Condado de Shanyin em Shuozhou, Shanxi}
    \definition{s.}{lado norte (ou sombreado) de uma montanha}
  \end{Phonetics}
\end{Entry}

\begin{Entry}{山体}{3,7}{⼭,⼈}
  \begin{Phonetics}{山体}{shan1ti3}
    \definition{s.}{corpo de uma montanha; maciço | forma de uma montanha}
  \end{Phonetics}
\end{Entry}

\begin{Entry}{山谷}{3,7}{⼭,⾕}
  \begin{Phonetics}{山谷}{shan1gu3}[][HSK 6]
    \definition[条,个]{s.}{vale; desfiladeiro; ravina; a área baixa e estreita entre duas montanhas geralmente tem riachos no meio}
  \end{Phonetics}
\end{Entry}

\begin{Entry}{山坡}{3,8}{⼭,⼟}
  \begin{Phonetics}{山坡}{shan1po1}[][HSK 6]
    \definition[个,座,片]{s.}{encosta; encosta da montanha; a inclinação entre o topo da montanha e o terreno plano}
  \end{Phonetics}
\end{Entry}

\begin{Entry}{山岭}{3,8}{⼭,⼭}
  \begin{Phonetics}{山岭}{shan1ling3}[][HSK 7-9]
    \definition[座,片,条]{s.}{cadeia de montanhas; cordilheira | crista; montanhas altas contínuas}
  \end{Phonetics}
\end{Entry}

\begin{Entry}{山顶}{3,8}{⼭,⾴}
  \begin{Phonetics}{山顶}{shan1ding3}[][HSK 7-9]
    \definition{s.}{topo de uma colina; cume de uma montanha; o pico de uma montanha; o topo da montanha}[他们距山顶还有100米远。===Eles ainda estavam a 100 metros do cume.]
  \end{Phonetics}
\end{Entry}

\begin{Entry}{山峰}{3,10}{⼭,⼭}
  \begin{Phonetics}{山峰}{shan1feng1}[][HSK 6]
    \definition[座,个]{s.}{pico (montanha); topo alto e pontudo da montanha}
  \end{Phonetics}
\end{Entry}

\begin{Entry}{山路}{3,13}{⼭,⾜}
  \begin{Phonetics}{山路}{shan1lu4}[][HSK 7-9]
    \definition{s.}{trilha de montanha | estrada de montanha}
  \end{Phonetics}
\end{Entry}

\begin{Entry}{山寨}{3,14}{⼭,⼧}
  \begin{Phonetics}{山寨}{shan1zhai4}[][HSK 7-9]
    \definition{s.}{fortaleza na montanha (especialmente de bandidos); aldeia fortificada na montanha; vila fortificada na colina | Humorístico: “bandido”; imitador | Figurativo: imitação (de produtos); falsificação}
  \end{Phonetics}
\end{Entry}

%%%%%%%%%% 川 %%%%%%%%%%
\subsection*{川}\addcontentsline{loh}{figure}{川}

\begin{Entry}{川}{3}{⼮}[Kangxi 47]
  \begin{Phonetics}{川}{chuan1}
    \definition*{s.}{Província de Sichuan, abreviação de 四川}
    \definition{s.}{rio; córrego | planície; terra plana}
  \seealsoref{四川}{si4chuan1}
  \end{Phonetics}
\end{Entry}

\begin{Entry}{川流不息}{3,10,4,10}{⼮,⽔,⼀,⼼}
  \begin{Phonetics}{川流不息}{chuan1liu2-bu4xi1}[][HSK 7-9]
    \definition{expr.}{``O fluxo nunca para de fluir.''; um fluxo contínuo (de); vem e vai em um fluxo sem fim; vem e vai o tempo todo; flui continuamente; continua continuamente; flui em um fluxo sem fim; flui sem cessar; flui perpetuamente; em um fluxo sem fim; sem fim; despeja em um fluxo sem fim; ``O fluxo flui incessantemente.''; descreve os pedestres, veículos, etc. indo e vindo tão continuamente quanto um fluxo de água}
  \end{Phonetics}
\end{Entry}

%%%%%%%%%% 工 %%%%%%%%%%
\subsection*{工}\addcontentsline{loh}{figure}{工}

\begin{Entry}{工}{3}{⼯}[Kangxi 48]
  \begin{Phonetics}{工}{gong1}
    \definition*{s.}{Sobrenome: Gong}
    \definition{adj.}{fino; requintado; delicado}
    \definition{s.}{trabalhador; operário; artesão | trabalho; labor; trabalho produtivo | projeto; construção; refere"-se à engenharia | indústria; refere"-se à indústria | homem"-dia; a quantidade de trabalho que um trabalhador faz em um dia | uma nota da escala em Gongchepu (工尺谱), correspondente a 3 na notação musical numerada | engenheiro; refere"-se a engenheiros}
    \definition{v.}{ser versado em; ser bom em | trabalhar em; agora geralmente escrito como 功}
  \seealsoref{功}{gong1}
  \seealsoref{工尺谱}{gong1che3 pu3}
  \end{Phonetics}
\end{Entry}

\begin{Entry}{工人}{3,2}{⼯,⼈}
  \begin{Phonetics}{工人}{gong1ren5}[][HSK 1]
    \definition[个,名]{s.}{trabalhador; operário; mão de obra; trabalhadores braçais que vivem do salário}
  \seealsoref{职员}{zhi2yuan2}
  \synonymref{员工}{yuan2gong1}
  \synonymref{职工}{zhi2gong1}
  \synonymref{职员}{zhi2yuan2}
  \end{Phonetics}
\end{Entry}

\begin{Entry}{工厂}{3,2}{⼯,⼚}
  \begin{Phonetics}{工厂}{gong1chang3}[][HSK 3]
    \definition[个,家,座,间]{s.}{fábrica; moinho; planta; unidades que realizam atividades de produção industrial diretamente, geralmente incluindo diferentes oficinas}
  \end{Phonetics}
\end{Entry}

\begin{Entry}{工夫}{3,4}{⼯,⼤}
  \begin{Phonetics}{工夫}{gong1fu1}
    \definition[个]{s.}{tempo | tempo livre; lazer}
  \end{Phonetics}
  \begin{Phonetics}{工夫}{gong1fu5}[][HSK 3]
    \definition[个]{s.}{(um período de) tempo; o tempo ou energia gastos para realizar uma tarefa | tempo livre}
  \end{Phonetics}
\end{Entry}

\begin{Entry}{工尺谱}{3,4,14}{⼯,⼫,⾔}
  \begin{Phonetics}{工尺谱}{gong1che3 pu3}
    \definition*{s.}{Gongchepu, notação musical tradicional chinesa}
    \definition{s.}{notação musical tradicional chinesa que usa caracteres chineses para representar notas musicais}
  \end{Phonetics}
\end{Entry}

\begin{Entry}{工艺}{3,4}{⼯,⾋}
  \begin{Phonetics}{工艺}{gong1yi4}[][HSK 5]
    \definition{s.}{técnica; tecnologia; arte industrial; técnicas ou métodos de fabricação e processamento de produtos | artesanato; arte artesanal}
  \end{Phonetics}
\end{Entry}

\begin{Entry}{工艺品}{3,4,9}{⼯,⾋,⼝}
  \begin{Phonetics}{工艺品}{gong1yi4pin3}
    \definition[个,件]{s.}{trabalho manual; artesanato; habilidade manual; artigo artesanal; itens delicados produzidos com técnicas artesanais. Por exemplo, esculturas em jade, esmaltes Jingtailan, bordados, etc.}
  \end{Phonetics}
\end{Entry}

\begin{Entry}{工艺流程}{3,4,10,12}{⼯,⾋,⽔,⽲}
  \begin{Phonetics}{工艺流程}{gong1yi4liu2cheng2}
    \definition{s.}{fluxograma do processo; fluxo do processo}
  \end{Phonetics}
\end{Entry}

\begin{Entry}{工业}{3,5}{⼯,⼀}
  \begin{Phonetics}{工业}{gong1ye4}[][HSK 3]
    \definition{s.}{indústria; utilização de recursos naturais; fabricação de meios de produção; meios de subsistência; ou processamento de produtos agrícolas, produtos semiacabados, etc.}
  \end{Phonetics}
\end{Entry}

\begin{Entry}{工会}{3,6}{⼯,⼈}
  \begin{Phonetics}{工会}{gong1hui4}[][HSK 7-9]
    \definition[个]{s.}{sindicato; sindicato trabalhista; organizações de massa criadas pelos trabalhadores para proteger os seus próprios interesses}
  \end{Phonetics}
\end{Entry}

\begin{Entry}{工地}{3,6}{⼯,⼟}
  \begin{Phonetics}{工地}{gong1di4}[][HSK 7-9]
    \definition{s.}{canteiro de obras; locais onde são realizadas construções, desenvolvimentos, produções, etc.}
  \end{Phonetics}
\end{Entry}

\begin{Entry}{工作}{3,7}{⼯,⼈}
  \begin{Phonetics}{工作}{gong1zuo4}[][HSK 1]
    \definition[份,个,分,项]{s.}{trabalho; emprego | dever; tarefa; negócio}
    \definition{v.}{trabalhar; operar (uma máquina); envolver"-se em trabalho físico ou intelectual, também se refere de maneira geral a máquinas e ferramentas operadas por pessoas para realizar funções produtivas}
  \seealsoref{活儿}{huo2r5}
  \seealsoref{事业}{shi4ye4}
  \synonymref{办事}{ban4shi4}
  \synonymref{就业}{jiu4/ye4}
  \synonymref{劳动}{lao2dong5}
  \synonymref{任务}{ren4wu5}
  \synonymref{上班}{shang4/ban1}
  \synonymref{事务}{shi4wu4}
  \synonymref{职业}{zhi2ye4}
  \antonymref{失业}{shi1/ye4}
  \antonymref{歇息}{xie1xi5}
  \antonymref{休憩}{xiu1qi4}
  \antonymref{休息}{xiu1xi5}
  \end{Phonetics}
\end{Entry}

\begin{Entry}{工作日}{3,7,4}{⼯,⼈,⽇}
  \begin{Phonetics}{工作日}{gong1zuo4ri4}[][HSK 5]
    \definition{s.}{dia de trabalho; dia útil; dias em que você deveria estar trabalhando de acordo com as regras | horas de trabalho por dia; horas do dia para fazer o trabalho necessário}
  \end{Phonetics}
\end{Entry}

\begin{Entry}{工作量}{3,7,12}{⼯,⼈,⾥}
  \begin{Phonetics}{工作量}{gong1zuo4liang4}[][HSK 7-9]
    \definition{s.}{quantidade de trabalho; volume de trabalho; carga de trabalho}
  \end{Phonetics}
\end{Entry}

\begin{Entry}{工序}{3,7}{⼯,⼴}
  \begin{Phonetics}{工序}{gong1xu4}[][HSK 7-9]
    \definition[道]{s.}{processo; procedimento de trabalho; sequência do processo de produção}
  \end{Phonetics}
\end{Entry}

\begin{Entry}{工具}{3,8}{⼯,⼋}
  \begin{Phonetics}{工具}{gong1ju4}[][HSK 3]
    \definition[个,件,套]{s.}{ferramenta; ferramentas utilizadas na produção | ferramenta; meio; instrumento; (metáfora) algo ou meio utilizado para atingir um determinado objetivo}
  \end{Phonetics}
\end{Entry}

\begin{Entry}{工科}{3,9}{⼯,⽲}
  \begin{Phonetics}{工科}{gong1ke1}[][HSK 7-9]
    \definition{s.}{curso de engenharia | engenharia como disciplina acadêmica; um termo geral para disciplinas de engenharia no ensino}
  \end{Phonetics}
\end{Entry}

\begin{Entry}{工资}{3,10}{⼯,⾙}
  \begin{Phonetics}{工资}{gong1zi1}[][HSK 3]
    \definition[份,笔,月,天]{s.}{pagamento; salário; remuneração; vencimentos; o pagamento em dinheiro ou em espécie feito ao trabalhador como remuneração pelo trabalho realizado}
  \end{Phonetics}
\end{Entry}

\begin{Entry}{工商}{3,11}{⼯,⼝}
  \begin{Phonetics}{工商}{gong1shang1}[][HSK 6]
    \definition{s.}{indústria e comércio; um termo combinado para indústria e comércio}
  \end{Phonetics}
\end{Entry}

\begin{Entry}{工商界}{3,11,9}{⼯,⼝,⽥}
  \begin{Phonetics}{工商界}{gong1shang1jie4}[][HSK 7-9]
    \definition{s.}{círculos industriais e comerciais; círculos de negócios | indústria | o mundo dos negócios}
  \end{Phonetics}
\end{Entry}

\begin{Entry}{工程}{3,12}{⼯,⽲}
  \begin{Phonetics}{工程}{gong1cheng2}[][HSK 4]
    \definition[个,项]{s.}{projeto; programa; trabalhos que utilizam equipamentos grandes e complexos, como projetos de reconstrução urbana e projetos de cestas de alimentos, etc. | engenharia; departamentos de produção e manufatura usam equipamentos grandes e complexos para realizar seu trabalho}
  \end{Phonetics}
\end{Entry}

\begin{Entry}{工程师}{3,12,6}{⼯,⽲,⼱}
  \begin{Phonetics}{工程师}{gong1cheng2shi1}[][HSK 3]
    \definition[个,位,名,些]{s.}{engenheiro; um dos cargos técnicos é o de especialista capaz de realizar de forma independente o projeto e a execução de uma tarefa técnica específica}
  \end{Phonetics}
\end{Entry}

\begin{Entry}{工龄}{3,13}{⼯,⿒}
  \begin{Phonetics}{工龄}{gong1ling2}
    \definition{s.}{tempo de serviço; posição; senioridade | duração do emprego; tempo de serviço como trabalhador; idade ativa; anos de trabalho | anos de serviço}
  \end{Phonetics}
\end{Entry}

\begin{Entry}{工整}{3,16}{⼯,⽁}
  \begin{Phonetics}{工整}{gong1zheng3}[][HSK 7-9]
    \definition{adj.}{limpo; organizado; meticuloso e organizado; não desleixado | fino; requintado}
  \end{Phonetics}
\end{Entry}

%%%%%%%%%% 已 %%%%%%%%%%
\subsection*{已}\addcontentsline{loh}{figure}{已}

\begin{Entry}{已}{3}{⼰}
  \begin{Phonetics}{已}{yi3}[][HSK 3]
    \definition{adv.}{já | posteriormente; mais tarde; depois de algum tempo | demasiadamente; excessivamente}
    \definition{v.}{terminar; parar; cessar}
  \end{Phonetics}
\end{Entry}

\begin{Entry}{已久}{3,3}{⼰,⼃}
  \begin{Phonetics}{已久}{yi3 jiu3}
    \definition{expr.}{``Já faz muito tempo.''}
  \end{Phonetics}
\end{Entry}

\begin{Entry}{已灭}{3,5}{⼰,⽕}
  \begin{Phonetics}{已灭}{yi3mie4}
    \definition{expr.}{extinto; destruído}
  \end{Phonetics}
\end{Entry}

\begin{Entry}{已知}{3,8}{⼰,⽮}
  \begin{Phonetics}{已知}{yi3zhi1}
    \definition{adj.}{conhecido (pela ciência)}
  \end{Phonetics}
\end{Entry}

\begin{Entry}{已经}{3,8}{⼰,⽷}
  \begin{Phonetics}{已经}{yi3jing1}[][HSK 2]
    \definition{adv.}{já; indica que uma ação ou mudança foi concluída ou atingiu um determinado nível}
  \end{Phonetics}
\end{Entry}

\begin{Entry}{已故}{3,9}{⼰,⽁}
  \begin{Phonetics}{已故}{yi3gu4}
    \definition{adj.}{falecido}
  \end{Phonetics}
\end{Entry}

\begin{Entry}{已婚}{3,11}{⼰,⼥}
  \begin{Phonetics}{已婚}{yi3hun1}
    \definition{adj.}{casado}
  \synonymref{成家}{cheng2/jia1}
  \end{Phonetics}
\end{Entry}

\begin{Entry}{已然}{3,12}{⼰,⽕}
  \begin{Phonetics}{已然}{yi3ran2}
    \definition{adv.}{já | já ser assim; já ter se tornado um fato}
    \definition{v.}{já se tornar um fato}
  \synonymref{既然}{ji4ran2}
  \synonymref{已经}{yi3jing1}
  \synonymref{早已}{zao3yi3}
  \antonymref{不曾}{bu4ceng2}
  \antonymref{即将}{ji2jiang1}
  \antonymref{尚未}{shang4wei4}
  \end{Phonetics}
\end{Entry}

%%%%%%%%%% 干 %%%%%%%%%%
\subsection*{干}\addcontentsline{loh}{figure}{干}

\begin{Entry}{干}{3}{⼲}[Kangxi 51]
  \begin{Phonetics}{干}{gan1}
    \definition*{s.}{Sobrenome: Gan}
    \definition{adj.}{seco | vazio; oco; seco | sem substância; vazio | de parentesco nominal; (parentes) não ligados por laços sanguíneos | sem água; (água) esgotada; completamente vazia | assumido como parente nominal; relação familiar reconhecida por adoção | rude; grosseiro; mal"-educado; descreve alguém que fala de forma muito direta e rude (sem delicadeza).}
    \definition{adv.}{em vão; fútil; sem propósito; para nada; sem resultado | apenas; sem nada mais | inutilmente; sem uso, sem aproveitamento | superficialmente; significa que não há conteúdo, apenas forma}
    \definition{s.}{Arcaico: escudo | margem; ribeira; margem das águas | alimentos desidratados | abreviação para os dez troncos celestiais}
    \definition{v.}{ofender; afrontar | ter a ver com; estar relacionado com; estar implicado em; interferir com | (antiquado) buscar (cargo público, remuneração, etc.) | Dialeto: deixar alguém de fora; tratar alguém com indiferença; desprezar | assediar; perturbar; criar confusão; causar estragos; bagunçar | solicitar; procurar; buscar (cargo, salário, etc.) | beber até o fim | tratar com indiferença; ignorar}
  \seealsoref{干儿}{gan1'er2}
  \seealsoref{干儿}{gan1r5}
  \antonymref{潮}{chao2}
  \antonymref{湿}{shi1}
  \antonymref{鲜}{xian1}
  \end{Phonetics}
  \begin{Phonetics}{干}{gan4}[][HSK 1]
    \definition{adj.}{capaz; competente; habilidoso}
    \definition{s.}{tronco; parte principal; corpo principal ou parte importante de algo | habilidade; capacidade; competência}
    \definition{v.}{fazer; trabalhar; cuidar; fazer coisas | ocupar o cargo de; estar envolvido em; assumir, exercer | lutar; golpear; esforçar"-se}
  \synonymref{从}{cong2}
  \synonymref{搞}{gao3}
  \synonymref{任}{ren4}
  \synonymref{糟}{zao1}
  \synonymref{做}{zuo4}
  \antonymref{潮}{chao2}
  \antonymref{好}{hao3}
  \antonymref{湿}{shi1}
  \antonymref{鲜}{xian1}
  \antonymref{支}{zhi1}
  \end{Phonetics}
\end{Entry}

\begin{Entry}{干儿}{3,2}{⼲,⼉}
  \begin{Phonetics}{干儿}{gan1'er2}
    \definition{s.}{filho adotivo (adoção tradicional, ou seja, sem implicações legais)}
  \end{Phonetics}
  \begin{Phonetics}{干儿}{gan1r5}
    \definition{s.}{alimentos secos, desidratados}
  \end{Phonetics}
\end{Entry}

\begin{Entry}{干与}{3,3}{⼲,⼀}
  \begin{Phonetics}{干与}{gan1yu4}
    \variantof{干预}
  \end{Phonetics}
\end{Entry}

\begin{Entry}{干什么}{3,4,3}{⼲,⼈,⼃}
  \begin{Phonetics}{干什么}{gan4shen2me5}[][HSK 1]
    \definition{adv.}{o que fazer; o que ele está fazendo?; o que você está fazendo?; perguntar a razão ou o objetivo}
  \synonymref{为何}{wei4he2}
  \synonymref{为什么}{wei4shen2me5}
  \end{Phonetics}
\end{Entry}

\begin{Entry}{干戈}{3,4}{⼲,⼽}
  \begin{Phonetics}{干戈}{gan1ge1}[][HSK 7-9]
    \definition{s.}{armas de guerra; armas; guerra}[化干戈为玉帛。===Transforme guerra em paz.]
  \antonymref{玉帛}{yu4bo2}
  \end{Phonetics}
\end{Entry}

\begin{Entry}{干吗}{3,6}{⼲,⼝}
  \begin{Phonetics}{干吗}{gan4ma2}[][HSK 3]
    \definition{pron.}{por que?}
    \definition{v.}{o que fazer?}
  \end{Phonetics}
\end{Entry}

\begin{Entry}{干你屁事}{3,7,7,8}{⼲,⼈,⼫,⼅}
  \begin{Phonetics}{干你屁事}{gan1 ni3 pi4shi4}
    \definition{interj.}{``Foda"-se!''; ``O que isso te importa?''}
  \end{Phonetics}
\end{Entry}

\begin{Entry}{干扰}{3,7}{⼲,⼿}
  \begin{Phonetics}{干扰}{gan1rao3}[][HSK 5]
    \definition{v.}{perturbar; incomodar | interferir; interromper o funcionamento adequado de equipamentos eletrônicos com sinais eletrônicos dispersos}
  \end{Phonetics}
\end{Entry}

\begin{Entry}{干旱}{3,7}{⼲,⽇}
  \begin{Phonetics}{干旱}{gan1han4}[][HSK 7-9]
    \definition{adj.}{árido; seco}
  \end{Phonetics}
\end{Entry}

\begin{Entry}{干事}{3,8}{⼲,⼅}
  \begin{Phonetics}{干事}{gan4shi5}[][HSK 7-9]
    \definition{s.}{secretário (ou funcionário administrativo) encarregado de algo | administrador | secretária executiva}
  \end{Phonetics}
\end{Entry}

\begin{Entry}{干净}{3,8}{⼲,⼎}
  \begin{Phonetics}{干净}{gan1jing4}[][HSK 1]
    \definition{adj.}{limpo; limpo e arrumado; sem poeira, impurezas, etc.}
    \definition{adv.}{completamente; totalmente; sem deixar nada para trás}
  \seealsoref{纯洁}{chun2jie2}
  \seealsoref{清洁}{qing1jie2}
  \synonymref{纯洁}{chun2jie2}
  \synonymref{洁净}{jie2jing4}
  \synonymref{清洁}{qing1jie2}
  \synonymref{清爽}{qing1shuang3}
  \synonymref{洗净}{xi3jing4}
  \synonymref{整洁}{zheng3jie2}
  \antonymref{灰尘}{hui1chen2}
  \antonymref{污秽}{wu1hui4}
  \antonymref{犹豫}{you2yu4}
  \end{Phonetics}
\end{Entry}

\begin{Entry}{干杯}{3,8}{⼲,⽊}
  \begin{Phonetics}{干杯}{gan1/bei1}[][HSK 2]
    \definition{interj.}{`Saúde!''}
    \definition{v.+compl.}{fazer um brinde;  brindar até a última gota}
  \end{Phonetics}
\end{Entry}

\begin{Entry}{干活}{3,9}{⼲,⽔}
  \begin{Phonetics}{干活}{gan4/huo2}
    \definition{v.+compl.}{trabalhar; trabalhar em um emprego; fazer algo que exige esforço físico ou mental, especialmente trabalho árduo ou laborioso}
  \synonymref{工作}{gong1zuo4}
  \synonymref{劳动}{lao2dong5}
  \antonymref{歇息}{xie1xi5}
  \antonymref{休憩}{xiu1qi4}
  \antonymref{休息}{xiu1xi5}
  \end{Phonetics}
\end{Entry}

\begin{Entry}{干活儿}{3,9,2}{⼲,⽔,⼉}
  \begin{Phonetics}{干活儿}{gan4huo2r5}[][HSK 2]
    \definition{v.}{trabalhar; gastar energia física ou mental para fazer algo, especialmente trabalho árduo ou esforçado.}
  \end{Phonetics}
\end{Entry}

\begin{Entry}{干涉}{3,10}{⼲,⽔}
  \begin{Phonetics}{干涉}{gan1she4}[][HSK 6]
    \definition{s.}{interferência; refere"-se ao ato ou comportamento de interferir nos assuntos dos outros}
    \definition{v.}{interferir; intervir; intrometer"-se; pedir ou impedir algo geralmente significa interferir quando não se deve}
  \end{Phonetics}
\end{Entry}

\begin{Entry}{干脆}{3,10}{⼲,⾁}
  \begin{Phonetics}{干脆}{gan1cui4}[][HSK 5]
    \definition{adj.}{claro; direto; (falar, fazer coisas) sem hesitação; atitude clara}
    \definition{adv.}{justamente; diretamente; sem maiores considerações}
  \end{Phonetics}
\end{Entry}

\begin{Entry}{干部}{3,10}{⼲,⾢}
  \begin{Phonetics}{干部}{gan4bu4}[][HSK 7-9]
    \definition[名,个,位]{s.}{quadro; funcionário público; funcionário do governo; servidores públicos (excluindo soldados e pessoal geral) em órgãos estatais, militares e organizações populares; refere"-se àqueles que desempenham determinado trabalho de liderança ou gestão | líder; pessoal que ocupa determinados cargos de liderança}
  \end{Phonetics}
\end{Entry}

\begin{Entry}{干预}{3,10}{⼲,⾴}
  \begin{Phonetics}{干预}{gan1yu4}[][HSK 5]
    \definition{s.}{intromissão; intervenção}
    \definition{v.}{intrometer"-se; intervir; interpor"-se;}
  \end{Phonetics}
\end{Entry}

\begin{Entry}{干燥}{3,17}{⼲,⽕}
  \begin{Phonetics}{干燥}{gan1zao4}[][HSK 7-9]
    \definition{adj.}{seco; árido; sem umidade ou muito pouca umidade | enfadonho; desinteressante; chato, sem graça}
  \end{Phonetics}
\end{Entry}

%%%%%%%%%% 广 %%%%%%%%%%
\subsection*{广}\addcontentsline{loh}{figure}{广}

\begin{Entry}{广}{3}{⼴}[Kangxi 53]
  \begin{Phonetics}{广}{an1}
    \definition{s.}{mais comum em nomes de pessoas; o mesmo que 庵}[广安是我的朋友。===An'an é meu amigo.]
  \seealsoref{庵}{an1}
  \end{Phonetics}
  \begin{Phonetics}{广}{guang3}[][HSK 5]
    \definition*{s.}{Sobrenome: Guang}
    \definition{adj.}{largo; vasto; amplo; extenso | numeroso | comum; universal}
    \definition{s.}{Guangdong, 广东, e Guangxi, 广州}
    \definition{v.}{expandir; espalhar; ampliar}
  \seealsoref{广东}{guang3dong1}
  \seealsoref{广州}{guang3zhou1}
  \antonymref{狭}{xia2}
  \antonymref{窄}{zhai3}
  \end{Phonetics}
  \begin{Phonetics}{广}{yan3}
    \definition[家]{s.}{casa ou edifício construído contra ou ao longo da encosta de uma montanha ou penhasco}
  \end{Phonetics}
\end{Entry}

\begin{Entry}{广义}{3,3}{⼴,⼂}
  \begin{Phonetics}{广义}{guang3yi4}[][HSK 7-9]
    \definition*{s.}{Província de Quang Ngai; nome de lugar vietnamita, uma das províncias do Vietnã Central}
    \definition{s.}{sentido amplo; sentido geral; definição mais ampla}
  \end{Phonetics}
\end{Entry}

\begin{Entry}{广大}{3,3}{⼴,⼤}
  \begin{Phonetics}{广大}{guang3da4}[][HSK 3]
    \definition{adj.}{muito difundido; enorme (alcance, escala) | (uma área ou espaço) vasto; extenso; em grande escala; amplo (área, espaço) | numeroso; muitos (número de pessoas)}
  \end{Phonetics}
\end{Entry}

\begin{Entry}{广东}{3,5}{⼴,⼀}
  \begin{Phonetics}{广东}{guang3dong1}
    \definition*{s.}{Província de Guangdong}
  \seealsoref{粤}{yue4}
  \end{Phonetics}
\end{Entry}

\begin{Entry}{广场}{3,6}{⼴,⼟}
  \begin{Phonetics}{广场}{guang3chang3}[][HSK 2]
    \definition{s.}{praça; praça pública; esplanada; área ampla, especificamente uma área ampla na cidade}
  \end{Phonetics}
\end{Entry}

\begin{Entry}{广场舞}{3,6,14}{⼴,⼟,⾇}
  \begin{Phonetics}{广场舞}{guang3chang3 wu3}
    \definition{s.}{quadrilha, uma rotina de exercícios tocada com música em quadrados públicos, parques e praças, popular especialmente entre mulheres de meia"-idade e aposentados na China}
  \end{Phonetics}
\end{Entry}

\begin{Entry}{广州}{3,6}{⼴,⼮}
  \begin{Phonetics}{广州}{guang3zhou1}
    \definition*{s.}{Guangzhou, antigamente Cantão; Capital da Província de Guangdong}
  \end{Phonetics}
\end{Entry}

\begin{Entry}{广西}{3,6}{⼴,⾑}
  \begin{Phonetics}{广西}{guang3xi1}
    \definition*{s.}{Guangxi (Região Autônoma de Zhuang)}
  \seealsoref{壮}{zhuang4}
  \end{Phonetics}
\end{Entry}

\begin{Entry}{广告}{3,7}{⼴,⼝}
  \begin{Phonetics}{广告}{guang3gao4}[][HSK 2]
    \definition[则,条,段,项,个]{s.}{anúncio; propaganda; uma forma de divulgação ao público de produtos, serviços ou programas culturais e esportivos, geralmente realizada por meio de jornais, televisão, rádio, cartazes, etc.}
    \definition{v.}{anunciar; a ação ou ato de promover ou divulgar algo}
  \end{Phonetics}
\end{Entry}

\begin{Entry}{广泛}{3,7}{⼴,⽔}
  \begin{Phonetics}{广泛}{guang3fan4}[][HSK 5]
    \definition{adj.}{amplo; extenso; de grande alcance; disseminado; escopo e cobertura amplos}
  \end{Phonetics}
\end{Entry}

\begin{Entry}{广阔}{3,12}{⼴,⾨}
  \begin{Phonetics}{广阔}{guang3kuo4}[][HSK 6]
    \definition{adj.}{vasto; largo; amplo}
  \end{Phonetics}
\end{Entry}

\begin{Entry}{广播}{3,15}{⼴,⼿}
  \begin{Phonetics}{广播}{guang3bo1}[][HSK 3]
    \definition[个,次,段,则,条]{s.}{programa de rádio; transmissão (de rádio); refere"-se a programas transmitidos por estações de rádio ou televisão a cabo}
    \definition{v.}{transmitir; estar no ar | espalhar"-se amplamente; ser conhecido em toda parte; divulgar amplamente}
  \end{Phonetics}
\end{Entry}

%%%%%%%%%% 弓 %%%%%%%%%%
\subsection*{弓}\addcontentsline{loh}{figure}{弓}

\begin{Entry}{弓}{3}{⼸}[Kangxi 57]
  \begin{Phonetics}{弓}{gong1}[][HSK 7-9]
    \definition*{s.}{Sobrenome: Gong}
    \definition{clas.}{uma antiga unidade de comprimento para medir a terra, igual a cinco 尺}
    \definition[张]{s.}{arco | qualquer coisa em forma de arco | Obsoleto: ferramenta de madeira de medição de terreno; divisores de madeira para medição de terrenos; arco de medição; régua escalonada (1,5m)}
    \definition{v.}{dobrar; arquear; curvar; entortar}
  \seealsoref{尺}{chi3}
  \end{Phonetics}
\end{Entry}

%%%%%%%%%% 才 %%%%%%%%%%
\subsection*{才}\addcontentsline{loh}{figure}{才}

\begin{Entry}{才}{3}{⼿}
  \begin{Phonetics}{才}{cai2}[][HSK 2,4]
    \definition*{s.}{Sobrenome: Cai}
    \definition{adv.}{há pouco; agora mesmo | (precedido por uma expressão de tempo) não até | (precedido por uma expressão de razão ou condição) não a menos que; não até que; então e somente então; por nenhuma outra razão | (seguido por uma expressão numérica) apenas; indica um intervalo pequeno ou uma quantidade reduzida, equivalente a 仅仅 ou 只 | (em uma afirmação ou negação, enfatizando o que vem antes de 才, geralmente com 呢 no final da frase) na verdade; realmente | dica que algo acontece tarde ou termina tarde | (precedido por uma expressão de tempo) não até; indicando que não era assim, mas agora surgiu uma nova situação | (precedido por uma expressão de razão ou condição) a menos que; indica que só em determinadas condições e, em seguida, como | (expressa ênfase )}
    \definition{s.}{habilidade; talento; dom | pessoa competente | pessoas de um determinado tipo (frequentemente usado como sufixo) | dotação; talento; habilidade}
  \seealsoref{呢}{ne5}
  \synonymref{刚}{gang1}
  \synonymref{就}{jiu4}
  \end{Phonetics}
\end{Entry}

\begin{Entry}{才华}{3,6}{⼿,⼗}
  \begin{Phonetics}{才华}{cai2hua2}[][HSK 7-9]
    \definition[份]{s.}{talento literário; talento artístico; talentos que são exibidos externamente (principalmente nas artes)}
  \end{Phonetics}
\end{Entry}

\begin{Entry}{才能}{3,10}{⼿,⾁}
  \begin{Phonetics}{才能}{cai2neng2}[][HSK 3]
    \definition[间]{s.}{talento; habilidade; dom; capacidade; inteligência e habilidade}
  \end{Phonetics}
\end{Entry}

\begin{Entry}{才略}{3,11}{⼿,⽥}
  \begin{Phonetics}{才略}{cai2lve4}
    \definition{s.}{habilidade e sabedoria (em assuntos políticos e militares); grande talento e estratégia | habilidade e sagacidade; talento e sabedoria política ou militar}
  \synonymref{本领}{ben3ling3}
  \synonymref{才华}{cai2hua2}
  \synonymref{才能}{cai2neng2}
  \synonymref{能力}{neng2li4}
  \end{Phonetics}
\end{Entry}

%%%%%%%%%% 门 %%%%%%%%%%
\subsection*{门}\addcontentsline{loh}{figure}{门}

\begin{Entry}{门}{3}{⾨}[Kangxi 169]
  \begin{Phonetics}{门}{men2}[][HSK 1]
    \definition*{s.}{Sobrenome: Men}
    \definition{clas.}{para equipamentos de artilharia (por exemplo: canhões) | para trabalhos escolares, ciência e tecnologia, etc. | para idiomas | para casamentos | para parentes}
    \definition[个,把,道,扇]{s.}{entradas e saídas de edifícios, veículos, navios, aviões, etc. | válvula; interruptor; algo que funciona como um interruptor ou como uma porta | habilidade; método; acesso; maneira de fazer algo | família; ramo de uma família ou clã | seita (religiosa); escola (de pensamento); faculdades acadêmicas, ideológicas ou religiosas | classe; categoria; ramo de estudo; refere"-se à categoria geral de coisas | filo; segundo nível da classificação biológica | (computador) \emph{gate}; porta (lógica) | porta; portão; entrada; refere"-se a uma porta que pode ser aberta e fechada, instalada na entrada e saída | qualquer abertura; partes de objetos que podem ser abertas e fechadas | orifício no corpo humano; refere"-se especificamente aos orifícios do corpo humano | estudar com o mesmo professor; refere"-se especificamente ao professor ou mestre | posição em um jogo de apostas (em relação ao local onde se senta ou onde se faz uma aposta)}
  \end{Phonetics}
\end{Entry}

\begin{Entry}{门口}{3,3}{⾨,⼝}
  \begin{Phonetics}{门口}{men2kou3}[][HSK 1]
    \definition[个]{s.}{porta; portão; entrada; porta de entrada}
  \synonymref{大门}{da4men2}
  \end{Phonetics}
\end{Entry}

\begin{Entry}{门当户对}{3,6,4,5}{⾨,⼹,⼾,⼨}
  \begin{Phonetics}{门当户对}{men2dang1-hu4dui4}[][HSK 7-9]
    \definition{expr.}{``Casar com alguém de posição social equivalente.''; compatibilidade social e econômica adequada (para casamento); (um possível parceiro para casamento) uma combinação adequada; as famílias são bem compatíveis em termos de status social}
  \end{Phonetics}
\end{Entry}

\begin{Entry}{门诊}{3,7}{⾨,⾔}
  \begin{Phonetics}{门诊}{men2zhen3}[][HSK 5]
    \definition{s.}{(no hospital) clínica ambulatorial; seção para pacientes ambulatoriais; local onde os médicos atendem pacientes que não estão internados no hospital}
  \end{Phonetics}
\end{Entry}

\begin{Entry}{门铃}{3,10}{⾨,⾦}
  \begin{Phonetics}{门铃}{men2ling2}[][HSK 7-9]
    \definition{s.}{campainha; uma campainha ou campainha elétrica para bater à porta}
  \end{Phonetics}
\end{Entry}

\begin{Entry}{门票}{3,11}{⾨,⽰}
  \begin{Phonetics}{门票}{men2piao4}[][HSK 1]
    \definition[张]{s.}{bilhete de entrada; bilhete de admissão; ingressos para locais de turismo, entretenimento, etc.}
  \synonymref{入场券}{ru4chang3quan4}
  \synonymref{通行证}{tong1xing2zheng4}
  \end{Phonetics}
\end{Entry}

\begin{Entry}{门道}{3,12}{⾨,⾡}
  \begin{Phonetics}{门道}{men2dao4}
    \definition{s.}{porta | portal}
  \end{Phonetics}
  \begin{Phonetics}{门道}{men2dao5}
    \definition{s.}{talento | a maneira de fazer algo}
  \end{Phonetics}
\end{Entry}

\begin{Entry}{门路}{3,13}{⾨,⾜}
  \begin{Phonetics}{门路}{men2lu5}[][HSK 7-9]
    \definition{s.}{maneira de fazer algo; jeito | conexões sociais (para conseguir empregos, etc.) | jeito; saber fazer; truque do ofício; o segredo para fazer as coisas acontecerem}
  \seealsoref{门道}{men2dao5}
  \end{Phonetics}
\end{Entry}

\begin{Entry}{门槛}{3,14}{⾨,⽊}
  \begin{Phonetics}{门槛}{men2kan3}[][HSK 7-9]
    \definition{s.}{soleira; batente da porta; a viga horizontal ou faixa de pedra na parte inferior do batente da porta, próxima ao chão, etc. | o requisito para entrar em um determinado campo; metaforicamente, refere"-se aos padrões ou condições para entrar em um determinado intervalo}
  \end{Phonetics}
\end{Entry}

%%%%%%%%%% 飞 %%%%%%%%%%
\subsection*{飞}\addcontentsline{loh}{figure}{飞}

\begin{Entry}{飞}{3}{⾶}[Kangxi 183]
  \begin{Phonetics}{飞}{fei1}[][HSK 1]
    \definition{adj.}{inesperado; acidental; surgido do nada}
    \definition{adv.}{rapidamente; velozmente}
    \definition{s.}{roda livre de uma bicicleta}
    \definition{v.}{voar; esvoaçar; (pássaros, insetos, etc.) voar pelo ar batendo as asas | voar; utilizar máquinas motorizadas para se deslocar no ar | voar; (objetos naturais) flutuar ou esvoaçar no ar | volatilizar; evaporar; um gás se dissipar no ar | ir muito rapidamente; movimentar"-se rapidamente, como se estivesse voando}
  \synonymref{飘}{piao1}
  \end{Phonetics}
\end{Entry}

\begin{Entry}{飞机}{3,6}{⾶,⽊}
  \begin{Phonetics}{飞机}{fei1ji1}[][HSK 1]
    \definition[架,个,班]{s.}{avião; aeronave; aroplano}
  \synonymref{客机}{ke4ji1}
  \end{Phonetics}
\end{Entry}

\begin{Entry}{飞机票}{3,6,11}{⾶,⽊,⽰}
  \begin{Phonetics}{飞机票}{fei1ji1piao4}
    \definition[张]{s.}{bilhete de avião; documento emitido mediante pagamento de passagem aérea, que autoriza o titular a viajar}
  \seealsoref{机票}{ji1piao4}
  \end{Phonetics}
\end{Entry}

\begin{Entry}{飞行}{3,6}{⾶,⾏}
  \begin{Phonetics}{飞行}{fei1xing2}[][HSK 3]
    \definition{s.}{voo; aviação}
    \definition{v.}{voar; fazer um voo; (aviões, foguetes, etc.) voar no ar}
  \end{Phonetics}
\end{Entry}

\begin{Entry}{飞行员}{3,6,7}{⾶,⾏,⼝}
  \begin{Phonetics}{飞行员}{fei1xing2yuan2}[][HSK 6]
    \definition[名,班]{s.}{piloto; aviador; pilotos de aeronaves}
  \end{Phonetics}
\end{Entry}

\begin{Entry}{飞往}{3,8}{⾶,⼻}
  \begin{Phonetics}{飞往}{fei1wang3}[][HSK 7-9]
    \definition{v.}{voar para}[飞机将径直飞往圣保罗。===O avião voará diretamente para São Paulo.]
  \end{Phonetics}
\end{Entry}

\begin{Entry}{飞速}{3,10}{⾶,⾡}
  \begin{Phonetics}{飞速}{fei1su4}[][HSK 7-9]
    \definition{adj.}{em velocidade máxima}
  \end{Phonetics}
\end{Entry}

\begin{Entry}{飞弹}{3,11}{⾶,⼸}
  \begin{Phonetics}{飞弹}{fei1dan4}
    \definition{s.}{míssil (guiado) | bala perdida; bala disparada aleatoriamente | avião de mísseis; bomba voadora; bombas com equipamento de voo automático | dardo}
  \end{Phonetics}
\end{Entry}

\begin{Entry}{飞船}{3,11}{⾶,⾈}
  \begin{Phonetics}{飞船}{fei1chuan2}[][HSK 6]
    \definition{s.}{nave espacial; espaçonave | dirigível; aerobarco}
  \end{Phonetics}
\end{Entry}

\begin{Entry}{飞跃}{3,11}{⾶,⾜}
  \begin{Phonetics}{飞跃}{fei1yue4}[][HSK 7-9]
    \definition{v.}{saltar; crescer rapidamente}
  \end{Phonetics}
\end{Entry}

\begin{Entry}{飞翔}{3,12}{⾶,⽻}
  \begin{Phonetics}{飞翔}{fei1xiang2}[][HSK 7-9]
    \definition{v.}{pairar; voar; circular no ar; voar em círculos}
  \end{Phonetics}
\end{Entry}

\begin{Entry}{飞碟}{3,14}{⾶,⽯}
  \begin{Phonetics}{飞碟}{fei1die2}
    \definition[个]{s.}{tiro ao prato; tiro ao alvo; tiro ao pombo de barro | disco"-voador, OVNI, \emph{UFO} | \emph{frisbee}}
  \synonymref{飞机}{fei1ji1}
  \end{Phonetics}
\end{Entry}

%%%%%%%%%% 马 %%%%%%%%%%
\subsection*{马}\addcontentsline{loh}{figure}{马}

\begin{Entry}{马}{3}{⾺}[Kangxi 187]
  \begin{Phonetics}{马}{ma3}[][HSK 3]
    \definition*{s.}{Sobrenome: Ma}
    \definition{adj.}{grande; extenso; amplo}
    \definition[匹,头,只,群]{s.}{cavalo | a peça do cavalo no xadrez chinês}
  \end{Phonetics}
\end{Entry}

\begin{Entry}{马力}{3,2}{⾺,⼒}
  \begin{Phonetics}{马力}{ma3li4}[][HSK 7-9]
    \definition{clas.}{Física: cavalos de potência, cavalo"-vapor (cv)}
  \end{Phonetics}
\end{Entry}

\begin{Entry}{马上}{3,3}{⾺,⼀}
  \begin{Phonetics}{马上}{ma3shang4}[][HSK 1]
    \definition{adv.}{imediatamente; de uma só vez; em um piscar de olhos | em breve; em um futuro próximo; em um curto espaço de tempo}
  \seealsoref{立刻}{li4ke4}
  \seealsoref{眼看}{yan3kan4}
  \synonymref{当即}{dang1ji2}
  \synonymref{顿时}{dun4shi2}
  \synonymref{赶紧}{gan3jin3}
  \synonymref{急忙}{ji2mang2}
  \synonymref{立即}{li4ji2}
  \synonymref{立刻}{li4ke4}
  \antonymref{稍后}{shao1hou4}
  \antonymref{逐渐}{zhu2jian4}
  \end{Phonetics}
\end{Entry}

\begin{Entry}{马马虎虎}{3,3,8,8}{⾺,⾺,⾌,⾌}
  \begin{Phonetics}{马马虎虎}{ma3ma3hu3hu3}
    \definition{adj.}{tolerável; aceitável; mais ou menos; razoável; não tão ruim | descuidado; casual; vago; descuidado e negligente na execução das tarefas}
  \end{Phonetics}
\end{Entry}

\begin{Entry}{马车}{3,4}{⾺,⾞}
  \begin{Phonetics}{马车}{ma3che1}[][HSK 6]
    \definition[辆]{s.}{carruagem (puxada por cavalo); carroça; charrete}
  \end{Phonetics}
\end{Entry}

\begin{Entry}{马列主义}{3,6,5,3}{⾺,⼑,⼂,⼂}
  \begin{Phonetics}{马列主义}{ma3-lie4 zhu3yi4}
    \definition*{s.}{Marxismo"-Leninismo}
  \end{Phonetics}
\end{Entry}

\begin{Entry}{马后炮}{3,6,9}{⾺,⼝,⽕}
  \begin{Phonetics}{马后炮}{ma3hou4pao4}[][HSK 7-9]
    \definition{s.}{(termo do xadrez chinês) ação ou conselho tardio; esforço tardio; tarde demais | Figurativo: ação tardia | visão retrospectiva}
  \end{Phonetics}
\end{Entry}

\begin{Entry}{马戏}{3,6}{⾺,⼽}
  \begin{Phonetics}{马戏}{ma3xi4}[][HSK 7-9]
    \definition[场,次]{s.}{circo | espetáculo de circo (apresentação)}
  \end{Phonetics}
\end{Entry}

\begin{Entry}{马耳他}{3,6,5}{⾺,⽿,⼈}
  \begin{Phonetics}{马耳他}{ma3'er3ta1}
    \definition*{s.}{Malta}
  \end{Phonetics}
\end{Entry}

\begin{Entry}{马克思列宁主义}{3,7,9,6,5,5,3}{⾺,⼗,⼼,⼑,⼧,⼂,⼂}
  \begin{Phonetics}{马克思列宁主义}{ma3ke4si1lie4ning2zhu3yi4}
    \definition*{s.}{Marxismo"-Leninismo}
  \end{Phonetics}
\end{Entry}

\begin{Entry}{马尾}{3,7}{⾺,⼫}
  \begin{Phonetics}{马尾}{ma3wei3}
    \definition*{s.}{Distrito de Mawei da cidade de Fuzhou (福州), Fujian (福建)}
    \definition[个]{s.}{cauda de cavalo | rabo de cavalo (penteado) | fibras finas como a cauda de um cavalo (aplica"-se a várias plantas)}
  \seealsoref{福建}{fu2jian4}
  \seealsoref{福州}{fu2zhou1}
  \end{Phonetics}
\end{Entry}

\begin{Entry}{马虎}{3,8}{⾺,⾌}
  \begin{Phonetics}{马虎}{ma3hu5}[][HSK 7-9]
    \definition{adj.}{descuidado; casual | superficial; apressado; descuidado}
    \definition{v.}{encarar de forma leviana; fazer um trabalho malfeito}
  \end{Phonetics}
\end{Entry}

\begin{Entry}{马桶}{3,11}{⾺,⽊}
  \begin{Phonetics}{马桶}{ma3tong3}[][HSK 7-9]
    \definition[个,台]{s.}{vaso sanitário}[小心别把手机掉进马桶。===Tenha cuidado para não deixar seu celular cair no vaso sanitário.]
  \end{Phonetics}
\end{Entry}

\begin{Entry}{马路}{3,13}{⾺,⾜}
  \begin{Phonetics}{马路}{ma3lu4}[][HSK 1]
    \definition[条,段]{s.}{estrada; rua; avenida; estradas largas e planas para o tráfego de carros e cavalos nas cidades ou nos subúrbios}
  \synonymref{车道}{che1dao4}
  \synonymref{大道}{da4dao4}
  \synonymref{道路}{dao4lu4}
  \synonymref{公路}{gong1lu4}
  \synonymref{街道}{jie1dao4}
  \end{Phonetics}
\end{Entry}

%%%%% EOF %%%%%

